\documentclass[11pt]{article}
\usepackage{ortak}
\title{\color{ana}\textbf{Mühendislik Matematiği}\\[4pt]\Large Soru Çözümü -- Derste Çözülen Örnekler ve Ek Alıştırmalar}
\author{Ders tahtası fotoğrafları (Ekim--Kasım 2026) temel alınmış;\\ Kreyszig, Zill \& Wright (çev. Cangül) ve Balcı kitaplarıyla genişletilmiştir}
\date{9 Kasım 2026}
\begin{document}
\maketitle
\thispagestyle{empty}
\begin{bilgi}[Bu belge hakkında]
Çıkmış sınav soruları bulunmadığından bu belgede derste tahtada çözülen \textbf{bütün} örnekler konu konu toplanmıştır: önce soru metni, ardından adım adım çözüm. Her sorunun yanında köşeli parantez içinde hangi tahta fotoğrafından alındığı yazılıdır. \ekal etiketli sorular derste çözülmemiş, kaynak kitaplardaki alıştırmalara benzer biçimde eklenmiş pekiştirme sorularıdır. Kendini sınamak için çözümleri kapatıp önce soruyu kendin çöz.
\end{bilgi}
\tableofcontents
\clearpage\section{Konu 1: Üç Boyutlu Koordinat Sistemi, Vektörler, Skaler ve Vektörel Çarpım}

\begin{bilgi}[Bu bölüm hakkında]
Bu bölümde 09.10 defter notlarında ve 23.10 tahtasında çözülen bütün örnekler, 02.10 tahtasındaki skaler çarpım uygulaması ve derste verilen çalışma soruları, işlenme sırasıyla ve adım adım çözümleriyle toplanmıştır. Aynı soru birden çok fotoğrafta görünüyorsa bütün fotoğraf adları verilmiştir. Bölüm sonunda kaynak kitaplardaki tiplere benzer ek alıştırmalar vardır.
\end{bilgi}

%==========================================================
\subsection{Düzlemde vektörler: uzunluk, birim vektör, skaler çarpım}

\begin{soru}[Soru 1 \foto{IMG-20261009-WA0026}]
$\mathbf u=(3,4)$ vektörünün uzunluğunu ve $\mathbf u$ yönündeki birim vektörü bulunuz.
\end{soru}
\begin{cevap}[Çözüm (adım adım)]
\textbf{1. Uzunluk.} Düzlemde $\|(a,b)\|=\sqrt{a^2+b^2}$ olduğundan
\[
\|\mathbf u\|=\sqrt{3^2+4^2}=\sqrt{9+16}=\sqrt{25}=5 .
\]
\textbf{2. Birim vektör.} Sıfırdan farklı bir vektörü kendi uzunluğuna bölersek yönü değişmez, uzunluğu 1 olur:
\[
\mathbf e=\frac{\mathbf u}{\|\mathbf u\|}=\frac{(3,4)}{5}=\left(\frac35,\ \frac45\right).
\]
\textbf{3. Kontrol.} $\|\mathbf e\|=\sqrt{\frac{9}{25}+\frac{16}{25}}=1$ \checkmark. Ayrıca $\mathbf e=(\cos\theta,\sin\theta)$ yazılışından $\mathbf u$'nun $x$-ekseniyle yaptığı açı $\theta=\arccos\frac35\approx53{,}13^\circ$'dir.
\end{cevap}

\begin{soru}[Soru 2 \foto{IMG-20261009-WA0026}]
$(1,2)$ ve $(4,-2)$ vektörlerinin dik olup olmadığını belirleyiniz.
\end{soru}
\begin{cevap}[Çözüm (adım adım)]
\textbf{1. Ölçüt.} Sıfırdan farklı iki vektör ancak ve ancak skaler çarpımları sıfırsa diktir ($\cos90^\circ=0$).

\textbf{2. Hesap.} $(1,2)\cdot(4,-2)=1\cdot4+2\cdot(-2)=4-4=0$.

\textbf{3. Sonuç.} Skaler çarpım sıfır olduğundan vektörler \textbf{diktir}.
\end{cevap}

\begin{soru}[Soru 3 \foto{IMG-20261009-WA0026}]
$(1,1)$ ve $(1,0)$ vektörleri dik midir? Aralarındaki açıyı bulunuz.
\end{soru}
\begin{cevap}[Çözüm (adım adım)]
\textbf{1. Skaler çarpım.} $(1,1)\cdot(1,0)=1\cdot1+1\cdot0=1\neq0$; vektörler dik değildir.

\textbf{2. Uzunluklar.} $\|(1,1)\|=\sqrt2$, \ $\|(1,0)\|=1$.

\textbf{3. Açı.}
\[
\cos\theta=\frac{\mathbf u\cdot\mathbf v}{\|\mathbf u\|\,\|\mathbf v\|}=\frac{1}{\sqrt2}=\frac{\sqrt2}{2}\ \Rightarrow\ \theta=45^\circ .
\]
\end{cevap}

\begin{dikkat}[Defterdeki hata]
Defterde bu soruda $\cos\theta=\frac12$ yazılmıştır; doğrusu $\cos\theta=\frac{1}{\sqrt2}$'dir. Defterdeki $\theta=45^\circ$ sonucu doğrudur ($\cos\theta=\frac12$ olsaydı açı $60^\circ$ olurdu).
\end{dikkat}

%==========================================================
\subsection{Lineer bağımlılık, bağımsızlık ve taban}

\begin{soru}[Soru 4 \foto{IMG-20261009-WA0033}, \foto{IMG-20261023-WA0008}, \foto{IMG-20261023-WA0010}, \foto{IMG-20261023-WA0015}]
$\mathbf i=(1,0)$, $\mathbf j=(0,1)$ vektörlerinin lineer bağımsız olduğunu gösteriniz; $\mathbf u=(-6,9)$, $\mathbf v=(4,-6)$ vektörlerinin lineer bağımlı olup olmadığını kontrol ediniz.
\end{soru}
\begin{cevap}[Çözüm (adım adım)]
\textbf{Yöntem.} $\alpha\mathbf u+\beta\mathbf v=\mathbf 0$ denklemini kurup çözeriz. Tek çözüm $\alpha=\beta=0$ ise bağımsız, sıfırdan farklı çözüm de varsa bağımlıdır.

\textbf{(a) $\mathbf i,\mathbf j$.} $\alpha(1,0)+\beta(0,1)=(\alpha,\beta)=(0,0)$; buradan doğrudan $\alpha=0$, $\beta=0$. Tek çözüm sıfır çözüm $\Rightarrow$ $\mathbf i,\mathbf j$ \textbf{lineer bağımsız}.

\textbf{(b) $\mathbf u,\mathbf v$.} $\alpha(-6,9)+\beta(4,-6)=(0,0)$:
\[
\begin{cases}-6\alpha+4\beta=0 & /\cdot 3\\ \ \ 9\alpha-6\beta=0 & /\cdot 2\end{cases}
\ \Rightarrow\
\begin{cases}-18\alpha+12\beta=0\\ \ \ 18\alpha-12\beta=0\end{cases}
\]
Toplayınca $0=0$: ikinci denklem yeni bilgi vermez, yalnızca $3\alpha=2\beta$ koşulu kalır. Sonsuz çözüm vardır, örneğin $\alpha=2$, $\beta=3$:
\[
2(-6,9)+3(4,-6)=(-12+12,\ 18-18)=(0,0).
\]
Sıfır olmayan katsayılarla sıfır vektörü elde edildiğinden $\mathbf u,\mathbf v$ \textbf{lineer bağımlıdır}.

\textbf{Kontrol.} $\begin{vmatrix}-6&9\\4&-6\end{vmatrix}=36-36=0$; ayrıca $\mathbf v=-\frac23\mathbf u$, yani vektörler paraleldir.
\end{cevap}

\begin{soru}[Soru 5 \foto{IMG-20261009-WA0034}, \foto{IMG-20261023-WA0010}, \foto{IMG-20261023-WA0015}]
$\mathbf u=(4,-3)$, $\mathbf v=(-3,4)$, $T=\{\mathbf u,\mathbf v\}$ olsun. $T$'nin $\mathbb R^2$'nin bir tabanı olduğunu gösteriniz.
\end{soru}
\begin{cevap}[Çözüm (adım adım)]
Taban olmak için (1) lineer bağımsızlık ve (2) $\mathbb R^2$'yi germe gösterilmelidir.

\textbf{1. Bağımsızlık.} $\alpha(4,-3)+\beta(-3,4)=(0,0)$:
\[
\begin{cases}\ \ 4\alpha-3\beta=0 & /\cdot 3\\ -3\alpha+4\beta=0 & /\cdot 4\end{cases}
\Rightarrow
\begin{cases}\ \ 12\alpha-9\beta=0\\ -12\alpha+16\beta=0\end{cases}
\Rightarrow 7\beta=0\Rightarrow\beta=0,\ \alpha=0 .
\]
Yalnızca sıfır çözüm var: lineer bağımsız.

\textbf{2. Germe.} Keyfi $(x,y)\in\mathbb R^2$ için $(x,y)=\alpha\mathbf u+\beta\mathbf v$ yazılabilmeli:
\[
\begin{cases}\ \ 4\alpha-3\beta=x & /\cdot 3\\ -3\alpha+4\beta=y & /\cdot 4\end{cases}
\Rightarrow
\begin{cases}\ \ 12\alpha-9\beta=3x\\ -12\alpha+16\beta=4y\end{cases}
\Rightarrow 7\beta=3x+4y .
\]
\[
\beta=\frac{3x+4y}{7},\qquad \alpha=\frac{4x+3y}{7}\quad(\text{ilk denklemler } /\cdot4 \text{ ve } /\cdot3 \text{ ile } \beta \text{ yok edilerek}).
\]
\textbf{3. Kontrol.} $4\alpha-3\beta=\frac{16x+12y-9x-12y}{7}=x$ \checkmark, \ $-3\alpha+4\beta=\frac{-12x-9y+12x+16y}{7}=y$ \checkmark.

\textbf{4. Sonuç.} Her $(x,y)$ için $\alpha,\beta$ bulunduğundan $T$ gerer; bağımsız ve geren olduğundan $\mathbb R^2$'nin bir \textbf{tabanıdır}. (Kısa yol: $\det\begin{pmatrix}4&-3\\-3&4\end{pmatrix}=7\neq0$.)
\end{cevap}

\begin{bilgi}[Defterdeki küçük yazım hataları]
Defterde ``$T$'nin $\mathbb R$'nin bir tabanı'' yazılmıştır; doğrusu $\mathbb R^2$. $12\alpha\ \_\ 9\beta$ satırındaki işaret silik görünmektedir; doğrusu $12\alpha-9\beta=0$.
\end{bilgi}

%==========================================================
\subsection{Uzayda koordinatlar ve küre}

\begin{soru}[Soru 6 \foto{IMG-20261009-WA0028}, \foto{IMG-20261023-WA0013}]
Denklemi $x^2+y^2+z^2-2x-4y-6z-2=0$ olan kürenin merkezini ve yarıçapını bulunuz. \okunamadi{23.10 tahtasında soru cümlesinin sonu kenardan taşmış; istenen, çözümden anlaşılmaktadır}
\end{soru}
\begin{cevap}[Çözüm (adım adım)]
\textbf{1. Genel denklemle karşılaştırma.} $(x-a)^2+(y-b)^2+(z-c)^2=r^2$ açılırsa
$x^2+y^2+z^2-2ax-2by-2cz+(a^2+b^2+c^2-r^2)=0$ elde edilir. Katsayıları eşleyelim:
\[
-2a=-2,\quad -2b=-4,\quad -2c=-6\ \Rightarrow\ a=1,\ b=2,\ c=3,\qquad M(1,2,3).
\]
\textbf{2. Yarıçap.} Sabit terim: $a^2+b^2+c^2-r^2=-2\Rightarrow 14-r^2=-2\Rightarrow r^2=16\Rightarrow r=4$.

\textbf{3. Kontrol (tam kare).} $(x-1)^2+(y-2)^2+(z-3)^2=2+1+4+9=16$ \checkmark.
\end{cevap}

%==========================================================
\subsection{Vektörel çarpım ve karma çarpım}

\begin{soru}[Soru 7 \foto{IMG-20261009-WA0035}, \foto{IMG-20261023-WA0029}, \foto{IMG-20261023-WA0014}]
$\mathbf u=(1,2,3)$, $\mathbf v=(2,0,1)$ için $\mathbf u\times\mathbf v$'yi bulunuz. Sonucun $\mathbf u$'ya dik olduğunu gösteriniz.
\end{soru}
\begin{cevap}[Çözüm (adım adım)]
\textbf{1. Determinantı kurma.} İlk satır $\mathbf i,\mathbf j,\mathbf k$; ikinci satır $\mathbf u$; üçüncü satır $\mathbf v$:
\[
\mathbf u\times\mathbf v=\begin{vmatrix}\mathbf i&\mathbf j&\mathbf k\\1&2&3\\2&0&1\end{vmatrix}.
\]
\textbf{2. İlk satıra göre açılım} (işaretler $+,-,+$):
\[
=\mathbf i\begin{vmatrix}2&3\\0&1\end{vmatrix}-\mathbf j\begin{vmatrix}1&3\\2&1\end{vmatrix}+\mathbf k\begin{vmatrix}1&2\\2&0\end{vmatrix}
=\mathbf i\,(2-0)-\mathbf j\,(1-6)+\mathbf k\,(0-4).
\]
\textbf{3. Sonuç.} $\mathbf u\times\mathbf v=2\mathbf i+5\mathbf j-4\mathbf k=(2,5,-4)$.

\textbf{4. Diklik.} $\mathbf u\cdot(\mathbf u\times\mathbf v)=(1,2,3)\cdot(2,5,-4)=2+10-12=0$, dolayısıyla $\mathbf u\times\mathbf v\perp\mathbf u$. Benzer biçimde $\mathbf v\cdot(\mathbf u\times\mathbf v)=4+0-4=0$.

\textbf{5. Ek bilgi.} $\|\mathbf u\times\mathbf v\|=\sqrt{4+25+16}=\sqrt{45}=3\sqrt5$; bu sayı $\mathbf u$ ve $\mathbf v$ üzerine kurulan paralelkenarın alanıdır.
\end{cevap}

\begin{soru}[Soru 8 \foto{IMG-20261009-WA0027}, \foto{IMG-20261023-WA0014}]
$\mathbf i,\mathbf j,\mathbf k$ vektörlerinin karma çarpımını bulunuz ve geometrik anlamını yorumlayınız.
\end{soru}
\begin{cevap}[Çözüm (adım adım)]
\textbf{1. Determinant.}
\[
\mathbf i\cdot(\mathbf j\times\mathbf k)=\begin{vmatrix}1&0&0\\0&1&0\\0&0&1\end{vmatrix}=1\cdot(1\cdot1-0\cdot0)-0+0=1 .
\]
\textbf{2. Doğrudan.} $\mathbf j\times\mathbf k=\mathbf i$ olduğundan $\mathbf i\cdot\mathbf i=1$.

\textbf{3. Yorum.} Teoreme göre $|\mathbf u\cdot(\mathbf v\times\mathbf w)|$, üç vektör üzerine kurulan paralelyüzlünün hacmidir. $\mathbf i,\mathbf j,\mathbf k$ üzerine kurulan paralelyüzlü, kenarları 1 olan \textbf{birim küptür} ve hacmi 1'dir. Sonucun sıfırdan farklı olması üç vektörün aynı düzlemde olmadığını, pozitif olması sağ el sistemi oluşturduklarını gösterir.
\end{cevap}

%==========================================================
\subsection{Skaler çarpım uygulaması: iki doğru arasındaki açı}

\begin{soru}[Soru 9 \foto{IMG-20261002-WA0018}]
$\dfrac{x-1}{1}=\dfrac{y}{-\sqrt2}=\dfrac{z-1}{1}$ ve $\dfrac{x}{-1}=\dfrac{y-1}{\sqrt2}=\dfrac{z+2}{1}$ doğrularının oluşturduğu açıyı bulunuz.
\end{soru}
\begin{cevap}[Çözüm (adım adım)]
\textbf{1. Doğrultman vektörleri.} $\frac{x-x_0}{a}=\frac{y-y_0}{b}=\frac{z-z_0}{c}$ biçimindeki bir doğrunun doğrultmanı $(a,b,c)$'dir (paydalar):
\[
\mathbf v=(1,-\sqrt2,1),\qquad \mathbf u=(-1,\sqrt2,1).
\]
\textbf{2. Skaler çarpım ve uzunluklar.}
$\mathbf u\cdot\mathbf v=-1-2+1=-2$, \ $\|\mathbf u\|=\|\mathbf v\|=\sqrt{1+2+1}=2$.

\textbf{3. Açı.} $\cos\theta=\dfrac{-2}{2\cdot2}=-\dfrac12\Rightarrow\theta=120^\circ$.

\textbf{4. Dar açı.} Kesişen iki doğru $120^\circ$ ve $60^\circ$'lik açılar oluşturur; doğrular arasındaki açı dar olan alınır: $180^\circ-120^\circ=\boxed{60^\circ}$.
\end{cevap}

%==========================================================
\subsection{Çalışma soruları}

\begin{soru}[Çalışma sorusu 1 (derste çalışma sorusu olarak verildi) \foto{IMG-20261009-WA0032}, \foto{IMG-20261023-WA0027}]
$(2,-1,3)$ ile $(1,4,2)$ vektörlerinin dik olup olmadığını bulunuz ve $\cos\theta$'yı hesaplayınız.
\end{soru}
\begin{cevap}[Çözüm (adım adım)]
\textbf{1.} $(2,-1,3)\cdot(1,4,2)=2\cdot1+(-1)\cdot4+3\cdot2=2-4+6=4\neq0$; vektörler \textbf{dik değildir}.

\textbf{2.} $\|(2,-1,3)\|=\sqrt{4+1+9}=\sqrt{14}$, \ $\|(1,4,2)\|=\sqrt{1+16+4}=\sqrt{21}$.

\textbf{3.}
\[
\cos\theta=\frac{4}{\sqrt{14}\cdot\sqrt{21}}=\frac{4}{\sqrt{294}}=\frac{4}{7\sqrt6}=\frac{4\sqrt6}{42}=\frac{2\sqrt6}{21}\approx0{,}233,\qquad \theta\approx76{,}5^\circ .
\]
Skaler çarpım pozitif olduğu için açı dardır.
\end{cevap}

\begin{soru}[Çalışma sorusu 2 (derste çalışma sorusu olarak verildi) \foto{IMG-20261009-WA0032}, \foto{IMG-20261023-WA0027}]
$x^2+y^2+z^2+4x-6y+2z-11=0$ kürenin merkezini ve yarıçapını bulunuz.
\end{soru}
\begin{cevap}[Çözüm (adım adım)]
\textbf{1. Tam kareye tamamlama.}
\[
(x^2+4x+4)+(y^2-6y+9)+(z^2+2z+1)=11+4+9+1 .
\]
\textbf{2.} $(x+2)^2+(y-3)^2+(z+1)^2=25$.

\textbf{3. Sonuç.} Merkez $M(-2,3,-1)$, yarıçap $r=5$.

\textbf{Kontrol (formül).} $D=4$, $E=-6$, $F=2$, $G=-11$: $M=\left(-\frac42,\frac62,-\frac22\right)=(-2,3,-1)$, \ $r^2=4+9+1-(-11)=25$ \checkmark.
\end{cevap}

\begin{soru}[Çalışma sorusu 3 (derste çalışma sorusu olarak verildi) \foto{IMG-20261009-WA0032}, \foto{IMG-20261023-WA0027}]
$A(0,1,2)$ ve $B(3,1,-1)$ noktalarından geçen doğrunun parametrik ve kartezyen denklemini bulunuz.
\end{soru}
\begin{cevap}[Çözüm (adım adım)]
Bu soru, vektörleri Konu 2'deki doğru denklemlerine bağlar.

\textbf{1. Doğrultman.} $\overrightarrow{AB}=(3-0,\ 1-1,\ -1-2)=(3,0,-3)$. Daha sade olarak buna paralel $(1,0,-1)$ de alınabilir.

\textbf{2. Vektörel denklem.} Doğru üzerindeki $P(x,y,z)$ için $\overrightarrow{AP}=t\,\overrightarrow{AB}$, $t\in\mathbb R$:
$(x,\ y-1,\ z-2)=t\,(3,0,-3)$.

\textbf{3. Parametrik denklem.}
\[
x=3t,\qquad y=1,\qquad z=2-3t\qquad(t\in\mathbb R).
\]
(Kontrol: $t=0\Rightarrow A$, $t=1\Rightarrow(3,1,-1)=B$ \checkmark.)

\textbf{4. Kartezyen denklem.} Doğrultmanın ikinci bileşeni $b=0$ olduğundan $y$ için kesir yazılamaz; doğru $Oy$ eksenine diktir ve
\[
\frac{x}{3}=\frac{z-2}{-3},\qquad y=1
\]
olur. Sadeleştirirsek $x+z=2$, $y=1$.
\end{cevap}

%==========================================================
\subsection{Ek alıştırmalar}

\begin{soru}[\ekal Ek alıştırma 1]
Merkezi $M(1,2,-2)$ olan ve orijinden geçen kürenin denklemini yazınız. $Q(3,3,0)$ noktası kürenin içinde mi, üzerinde mi, dışında mı?
\end{soru}
\begin{cevap}
\textbf{1.} Küre orijinden geçtiğinden $r=|MO|=\sqrt{1+4+4}=3$.
\[
(x-1)^2+(y-2)^2+(z+2)^2=9\quad\Longleftrightarrow\quad x^2+y^2+z^2-2x-4y+4z=0 .
\]
\textbf{2.} $|MQ|^2=(3-1)^2+(3-2)^2+(0+2)^2=4+1+4=9=r^2$. $Q$ küre \textbf{üzerindedir}. (Genel kural: $|MQ|<r$ içte, $|MQ|>r$ dışta.)
\end{cevap}

\begin{soru}[\ekal Ek alıştırma 2]
$\mathbf u=(x,2,-1)$ ve $\mathbf v=(3,x,4)$ olsun. (a) Hangi $x$ için $\mathbf u\perp\mathbf v$'dir? (b) $\mathbf u\parallel\mathbf v$ olacak bir $x$ var mıdır?
\end{soru}
\begin{cevap}
(a) $\mathbf u\cdot\mathbf v=3x+2x-4=5x-4=0\Rightarrow x=\dfrac45$.

(b) Paralellik için $\mathbf u=\lambda\mathbf v$ olmalı: $x=3\lambda$, $2=\lambda x$, $-1=4\lambda$. Son denklemden $\lambda=-\frac14$, ilkinden $x=-\frac34$; ikinci denklem $2=\left(-\frac14\right)\left(-\frac34\right)=\frac{3}{16}$ olur ki bu yanlıştır. Dolayısıyla hiçbir $x$ için paralel değildir.
\end{cevap}

\begin{soru}[\ekal Ek alıştırma 3]
$\mathbf a=(4,0,-2)$ vektörünün $\mathbf b=(2,1,2)$ üzerine skaler ve vektörel izdüşümünü bulunuz; $\mathbf a$'yı $\mathbf b$'ye paralel ve dik iki vektörün toplamı olarak yazınız.
\end{soru}
\begin{cevap}
\textbf{1.} $\mathbf a\cdot\mathbf b=8+0-4=4$, \ $\|\mathbf b\|=\sqrt{4+1+4}=3$.

\textbf{2.} $\text{comp}_{\mathbf b}\mathbf a=\dfrac{4}{3}$, \quad $\text{proj}_{\mathbf b}\mathbf a=\dfrac{4}{9}(2,1,2)=\left(\dfrac89,\dfrac49,\dfrac89\right)$.

\textbf{3.} Dik bileşen: $\mathbf a-\text{proj}_{\mathbf b}\mathbf a=\left(\dfrac{28}{9},-\dfrac49,-\dfrac{26}{9}\right)$. Kontrol: $\frac{56-4-52}{9}=0$ \checkmark.
\[
\mathbf a=\left(\tfrac89,\tfrac49,\tfrac89\right)+\left(\tfrac{28}{9},-\tfrac49,-\tfrac{26}{9}\right).
\]
\end{cevap}

\begin{soru}[\ekal Ek alıştırma 4]
$\mathbf a=(1,2,-1)$ ve $\mathbf b=(2,0,3)$ üzerine kurulan paralelkenarın ve bu vektörlerin belirttiği üçgenin alanını bulunuz.
\end{soru}
\begin{cevap}
\[
\mathbf a\times\mathbf b=\begin{vmatrix}\mathbf i&\mathbf j&\mathbf k\\1&2&-1\\2&0&3\end{vmatrix}
=\mathbf i\,(6-0)-\mathbf j\,(3+2)+\mathbf k\,(0-4)=(6,-5,-4).
\]
Kontrol: $(1,2,-1)\cdot(6,-5,-4)=6-10+4=0$ \checkmark, \ $(2,0,3)\cdot(6,-5,-4)=12-12=0$ \checkmark.

Paralelkenar alanı $=\|(6,-5,-4)\|=\sqrt{36+25+16}=\sqrt{77}$; üçgen alanı $=\dfrac{\sqrt{77}}{2}$.
\end{cevap}

\begin{soru}[\ekal Ek alıştırma 5]
$\mathbf a=(1,2,3)$, $\mathbf b=(2,k,1)$, $\mathbf c=(1,1,1)$ vektörlerinin aynı düzlemde olması için $k$ ne olmalıdır?
\end{soru}
\begin{cevap}
Düzlemdeşlik $\iff$ karma çarpım sıfır:
\[
\begin{vmatrix}1&2&3\\2&k&1\\1&1&1\end{vmatrix}
=1\,(k-1)-2\,(2-1)+3\,(2-k)=k-1-2+6-3k=3-2k .
\]
$3-2k=0\Rightarrow k=\dfrac32$. Diğer $k$ değerleri için üç vektör bir paralelyüzlü oluşturur; hacmi $|3-2k|$'dir.
\end{cevap}

\begin{soru}[\ekal Ek alıştırma 6]
$\mathbf a=(1,2,3)$, $\mathbf b=(4,5,6)$, $\mathbf c=(7,8,9)$ vektörleri $\mathbb R^3$'ün bir tabanı mıdır? Değilse aralarındaki lineer bağıntıyı bulunuz.
\end{soru}
\begin{cevap}
\[
\begin{vmatrix}1&2&3\\4&5&6\\7&8&9\end{vmatrix}=1\,(45-48)-2\,(36-42)+3\,(32-35)=-3+12-9=0 .
\]
Determinant sıfır: vektörler lineer bağımlıdır, dolayısıyla taban \textbf{değildir} (hepsi aynı düzlemdedir).

Bağıntı: $\alpha\mathbf a+\beta\mathbf b+\gamma\mathbf c=\mathbf 0$ için ilk iki bileşen
$\alpha+4\beta+7\gamma=0$, $2\alpha+5\beta+8\gamma=0$. İlkinin 2 katını ikinciden çıkarırsak $-3\beta-6\gamma=0\Rightarrow\beta=-2\gamma$; ilkinden $\alpha=-4\beta-7\gamma=\gamma$. $\gamma=1$ alırsak $\alpha=1$, $\beta=-2$:
\[
\mathbf a-2\mathbf b+\mathbf c=(1-8+7,\ 2-10+8,\ 3-12+9)=(0,0,0)\ \checkmark
\]
\end{cevap}

\begin{soru}[\ekal Ek alıştırma 7]
Her $\mathbf u,\mathbf v$ için $\|\mathbf u+\mathbf v\|^2+\|\mathbf u-\mathbf v\|^2=2\left(\|\mathbf u\|^2+\|\mathbf v\|^2\right)$ olduğunu (paralelkenar kuralı) gösteriniz. Geometrik anlamı nedir?
\end{soru}
\begin{cevap}
$\|\mathbf w\|^2=\mathbf w\cdot\mathbf w$ ve skaler çarpımın dağılma özelliğini kullanalım:
\begin{align*}
\|\mathbf u+\mathbf v\|^2&=\mathbf u\cdot\mathbf u+2\,\mathbf u\cdot\mathbf v+\mathbf v\cdot\mathbf v,\\
\|\mathbf u-\mathbf v\|^2&=\mathbf u\cdot\mathbf u-2\,\mathbf u\cdot\mathbf v+\mathbf v\cdot\mathbf v .
\end{align*}
Toplarsak $\mathbf u\cdot\mathbf v$ terimleri sadeleşir: $2\|\mathbf u\|^2+2\|\mathbf v\|^2$.

\textbf{Geometrik anlam:} Bir paralelkenarın köşegen uzunluklarının kareleri toplamı, dört kenar uzunluğunun kareleri toplamına eşittir. ($\mathbf u+\mathbf v$ ve $\mathbf u-\mathbf v$ köşegenlerdir.) $\mathbf u\perp\mathbf v$ ise iki köşegen eşit olur; bu dikdörtgen durumudur.
\end{cevap}

\begin{soru}[\ekal Ek alıştırma 8]
$\mathbf u=(1,1,1)$ vektörünün koordinat eksenleriyle yaptığı açıları bulunuz.
\end{soru}
\begin{cevap}
$\|\mathbf u\|=\sqrt3$; $\cos\alpha=\cos\beta=\cos\gamma=\dfrac{1}{\sqrt3}$. Kontrol: $3\cdot\frac13=1$ \checkmark.
\[
\alpha=\beta=\gamma=\arccos\frac{1}{\sqrt3}\approx54{,}74^\circ .
\]
Bu, bir küpün köşegeninin kenarlarıyla yaptığı açıdır.
\end{cevap}

\clearpage\section{Konu 2: Doğru ve Düzlem Denklemleri, Silindirler ve Kuadratik Yüzeyler}

Bu bölümde Konu~2 kapsamında derste (01.10, 02.10, 09.10, 15.10, 22.10, 23.10 ve 08.11 tarihli tahta/defter fotoğrafları) çözülen bütün örnekler, çözümsüz bırakılan çalışma soruları ve kaynak kitaplara benzer ek alıştırmalar adım adım çözülmüştür. Aynı örnek birden çok derste tekrarlandıysa bir kez yazılmış, tüm fotoğraf kaynakları verilmiştir. Sorular, konunun işleniş sırasına (doğru $\to$ düzlem $\to$ doğru--düzlem $\to$ yüzeyler) göre dizilmiştir.

\begin{bilgi}[Kullanılan temel formüller]
Doğru: $\frac{x-x_0}{a}=\frac{y-y_0}{b}=\frac{z-z_0}{c}$, $\mathbf v=(a,b,c)$.\quad
Düzlem: $Ax+By+Cz+D=0$, $\mathbf n=(A,B,C)$.\quad
Nokta--doğru: $\ell=\frac{\|\overrightarrow{AP}\times\mathbf v\|}{\|\mathbf v\|}$.\quad
Nokta--düzlem: $d=\frac{|Ax_0+By_0+Cz_0+D|}{\sqrt{A^2+B^2+C^2}}$.
\end{bilgi}

%======================================================================
\subsection{Uzayda doğru denklemleri}

\begin{soru}[Soru 1 \foto{IMG-20261009-WA0030}, \foto{IMG-20261023-WA0018}]
$A(1,2,3)$ noktasından geçen ve $\mathbf v=(2,3,0)$ vektörüne paralel olan doğrunun denklemini yazınız.
\end{soru}
\begin{cevap}[Çözüm (adım adım)]
\textbf{1. Verileri okuyalım.} Nokta $(x_0,y_0,z_0)=(1,2,3)$, doğrultman parametreleri $a=2$, $b=3$, $c=0$.

\textbf{2. Parametrik denklemler.} $x=1+2t,\quad y=2+3t,\quad z=3+0\cdot t=3$.

\textbf{3. Kartezyen biçim.} $c=0$ olduğundan $\frac{z-3}{0}$ yazılamaz; $z$ sabittir. İlk iki denklemden $t$ çekilip eşitlenir:
\[ \boxed{\frac{x-1}{2}=\frac{y-2}{3},\qquad z=3}. \]
\textbf{4. Sadeleştirme.} $3(x-1)=2(y-2)\Rightarrow3x-2y+1=0$. Doğru $\{3x-2y+1=0,\ z=3\}$ sistemiyle verilir.

\textit{Not:} Tahtada sonuç yalnız $3x-2y+1=0$ olarak kutulanmıştı; bu tek başına $Oz$'ye paralel bir düzlemdir. Doğru için $z=3$ de mutlaka yazılmalıdır.
\end{cevap}

\begin{soru}[Soru 2 \foto{IMG-20261009-WA0030}, \foto{IMG-20261023-WA0017}]
$A(1,3,-1)$ ve $B(3,5,6)$ noktalarından geçen doğrunun denklemini bulunuz.
\end{soru}
\begin{cevap}[Çözüm (adım adım)]
\textbf{1. Doğrultman.} $\overrightarrow{AB}=(3-1,\ 5-3,\ 6-(-1))=(2,2,7)$.

\textbf{2. Denklem.} $A$ noktası ve $\overrightarrow{AB}$ ile:
\[ \frac{x-1}{3-1}=\frac{y-3}{5-3}=\frac{z-(-1)}{6-(-1)}\ \Longrightarrow\ \boxed{\frac{x-1}{2}=\frac{y-3}{2}=\frac{z+1}{7}} \]
\textbf{3. Kontrol.} $B$'yi yerine koyalım: $\frac{3-1}{2}=1$, $\frac{5-3}{2}=1$, $\frac{6+1}{7}=1$; üç oran eşit \checkmark. Parametrik biçim $x=1+2t$, $y=3+2t$, $z=-1+7t$.
\end{cevap}

\begin{soru}[Soru 3 \foto{IMG-20261009-WA0030}, \foto{IMG-20261023-WA0017}]
$\dfrac{x-1}{3}=\dfrac{y+2}{m}=\dfrac{z-2}{4}$ ve $\dfrac{x-2}{6}=\dfrac{y-3}{8}=\dfrac{z+1}{n}$ doğruları paralel ise $m$ ve $n$ kaçtır?
\end{soru}
\begin{cevap}[Çözüm (adım adım)]
\textbf{1. Doğrultmanlar.} $\mathbf v=(3,m,4)$, $\mathbf u=(6,8,n)$.

\textbf{2. Paralellik şartı.} Doğrultmanlar orantılı olmalı: $\dfrac36=\dfrac m8=\dfrac4n$. İlk oran bilinen: $\frac36=\frac12$.

\textbf{3. Çözüm.} $\frac m8=\frac12\Rightarrow m=4$;\quad $\frac4n=\frac12\Rightarrow n=8$. Yani $\boxed{m=4,\ n=8}$.

\textbf{4. Çakışık mı?} $d_1$'in noktası $(1,-2,2)$: $\frac{1-2}{6}=-\frac16$, $\frac{-2-3}{8}=-\frac58$; oranlar eşit değil, nokta $d_2$'de değil. Doğrular paralel ve farklıdır.
\end{cevap}

%======================================================================
\subsection{İki doğru arasındaki açı ve noktanın doğruya uzaklığı}

\begin{soru}[Soru 4 \foto{IMG-20261002-WA0018}, \foto{IMG-20261023-WA0019}]
$\dfrac{x-1}{1}=\dfrac{y}{-\sqrt2}=\dfrac{z-1}{1}$ ve $\dfrac{x}{-1}=\dfrac{y-1}{\sqrt2}=\dfrac{z+2}{1}$ doğrularının oluşturduğu açıyı bulunuz.
\end{soru}
\begin{cevap}[Çözüm (adım adım)]
\textbf{1. Doğrultmanlar.} Paydalardan $\mathbf v=(1,-\sqrt2,1)$, $\mathbf u=(-1,\sqrt2,1)$.

\textbf{2. Skaler çarpım ve normlar.}
$\mathbf u\cdot\mathbf v=-1+(-\sqrt2)(\sqrt2)+1=-1-2+1=-2$;\quad $\|\mathbf u\|=\|\mathbf v\|=\sqrt{1+2+1}=2$.

\textbf{3. Açı.} $\cos\theta=\dfrac{-2}{2\cdot2}=-\dfrac12\Rightarrow\theta=120^\circ$ (doğrultman vektörleri arasındaki açı).

\textbf{4. Doğrular arasındaki açı.} Kesişen iki doğru $120^\circ$ ve $180^\circ-120^\circ=60^\circ$'lik açılar oluşturur; doğrular arasındaki açı dar açı olarak verilir: $\boxed{60^\circ}$. (Mutlak değerli formülle $\cos\theta=\frac{|-2|}{4}=\frac12$ doğrudan $60^\circ$ verir.)
\end{cevap}

\begin{soru}[Soru 5 \foto{IMG-20261002-WA0019}, \foto{IMG-20261023-WA0020}, \foto{IMG-20261023-WA0057}]
$P(2,3,-1)$ noktasının $\dfrac{x-1}{1}=\dfrac{y-2}{1}=\dfrac{z-13}{4}$ doğrusuna uzaklığını bulunuz.
\end{soru}
\begin{cevap}[Çözüm (adım adım)]
\textbf{1. Formül.} $A\in d$, $\mathbf v$ doğrultman ise $\ell=\dfrac{\|\overrightarrow{AP}\times\mathbf v\|}{\|\mathbf v\|}$ (paralelkenar alanı / taban).

\textbf{2. Nokta ve doğrultman.} $A(1,2,13)$, $\mathbf v=(1,1,4)$.

\textbf{3. $\overrightarrow{AP}$.} $(2-1,\ 3-2,\ -1-13)=(1,1,-14)$.

\textbf{4. Vektörel çarpım.}
\[ \overrightarrow{AP}\times\mathbf v=\begin{vmatrix}\mathbf i&\mathbf j&\mathbf k\\1&1&-14\\1&1&4\end{vmatrix}
=\mathbf i\,(1\cdot4-(-14)\cdot1)-\mathbf j\,(1\cdot4-(-14)\cdot1)+\mathbf k\,(1\cdot1-1\cdot1)=(18,-18,0). \]
\textbf{5. Uzaklık.}
\[ \ell=\frac{\sqrt{18^2+(-18)^2+0}}{\sqrt{1+1+16}}=\frac{18\sqrt2}{\sqrt{18}}=\frac{18\sqrt2}{3\sqrt2}=\boxed{6}. \]
\textit{Not:} Defterde önce $A(2,3,17)$ yazılmıştı; bu da doğrunun noktasıdır ($t=1$) ve aynı sonucu verir: $\overrightarrow{AP}=(0,0,-18)$, $\overrightarrow{AP}\times\mathbf v=(18,-18,0)$.
\end{cevap}

%======================================================================
\subsection{Düzlem denklemleri}

\begin{soru}[Soru 6 \foto{IMG-20261002-WA0020}, \foto{IMG-20261023-WA0030}]
$A(2,3,-5)$ noktasından geçen ve $\mathbf n=(7,3,4)$ vektörüne dik olan düzlemin denklemini bulunuz.
\end{soru}
\begin{cevap}[Çözüm (adım adım)]
\textbf{1.} Düzlem $\mathbf n$'ye dik $\Rightarrow$ $\mathbf n$ düzlemin normalidir: $A=7,\ B=3,\ C=4$.

\textbf{2.} $A(x-x_0)+B(y-y_0)+C(z-z_0)=0$:\quad $7(x-2)+3(y-3)+4(z+5)=0$.

\textbf{3.} $7x-14+3y-9+4z+20=0\Rightarrow\boxed{7x+3y+4z-3=0}$.

\textbf{4. Kontrol.} $14+9-20-3=0$ \checkmark.
\end{cevap}

\begin{soru}[Soru 7 \foto{IMG-20261002-WA0024}, \foto{IMG-20261002-WA0021}, \foto{IMG-20261023-WA0032}]
$P(2,3,5)$ noktasından geçen ve $\dfrac{x-1}{4}=\dfrac{y-2}{-1}=\dfrac{z+1}{6}$ doğrusuna dik olan düzlemin denklemini bulunuz.
\end{soru}
\begin{cevap}[Çözüm (adım adım)]
\textbf{1. Fikir.} Düzlem doğruya dik ise doğrunun doğrultmanı düzleme diktir, yani normale paraleldir: $\mathbf n=\mathbf v$.

\textbf{2.} $\mathbf v=(4,-1,6)=\mathbf n$, nokta $(2,3,5)$.

\textbf{3.} $4(x-2)-1(y-3)+6(z-5)=0\Rightarrow4x-8-y+3+6z-30=0\Rightarrow\boxed{4x-y+6z-35=0}$.

\textbf{4. Kontrol.} $8-3+30-35=0$ \checkmark.

\textit{Tahtadaki hata:} 23.10 tahtasında sonuç $4x-y+6z-3=0$ yazılmıştı. Sabit terim $-8+3-30=-35$'tir; $-3$ ile $P$ noktası denklemi sağlamaz ($8-3+30-3=32\neq0$). 02.10 tahtasında doğru sonuç ($-35$) yazılmıştı.
\end{cevap}

\begin{soru}[Soru 8 \foto{IMG-20261002-WA0022}, \foto{IMG-20261023-WA0050}]
$A(1,3,5)$, $B(4,2,1)$, $C(-1,2,-3)$ noktalarından geçen düzlemin denklemini bulunuz.
\end{soru}
\begin{cevap}[Çözüm (adım adım)]
\textbf{1. Fikir.} $P(x,y,z)$ düzlemde $\iff$ $\overrightarrow{AP},\overrightarrow{AB},\overrightarrow{AC}$ aynı düzlemde $\iff$ karma çarpım (determinant) $0$.

\textbf{2. Vektörler.} $\overrightarrow{AB}=(3,-1,-4)$, $\overrightarrow{AC}=(-2,-1,-8)$, $\overrightarrow{AP}=(x-1,y-3,z-5)$.

\textbf{3. Determinant.}
\[ \begin{vmatrix}x-1&y-3&z-5\\3&-1&-4\\-2&-1&-8\end{vmatrix}=0 . \]
Birinci satıra göre açılım (eş çarpanlar):
\begin{align*}
(x-1)\cdot\big[(-1)(-8)-(-4)(-1)\big]&=(x-1)(8-4)=4(x-1),\\
-(y-3)\cdot\big[3(-8)-(-4)(-2)\big]&=-(y-3)(-24-8)=32(y-3),\\
(z-5)\cdot\big[3(-1)-(-1)(-2)\big]&=(z-5)(-3-2)=-5(z-5).
\end{align*}
\textbf{4. Toplam.} $4(x-1)+32(y-3)-5(z-5)=0\Rightarrow4x+32y-5z-4-96+25=0\Rightarrow\boxed{4x+32y-5z-75=0}$.

\textbf{5. Sağlama.} $A$: $4+96-25-75=0$; $B$: $16+64-5-75=0$; $C$: $-4+64+15-75=0$ \checkmark.
\end{cevap}

%======================================================================
\subsection{İki düzlemin durumu ve noktanın düzleme uzaklığı}

\begin{soru}[Soru 9 \foto{IMG-20261002-WA0022}, \foto{IMG-20261023-WA0024}]
$2x+my+3z+4=0$ ve $-nx+4y+z+2=0$ düzlemleri birbirine dik ise $m$ ile $n$ arasındaki bağıntıyı bulunuz.
\end{soru}
\begin{cevap}[Çözüm (adım adım)]
\textbf{1.} Normaller $\mathbf n_1=(2,m,3)$, $\mathbf n_2=(-n,4,1)$.

\textbf{2.} Düzlemler dik $\iff$ normaller dik $\iff$ $A_1A_2+B_1B_2+C_1C_2=0$:
\[ 2(-n)+m\cdot4+3\cdot1=0\ \Longrightarrow\ \boxed{4m-2n+3=0}. \]
Bu bir tek çözüm değil, $m$ ile $n$ arasında bir bağıntıdır: $n=2m+\frac32$. Örneğin $m=1$ için $n=\frac72$: $(2,1,3)\cdot(-\frac72,4,1)=-7+4+3=0$ \checkmark.
\end{cevap}

\begin{soru}[Soru 10 \foto{IMG-20261002-WA0023}, \foto{IMG-20261023-WA0023}]
\textbf{(a)} $x+y+z=0$ ile $2x+2y+2z-5=0$ düzlemlerinin paralel mi, çakışık mı olduğunu belirleyiniz.
\textbf{(b)} $x+y+z=0$ ile $x-y=0$ düzlemlerinin dik olup olmadığını belirleyiniz.
\end{soru}
\begin{cevap}[Çözüm (adım adım)]
\textbf{(a)} $A_1=B_1=C_1=1$, $D_1=0$; $A_2=B_2=C_2=2$, $D_2=-5$.
\[ \frac{A_1}{A_2}=\frac{B_1}{B_2}=\frac{C_1}{C_2}=\frac12\ \Rightarrow\ \text{paralel};\qquad \frac{D_1}{D_2}=\frac{0}{-5}=0\neq\frac12\ \Rightarrow\ \text{çakışık değil.} \]
Düzlemler paralel ve farklıdır. (Aralarındaki uzaklık: ikinciyi $x+y+z-\frac52=0$ yazarsak $d=\frac{|{-\frac52}-0|}{\sqrt3}=\frac{5}{2\sqrt3}$.)

\textit{Not:} 02.10 tahtasında ikinci düzlem $2x+2y+2z=5$ biçiminde yazılıp kontrol $\frac12\neq\frac05$ diye yapılmıştı; $Ax+By+Cz+D=0$ biçiminde $D_2=-5$'tir. Sonuç değişmez.

\textbf{(b)} $\mathbf n_1=(1,1,1)$, $\mathbf n_2=(1,-1,0)$: $\mathbf n_1\cdot\mathbf n_2=1-1+0=0$ $\Rightarrow$ düzlemler \textbf{diktir}.
\end{cevap}

\begin{soru}[Soru 11 \foto{IMG-20261002-WA0023}, \foto{IMG-20261023-WA0025}, \foto{IMG-20261023-WA0031}]
$P(2,3,4)$ noktasının $3x-4y+12z+10=0$ düzlemine uzaklığını bulunuz.
\end{soru}
\begin{cevap}[Çözüm (adım adım)]
\textbf{1.} $A=3$, $B=-4$, $C=12$, $D=10$; $(x_0,y_0,z_0)=(2,3,4)$.

\textbf{2. Pay.} $|3\cdot2+(-4)\cdot3+12\cdot4+10|=|6-12+48+10|=52$.

\textbf{3. Payda.} $\sqrt{3^2+(-4)^2+12^2}=\sqrt{9+16+144}=\sqrt{169}=13$.

\textbf{4.} $d=\dfrac{52}{13}=\boxed{4}$.
\end{cevap}

\begin{soru}[Soru 12 \foto{IMG-20261002-WA0026}, \foto{IMG-20261023-WA0056}]
$P_1:\ Ax+By+Cz+D_1=0$ ve $P_2:\ Ax+By+Cz+D_2=0$ paralel düzlemleri arasındaki uzaklığın $d=\dfrac{|D_2-D_1|}{\sqrt{A^2+B^2+C^2}}$ olduğunu gösteriniz.
\end{soru}
\begin{cevap}[Çözüm (adım adım)]
\textbf{1.} Paralel düzlemlerde $P_1$'in her noktası $P_2$'ye aynı uzaklıktadır; $P_1$'den bir $K(x_0,y_0,z_0)$ alalım.

\textbf{2.} Noktanın düzleme uzaklığı: $d=\dfrac{|Ax_0+By_0+Cz_0+D_2|}{\sqrt{A^2+B^2+C^2}}$.

\textbf{3.} $K\in P_1$ olduğundan $Ax_0+By_0+Cz_0=-D_1$. Yerine yazarsak
\[ d=\frac{|-D_1+D_2|}{\sqrt{A^2+B^2+C^2}}=\frac{|D_2-D_1|}{\sqrt{A^2+B^2+C^2}} . \]
\textit{Uyarı:} Formül yalnızca iki denklem \emph{aynı} $A,B,C$ katsayılarıyla yazıldığında geçerlidir.
\end{cevap}

%======================================================================
\subsection{Doğru ile düzlem; arakesit doğrusu}

\begin{soru}[Soru 13 \foto{IMG-20261023-WA0023}, \foto{IMG-20261023-WA0025}]
$d:\ \dfrac{x-1}{2}=\dfrac{y}{3}=\dfrac{z+1}{-1}$ doğrusunun
\textbf{(a)} $4x+6y-2z+5=0$ düzlemine,
\textbf{(b)} $x-y-z+7=0$ düzlemine göre durumunu inceleyiniz.
\end{soru}
\begin{cevap}[Çözüm (adım adım)]
Doğrunun doğrultmanı $\mathbf v=(2,3,-1)$, üzerindeki bir nokta $(1,0,-1)$.

\textbf{(a)} $\mathbf n=(4,6,-2)$. Diklik testi: $\frac24=\frac36=\frac{-1}{-2}=\frac12$; oranlar eşit $\Rightarrow$ $\mathbf v\parallel\mathbf n$ $\Rightarrow$ \textbf{doğru düzleme diktir}.

\textbf{(b)} $\mathbf n=(1,-1,-1)$. Oranlar $\frac21\neq\frac3{-1}$, dik değil. Paralellik testi: $\mathbf n\cdot\mathbf v=2-3+1=0$ $\Rightarrow$ $\mathbf v\perp\mathbf n$: doğru düzleme paralel \emph{ya da} düzlemin içinde. Noktayı deneyelim: $1-0-(-1)+7=9\neq0$, nokta düzlemde değil. Sonuç: \textbf{doğru düzleme paraleldir}, düzlemin içinde değildir.

\textit{Not:} 23.10 tahtasında şart ``$aA+b.A+cC=0$'' yazılmıştı; doğrusu $aA+bB+cC=0$.
\end{cevap}

\begin{soru}[Soru 14 \foto{IMG-20261002-WA0026}, \foto{IMG-20261002-WA0027}, \foto{IMG-20261023-WA0061}]
$\dfrac{x-1}{2}=\dfrac{y-1}{-1}=\dfrac{z}{3}$ doğrusunun $2x+3y+z+11=0$ düzlemini kestiği noktayı bulunuz.
\end{soru}
\begin{cevap}[Çözüm (adım adım)]
\textbf{1. Parametrik.} Ortak oran $t$: $x=1+2t,\ y=1-t,\ z=3t$.

\textbf{2. Düzlemde yerine.} $2(1+2t)+3(1-t)+3t+11=0\Rightarrow2+4t+3-3t+3t+11=0\Rightarrow4t+16=0\Rightarrow t=-4$.

\textbf{3. Nokta.} $x=1-8=-7$, $y=1+4=5$, $z=-12$: $\boxed{(-7,5,-12)}$.

\textbf{4. Sağlama.} Düzlem: $-14+15-12+11=0$ \checkmark. Doğru: $\frac{-8}{2}=\frac{4}{-1}=\frac{-12}{3}=-4$ \checkmark.
\end{cevap}

\begin{soru}[Soru 15 \foto{IMG-20261002-WA0026}, \foto{IMG-20261023-WA0033}]
$x+y+z=3$ ve $x-y+z=1$ düzlemlerinin arakesit doğrusunu bulunuz.
\end{soru}
\begin{cevap}[Çözüm (adım adım)]
\textbf{1. Doğrultman.} Arakesit doğrusu iki düzlemde de yatar, iki normale de diktir:
\[ \mathbf v=\mathbf n_1\times\mathbf n_2=\begin{vmatrix}\mathbf i&\mathbf j&\mathbf k\\1&1&1\\1&-1&1\end{vmatrix}
=\mathbf i(1\cdot1-1\cdot(-1))-\mathbf j(1\cdot1-1\cdot1)+\mathbf k(1\cdot(-1)-1\cdot1)=(2,0,-2). \]
Sadeleştirerek $\mathbf v=(1,0,-1)$ alalım.

\textbf{2. Ortak nokta.} $z=0$: $x+y=3$, $x-y=1$ $\Rightarrow$ $x=2$, $y=1$; $A(2,1,0)$.

\textbf{3. Doğru.} $x=2+t,\ y=1,\ z=-t$, yani
\[ \boxed{\frac{x-2}{1}=\frac{z}{-1},\qquad y=1}. \]
\textbf{4. Sağlama.} $(2+t)+1+(-t)=3$ \checkmark; $(2+t)-1+(-t)=1$ \checkmark (her $t$ için).
\end{cevap}

%======================================================================
\subsection{Çalışma soruları}

\begin{soru}[Çalışma 1 (derste çalışma sorusu olarak verildi) \foto{IMG-20261001-WA0014}, \foto{IMG-20261002-WA0027}, \foto{IMG-20261023-WA0048}]
$P(1,0,2)$ noktasının $x+2y-2z+4=0$ düzlemine uzaklığını bulunuz.
\end{soru}
\begin{cevap}
$A=1,\ B=2,\ C=-2,\ D=4$:
\[ d=\frac{|1\cdot1+2\cdot0-2\cdot2+4|}{\sqrt{1+4+4}}=\frac{|1+0-4+4|}{3}=\boxed{\frac13}. \]
\end{cevap}

\begin{soru}[Çalışma 2 (derste çalışma sorusu olarak verildi) \foto{IMG-20261001-WA0014}, \foto{IMG-20261002-WA0027}, \foto{IMG-20261023-WA0048}]
$A(1,0,0)$, $B(0,1,0)$, $C(0,0,1)$ noktalarından geçen düzlemin denklemini bulunuz.
\end{soru}
\begin{cevap}[Çözüm (adım adım)]
$\overrightarrow{AB}=(-1,1,0)$, $\overrightarrow{AC}=(-1,0,1)$, $\overrightarrow{AP}=(x-1,y,z)$.
\[ \begin{vmatrix}x-1&y&z\\-1&1&0\\-1&0&1\end{vmatrix}=(x-1)(1-0)-y(-1-0)+z(0+1)=x-1+y+z=0 . \]
$\boxed{x+y+z=1}$. Kontrol: üç nokta da sağlar \checkmark. Başka yol: eksenleri $1,1,1$'de kesen düzlem, eksen kesişim biçiminden $\frac x1+\frac y1+\frac z1=1$.
\end{cevap}

\begin{soru}[Çalışma 3 (derste çalışma sorusu olarak verildi) \foto{IMG-20261001-WA0014}, \foto{IMG-20261002-WA0027}, \foto{IMG-20261023-WA0048}]
$\dfrac{x-2}{1}=\dfrac{y}{2}=\dfrac{z+1}{-1}$ doğrusu ile $x-y+2z=3$ düzleminin kesim noktasını bulunuz.
\end{soru}
\begin{cevap}[Çözüm (adım adım)]
Parametrik: $x=2+t,\ y=2t,\ z=-1-t$. Düzlemde yerine:
\[ (2+t)-2t+2(-1-t)=3\ \Rightarrow\ 2+t-2t-2-2t=3\ \Rightarrow\ -3t=3\ \Rightarrow\ t=-1 . \]
Nokta: $x=1$, $y=-2$, $z=0$: $\boxed{(1,-2,0)}$. Sağlama: $1-(-2)+0=3$ \checkmark.
\end{cevap}

\begin{soru}[Çalışma 4 (derste çalışma sorusu olarak verildi) \foto{IMG-20261001-WA0014}, \foto{IMG-20261002-WA0027}, \foto{IMG-20261023-WA0048}]
$\dfrac{x}{1}=\dfrac{y}{1}=\dfrac{z}{0}$ ve $\dfrac{x-1}{1}=\dfrac{y}{-1}=\dfrac{z}{0}$ doğruları arasındaki açıyı bulunuz.
\end{soru}
\begin{cevap}[Çözüm (adım adım)]
\textbf{1. Yazımın anlamı.} Paydadaki $0$ bir bölme değildir; doğrultmanın üçüncü bileşeni $0$ demektir, yani $z=0$. Doğrular: $d_1:\ x=t,\ y=t,\ z=0$ ve $d_2:\ x=1+s,\ y=-s,\ z=0$.

\textbf{2. Doğrultmanlar.} $\mathbf v=(1,1,0)$, $\mathbf u=(1,-1,0)$; $\mathbf u\cdot\mathbf v=1-1+0=0$.

\textbf{3. Sonuç.} $\cos\theta=0\Rightarrow\boxed{\theta=90^\circ}$: doğrular diktir. (Üstelik ikisi de $z=0$ düzleminde olup $t=s+1$, $t=-s$ $\Rightarrow$ $s=-\frac12$, $t=\frac12$ için $(\frac12,\frac12,0)$ noktasında kesişirler.)
\end{cevap}

\begin{soru}[Çalışma 5 (derste çalışma sorusu olarak verildi) \foto{IMG-20261009-WA0032}, \foto{IMG-20261023-WA0027}]
$(2,-1,3)$ ile $(1,4,2)$ vektörlerinin dik olup olmadığını belirleyiniz ve aralarındaki açının kosinüsünü bulunuz.
\end{soru}
\begin{cevap}
(Bu soru Konu~1 bilgisiyle çözülür; derste Konu~2 çalışma sorularıyla birlikte verildi.)
$(2,-1,3)\cdot(1,4,2)=2-4+6=4\neq0$ $\Rightarrow$ \textbf{dik değildir}. $\|(2,-1,3)\|=\sqrt{14}$, $\|(1,4,2)\|=\sqrt{21}$:
\[ \cos\theta=\frac{4}{\sqrt{14}\sqrt{21}}=\frac{4}{\sqrt{294}}=\frac{4}{7\sqrt6}=\frac{2\sqrt6}{21}\approx0{,}233\qquad(\theta\approx76{,}5^\circ). \]
\end{cevap}

\begin{soru}[Çalışma 6 (derste çalışma sorusu olarak verildi) \foto{IMG-20261009-WA0032}, \foto{IMG-20261023-WA0027}]
$x^2+y^2+z^2+4x-6y+2z-11=0$ küresinin merkezini ve yarıçapını bulunuz.
\end{soru}
\begin{cevap}
Kareye tamamlama: $(x^2+4x+4)+(y^2-6y+9)+(z^2+2z+1)=11+4+9+1$, yani
\[ (x+2)^2+(y-3)^2+(z+1)^2=25 . \]
Merkez $\boxed{M(-2,3,-1)}$, yarıçap $\boxed{r=5}$. (Küre, elipsoidin $a=b=c$ özel hâlidir.)
\end{cevap}

\begin{soru}[Çalışma 7 (derste çalışma sorusu olarak verildi) \foto{IMG-20261009-WA0032}, \foto{IMG-20261023-WA0027}]
$A(0,1,2)$ ve $B(3,1,-1)$ noktalarından geçen doğrunun parametrik ve kartezyen denklemlerini bulunuz.
\end{soru}
\begin{cevap}[Çözüm (adım adım)]
\textbf{1.} $\overrightarrow{AB}=(3,0,-3)\parallel(1,0,-1)$.

\textbf{2. Parametrik.} $A$ ve $\overrightarrow{AB}$ ile: $x=3t,\ y=1,\ z=2-3t$ ($t\in\mathbb R$). (Sadeleşmiş doğrultmanla: $x=t,\ y=1,\ z=2-t$.)

\textbf{3. Kartezyen.} $y$ bileşeni $0$ olduğundan $y=1$ sabit:
\[ \boxed{\frac{x}{3}=\frac{z-2}{-3},\qquad y=1}\quad(\text{yani } x+z=2,\ y=1). \]
Kontrol: $B$: $\frac33=1$, $\frac{-1-2}{-3}=1$, $y=1$ \checkmark.
\end{cevap}

%======================================================================
\subsection{Silindirler ve kuadratik yüzeyler}

\begin{soru}[Soru 16 \foto{IMG-20261022-WA0050}, \foto{IMG-20261022-WA0052}, \foto{IMG-20261108-WA0096}]
Tahtada örnek olarak verilen $x^2+y^2=a^2$, $z=4x^2$ ve $-\dfrac{x^2}{a^2}+\dfrac{y^2}{b^2}=1$ yüzeylerinin neden silindir olduğunu açıklayınız ve her birinin ana doğrularının yönünü ve dayanak eğrisini belirtiniz.
\end{soru}
\begin{cevap}[Çözüm (adım adım)]
Bir denklemde değişkenlerden biri yoksa, o değişken serbesttir: dayanak eğrisi üzerindeki her noktadan o eksene paralel doğru tamamen yüzeyde kalır.
\begin{itemize}
  \item $x^2+y^2=a^2$: $z$ yok $\Rightarrow$ ana doğrular $Oz$'ye paralel; dayanak eğrisi $xy$-düzleminde yarıçapı $a$ olan çember $\Rightarrow$ \textbf{dairesel silindir}.
  \item $z=4x^2$: $y$ yok $\Rightarrow$ ana doğrular $Oy$'ye paralel; dayanak eğrisi $xz$-düzleminde $z=4x^2$ parabolü $\Rightarrow$ \textbf{parabolik silindir} ($y$ ekseni boyunca uzanan oluk).
  \item $-\frac{x^2}{a^2}+\frac{y^2}{b^2}=1$: $z$ yok $\Rightarrow$ ana doğrular $Oz$'ye paralel; dayanak eğrisi $xy$-düzleminde hiperbol. $\frac{y^2}{b^2}=1+\frac{x^2}{a^2}\ge1$ olduğundan $|y|\ge b$: hiperbolün kolları $\pm y$ yönünde açılır ($y=\pm b$ tepe noktaları). Yüzey $y$ ekseninin iki yanında, $z$ yönünde uzanan iki ayrı yapraktan oluşan \textbf{hiperbolik silindir}dir.
\end{itemize}
\end{cevap}

\begin{soru}[Soru 17 \foto{IMG-20261022-WA0049}, \foto{IMG-20261108-WA0097}]
Tek kanatlı hiperboloid $\frac{x^2}{a^2}+\frac{y^2}{b^2}-\frac{z^2}{c^2}=1$ için ``$z=0$'da elips, $x=0$'da hiperbol'' olduğunu; çift kanatlı hiperboloid $\frac{x^2}{a^2}+\frac{y^2}{b^2}-\frac{z^2}{c^2}=-1$ için de ``$-c<z<c$ aralığında yüzey yok'' olduğunu gösteriniz.
\end{soru}
\begin{cevap}[Çözüm (adım adım)]
\textbf{Tek kanatlı.} $z=0$ koyunca $\frac{x^2}{a^2}+\frac{y^2}{b^2}=1$: yarı eksenleri $a,b$ olan \textbf{elips} (bel elipsi). $x=0$ koyunca $\frac{y^2}{b^2}-\frac{z^2}{c^2}=1$: $yz$-düzleminde \textbf{hiperbol}. Genel $z=k$ izi $\frac{x^2}{a^2}+\frac{y^2}{b^2}=1+\frac{k^2}{c^2}>0$, her $k$ için elips olduğundan yüzey tek parçadır.

\textbf{Çift kanatlı.} $z=k$ koyunca
\[ \frac{x^2}{a^2}+\frac{y^2}{b^2}=\frac{k^2}{c^2}-1 . \]
Sol taraf $\ge0$ olduğundan $\frac{k^2}{c^2}\ge1$, yani $|k|\ge c$ olmalıdır. $-c<k<c$ ise sağ taraf negatiftir ve hiçbir $(x,y)$ denklemi sağlamaz: bu şeritte yüzeyin noktası yoktur. $|k|=c$'de tek nokta $(0,0,\pm c)$ (tepe noktaları), $|k|>c$'de elipsler vardır. Böylece yüzey $z\ge c$ ve $z\le-c$'de iki ayrı kanattan oluşur.
\end{cevap}

\begin{soru}[Soru 18 \foto{IMG-20261015-WA0043}, \foto{IMG-20261023-WA0063}]
Elipsoid $\frac{x^2}{a^2}+\frac{y^2}{b^2}+\frac{z^2}{c^2}=1$'in eksenleri $(\mp a,0,0)$, $(0,\mp b,0)$, $(0,0,\mp c)$ noktalarında kestiğini gösteriniz ve üst/alt yarı elipsoidin denklemlerini yazınız.
\end{soru}
\begin{cevap}
$Ox$ ekseninde $y=z=0$: $\frac{x^2}{a^2}=1\Rightarrow x=\mp a$. Benzer şekilde $Oy$ üzerinde $y=\mp b$, $Oz$ üzerinde $z=\mp c$.
$z$ çekilirse $z^2=c^2\left(1-\frac{x^2}{a^2}-\frac{y^2}{b^2}\right)$, yani
\[ z=c\sqrt{1-\frac{x^2}{a^2}-\frac{y^2}{b^2}}\ \ (\text{üst yarı}),\qquad z=-c\sqrt{1-\frac{x^2}{a^2}-\frac{y^2}{b^2}}\ \ (\text{alt yarı}), \]
tanım kümesi $\frac{x^2}{a^2}+\frac{y^2}{b^2}\le1$ elips bölgesidir.
\end{cevap}

\begin{soru}[Soru 19 \foto{IMG-20261022-WA0052}, \foto{IMG-20261108-WA0096}, \foto{IMG-20261108-WA0095}]
\textbf{(a)} $\frac{z^2}{c^2}=\frac{x^2}{a^2}+\frac{y^2}{b^2}$ konisinin $x=0$ izinin bir doğru çifti olduğunu gösteriniz.
\textbf{(b)} $\frac zc=\frac{x^2}{a^2}-\frac{y^2}{b^2}$ ($c>0$) eyer yüzeyinin $y=0$ ve $x=0$ izlerini bulunuz.
\end{soru}
\begin{cevap}
\textbf{(a)} $x=0$: $\frac{z^2}{c^2}=\frac{y^2}{b^2}\Rightarrow z=\pm\frac cb\,y$; orijinden geçen iki doğru. Koni bu tür doğrulardan (tepeden geçen ana doğrulardan) oluşur.

\textbf{(b)} $y=0$: $z=\frac{c}{a^2}x^2$, yukarı açılan parabol. $x=0$: $z=-\frac{c}{b^2}y^2$, aşağı açılan parabol. Orijin, bir yönde en alçak, diğer yönde en yüksek noktadır: \textbf{eyer noktası}. $z=0$ izi $y=\pm\frac ba x$ doğru çiftidir.
\end{cevap}

%======================================================================
\subsection{Ek alıştırmalar}

\begin{soru}[\ekal Ek alıştırma 1]
$d_1:\ \mathbf r=(1,1,0)+t(1,0,1)$ ve $d_2:\ \mathbf r=(0,2,1)+s(0,1,1)$ doğrularının durumunu belirleyiniz; aykırı iseler aralarındaki uzaklığı bulunuz.
\end{soru}
\begin{cevap}[Çözüm (adım adım)]
\textbf{1. Paralel mi?} $(1,0,1)$ ile $(0,1,1)$ orantılı değil $\Rightarrow$ paralel değil.

\textbf{2. Kesişiyor mu?} $x$: $1+t=0\Rightarrow t=-1$; $y$: $1=2+s\Rightarrow s=-1$; $z$: $t=-1$ fakat $1+s=0$; $-1\neq0$ $\Rightarrow$ kesişmez. Doğrular \textbf{aykırıdır}.

\textbf{3. Ortak dik.} $\mathbf v\times\mathbf u=\begin{vmatrix}\mathbf i&\mathbf j&\mathbf k\\1&0&1\\0&1&1\end{vmatrix}=(0-1,\ -(1-0),\ 1-0)=(-1,-1,1)$.

\textbf{4. Uzaklık.} $\overrightarrow{P_1P_2}=(0,2,1)-(1,1,0)=(-1,1,1)$;
\[ d=\frac{|(-1)(-1)+(1)(-1)+(1)(1)|}{\sqrt{1+1+1}}=\frac{1}{\sqrt3}=\frac{\sqrt3}{3}. \]
\end{cevap}

\begin{soru}[\ekal Ek alıştırma 2]
$\dfrac{x-1}{2}=\dfrac{y+2}{1}=\dfrac{z}{3}$ doğrusunu ve $Q(0,1,1)$ noktasını içeren düzlemin denklemini bulunuz.
\end{soru}
\begin{cevap}[Çözüm (adım adım)]
\textbf{1.} Doğrunun noktası $A(1,-2,0)$, doğrultmanı $\mathbf v=(2,1,3)$; $\overrightarrow{AQ}=(-1,3,1)$. Her iki vektör de düzlemde yatar.

\textbf{2. Normal.}
\[ \mathbf n=\mathbf v\times\overrightarrow{AQ}=\begin{vmatrix}\mathbf i&\mathbf j&\mathbf k\\2&1&3\\-1&3&1\end{vmatrix}=(1-9,\ -(2+3),\ 6+1)=(-8,-5,7). \]
\textbf{3. Düzlem.} $-8(x-1)-5(y+2)+7(z-0)=0\Rightarrow-8x-5y+7z-2=0\Rightarrow\boxed{8x+5y-7z+2=0}$.

\textbf{4. Kontrol.} $A$: $8-10-0+2=0$; $Q$: $0+5-7+2=0$; doğrunun başka bir noktası $(3,-1,3)$: $24-5-21+2=0$ \checkmark.
\end{cevap}

\begin{soru}[\ekal Ek alıştırma 3]
$2x-y+2z=3$ ile $-4x+2y-4z=6$ düzlemlerinin paralel olduğunu gösterip aralarındaki uzaklığı bulunuz.
\end{soru}
\begin{cevap}[Çözüm (adım adım)]
\textbf{1.} $\frac{2}{-4}=\frac{-1}{2}=\frac{2}{-4}=-\frac12$ $\Rightarrow$ paralel. İkinci denklemi $-2$'ye bölersek $2x-y+2z=-3$; birincisiyle aynı değil $\Rightarrow$ çakışık değil.

\textbf{2.} Aynı katsayılarla: $P_1:\ 2x-y+2z-3=0$ ($D_1=-3$), $P_2:\ 2x-y+2z+3=0$ ($D_2=3$).

\textbf{3.} $d=\dfrac{|3-(-3)|}{\sqrt{4+1+4}}=\dfrac63=\boxed{2}$.
\end{cevap}

\begin{soru}[\ekal Ek alıştırma 4]
$2x-y+z=4$ ve $x+y+2z=1$ düzlemleri arasındaki açıyı ve arakesit doğrusunu bulunuz.
\end{soru}
\begin{cevap}[Çözüm (adım adım)]
\textbf{1. Açı.} $\mathbf n_1=(2,-1,1)$, $\mathbf n_2=(1,1,2)$, $\mathbf n_1\cdot\mathbf n_2=2-1+2=3$, $\|\mathbf n_1\|=\|\mathbf n_2\|=\sqrt6$:
$\cos\theta=\frac{3}{6}=\frac12\Rightarrow\theta=60^\circ$.

\textbf{2. Doğrultman.}
$\mathbf n_1\times\mathbf n_2=\begin{vmatrix}\mathbf i&\mathbf j&\mathbf k\\2&-1&1\\1&1&2\end{vmatrix}=(-2-1,\ -(4-1),\ 2+1)=(-3,-3,3)\parallel(1,1,-1)$.

\textbf{3. Nokta.} $z=0$: $2x-y=4$, $x+y=1$ $\Rightarrow$ $3x=5$, $x=\frac53$, $y=-\frac23$.

\textbf{4. Doğru.} $x=\frac53+t,\ y=-\frac23+t,\ z=-t$. Kontrol: $2(\frac53+t)-(-\frac23+t)-t=\frac{12}{3}=4$ \checkmark; $(\frac53+t)+(-\frac23+t)-2t=1$ \checkmark.
\end{cevap}

\begin{soru}[\ekal Ek alıştırma 5]
$P(1,2,3)$ noktasının $x+y+z=0$ düzlemi üzerindeki dik izdüşümünü (dikme ayağını) ve noktanın düzleme uzaklığını bulunuz.
\end{soru}
\begin{cevap}[Çözüm (adım adım)]
\textbf{1.} $P$'den geçen ve normale $\mathbf n=(1,1,1)$ paralel doğru: $x=1+t,\ y=2+t,\ z=3+t$.

\textbf{2.} Düzlemde yerine: $(1+t)+(2+t)+(3+t)=0\Rightarrow6+3t=0\Rightarrow t=-2$.

\textbf{3.} Dikme ayağı $K(-1,0,1)$.

\textbf{4.} Uzaklık: $|PK|=\sqrt{4+4+4}=2\sqrt3$. Formülle: $\frac{|1+2+3|}{\sqrt3}=\frac{6}{\sqrt3}=2\sqrt3$ \checkmark.
\end{cevap}

\begin{soru}[\ekal Ek alıştırma 6]
$x^2+4y^2+z^2-2x+8y+1=0$ yüzeyini tanıyınız.
\end{soru}
\begin{cevap}[Çözüm (adım adım)]
\textbf{1. Kareye tamamlama.} $x^2-2x=(x-1)^2-1$, $4y^2+8y=4(y+1)^2-4$.

\textbf{2.} $(x-1)^2-1+4(y+1)^2-4+z^2+1=0\Rightarrow(x-1)^2+4(y+1)^2+z^2=4$.

\textbf{3.} $4$'e bölersek
\[ \frac{(x-1)^2}{4}+\frac{(y+1)^2}{1}+\frac{z^2}{4}=1 . \]
Üç kare, hepsi pozitif $\Rightarrow$ \textbf{elipsoid}; merkezi $(1,-1,0)$, yarı eksenleri $x$ yönünde $2$, $y$ yönünde $1$, $z$ yönünde $2$. ($a=c$ olduğundan $Oy$'ye paralel eksen etrafında dönel bir elipsoiddir.)
\end{cevap}

\begin{soru}[\ekal Ek alıştırma 7]
$y=z^2-x^2$ yüzeyini izleriyle tanıyınız.
\end{soru}
\begin{cevap}[Çözüm (adım adım)]
İki değişken kare, zıt işaretli; $y$ birinci dereceden $\Rightarrow$ \textbf{hiperbolik paraboloid} (eyer), ekseni $Oy$.
\begin{itemize}
  \item $y=k$: $z^2-x^2=k$; $k\neq0$ için hiperbol ($k>0$'da $z$ yönünde, $k<0$'da $x$ yönünde açılır), $k=0$ için $z=\pm x$ doğru çifti.
  \item $x=k$: $y=z^2-k^2$, $+y$ yönünde açılan parabol.
  \item $z=k$: $y=k^2-x^2$, $-y$ yönünde açılan parabol.
\end{itemize}
Orijin eyer noktasıdır: $z$ yönünde yürüyünce $y$ artar, $x$ yönünde yürüyünce azalır.
\end{cevap}

\clearpage\section{Konu 3: Vektör Fonksiyonları ve Uzay Eğrileri}

Bu bölümde 09.10 ve 15.10 tarihli defter sayfalarında ve 23.10 tarihli tahta fotoğraflarında çözülen bütün sorular, derste
işlenme sırasıyla verilmiştir. Aynı soru birden fazla fotoğrafta geçiyorsa bütün kaynaklar başlıkta yazılıdır. Sonda, derste
çalışma sorusu olarak verilen problemlerin ve kaynak kitaplardaki tipik alıştırmalara benzeyen ek alıştırmaların tam çözümleri vardır.

\begin{bilgi}[Hatırlatma]
Vektör fonksiyonlarda limit, süreklilik, türev ve integral \emph{bileşen bileşen} yapılır:
$\mathbf r(t)=x(t)\mathbf i+y(t)\mathbf j+z(t)\mathbf k$ ise $\mathbf r'(t)=x'(t)\mathbf i+y'(t)\mathbf j+z'(t)\mathbf k$,
$\int\mathbf r\,dt=\big(\int x\,dt\big)\mathbf i+\big(\int y\,dt\big)\mathbf j+\big(\int z\,dt\big)\mathbf k+\mathbf C$.
\end{bilgi}

%======================================================================
\subsection{Vektör fonksiyonları: tanım kümesi, işlemler, limit, süreklilik}

\begin{soru}[Soru 1 \foto{IMG-20261009-WA0016}, \foto{IMG-20261023-WA0060}]
$\mathbf F(t)=(t+1)\,\mathbf i+\ln(1-t)\,\mathbf j+\sqrt{4-t^2}\,\mathbf k$ fonksiyonunun tanım kümesini bulunuz.
\end{soru}
\begin{cevap}[Çözüm (adım adım)]
\textbf{Fikir:} $\mathbf F(t)$'nin tanımlı olması için üç bileşenin de tanımlı olması gerekir: $T_{\mathbf F}=T_f\cap T_g\cap T_h$.

\textbf{Adım 1.} $f(t)=t+1$ polinomdur, her reel sayıda tanımlı: $T_f=(-\infty,+\infty)$.

\textbf{Adım 2.} $g(t)=\ln(1-t)$: logaritmanın argümanı pozitif olmalı. $1-t>0\iff t<1$, \ $T_g=(-\infty,1)$.

\textbf{Adım 3.} $h(t)=\sqrt{4-t^2}$: kök içi negatif olamaz. $4-t^2\ge0\iff t^2\le4\iff-2\le t\le2$, \ $T_h=[-2,2]$.

\textbf{Adım 4.} Kesişim (sayı doğrusunda üçünün ortak kısmı):
\[ T_{\mathbf F}=(-\infty,+\infty)\cap(-\infty,1)\cap[-2,2]=\boxed{[-2,1)} . \]
$-2$ dahildir ($\sqrt0=0$), $1$ dahil değildir ($\ln0$ tanımsız).
\end{cevap}

\begin{soru}[Soru 2 \foto{IMG-20261009-WA0017}, \foto{IMG-20261023-WA0060}, \foto{IMG-20261023-WA0035}]
$\mathbf F(t)=t\,\mathbf i+t^2\mathbf j+t^3\mathbf k$, $\mathbf G(t)=(1-t)\,\mathbf i+(1+t^2)\,\mathbf j+t\,\mathbf k$, $f(t)=t+1$, $g(t)=t-2$ için
$(\mathbf F+\mathbf G)(t)$, $(\mathbf F\cdot\mathbf G)(t)$, $(\mathbf F\times\mathbf G)(t)$, $(f\mathbf G)(t)$ ve $(\mathbf F\circ g)(t)$'yi bulunuz.
\end{soru}
\begin{cevap}[Çözüm (adım adım)]
\textbf{Toplam:} karşılıklı bileşenler toplanır.
\[ (\mathbf F+\mathbf G)(t)=(t+1-t)\,\mathbf i+(t^2+1+t^2)\,\mathbf j+(t^3+t)\,\mathbf k=\mathbf i+(2t^2+1)\,\mathbf j+(t^3+t)\,\mathbf k . \]
\textbf{Skaler çarpım:} $x$'ler, $y$'ler, $z$'ler çarpılıp toplanır; sonuç reel fonksiyondur.
\[ (\mathbf F\cdot\mathbf G)(t)=t(1-t)+t^2(1+t^2)+t^3\cdot t=t-t^2+t^2+t^4+t^4=2t^4+t . \]
\textbf{Vektörel çarpım:} determinant birinci satıra göre açılır.
\begin{align*}
(\mathbf F\times\mathbf G)(t)&=\begin{vmatrix}\mathbf i&\mathbf j&\mathbf k\\t&t^2&t^3\\1-t&1+t^2&t\end{vmatrix}\\
&=\mathbf i\big(t^3-t^3(1+t^2)\big)-\mathbf j\big(t^2-t^3(1-t)\big)+\mathbf k\big(t(1+t^2)-t^2(1-t)\big)\\
&=-t^5\,\mathbf i-(t^4-t^3+t^2)\,\mathbf j+(2t^3-t^2+t)\,\mathbf k .
\end{align*}
\textbf{Skalerle çarpım:}
\[ (f\mathbf G)(t)=(t+1)\big[(1-t)\mathbf i+(1+t^2)\mathbf j+t\,\mathbf k\big]=(1-t^2)\,\mathbf i+(t^3+t^2+t+1)\,\mathbf j+(t^2+t)\,\mathbf k . \]
\textbf{Bileşke:} $\mathbf F$'de $t$ yerine $g(t)=t-2$ yazılır.
\[ (\mathbf F\circ g)(t)=\mathbf F(t-2)=(t-2)\,\mathbf i+(t-2)^2\,\mathbf j+(t-2)^3\,\mathbf k . \]
\textbf{Kontrol:} $t=1$'de $\mathbf F=(1,1,1)$, $\mathbf G=(0,2,1)$, $\mathbf F\times\mathbf G=(1\cdot1-1\cdot2,\ 1\cdot0-1\cdot1,\ 1\cdot2-1\cdot0)=(-1,-1,2)$;
formülümüz de $(-1,-(1-1+1),2-1+1)=(-1,-1,2)$ verir. \checkmark
\end{cevap}

\begin{soru}[Soru 3 \foto{IMG-20261009-WA0017}, \foto{IMG-20261023-WA0035}]
$\mathbf F(t)=(t+t^2)\,\mathbf i+2t\,\mathbf j-4t\,\mathbf k$ ve $\mathbf G(t)=(2+4t)\,\mathbf i+(1-3t)\,\mathbf j+(1+t^2)\,\mathbf k$ ortogonal midir?
\end{soru}
\begin{cevap}[Çözüm (adım adım)]
Ortogonallik için \emph{her} $t$'de $\mathbf F(t)\cdot\mathbf G(t)=0$ olmalıdır.
\begin{align*}
\mathbf F\cdot\mathbf G&=(t+t^2)(2+4t)+2t(1-3t)+(-4t)(1+t^2)\\
&=(2t+4t^2+2t^2+4t^3)+(2t-6t^2)+(-4t-4t^3)\\
&=(2t+2t-4t)+(6t^2-6t^2)+(4t^3-4t^3)=0 .
\end{align*}
Özdeş olarak sıfırdır: fonksiyonlar \textbf{ortogonaldir}.
\end{cevap}

\begin{soru}[Soru 4 \foto{IMG-20261009-WA0018}, \foto{IMG-20261023-WA0035}, \foto{IMG-20261023-WA0046}]
$\mathbf F(t)=2\cos t\,\mathbf i+2\sin t\,\mathbf j+3t\,\mathbf k$ fonksiyonunun normunu bulunuz.
\end{soru}
\begin{cevap}
\[ \|\mathbf F(t)\|=\sqrt{(2\cos t)^2+(2\sin t)^2+(3t)^2}=\sqrt{4(\cos^2t+\sin^2t)+9t^2}=\sqrt{4+9t^2} . \]
\end{cevap}

\begin{soru}[Soru 5 \foto{IMG-20261009-WA0018}, \foto{IMG-20261023-WA0036}]
$\mathbf F(t)=(1-t^2)\,\mathbf i+\dfrac{\sin8t}{t}\,\mathbf j+(e^t+1)\,\mathbf k$ için $\displaystyle\lim_{t\to0}\mathbf F(t)$'yi bulunuz.
\end{soru}
\begin{cevap}[Çözüm (adım adım)]
Limit bileşen bileşen alınır.
\begin{itemize}
\item $\mathbf i$: $\lim_{t\to0}(1-t^2)=1$.
\item $\mathbf j$: $\dfrac{\sin8t}{t}=8\cdot\dfrac{\sin8t}{8t}$ ve $u=8t\to0$ iken $\dfrac{\sin u}{u}\to1$; limit $8$.
\item $\mathbf k$: $\lim_{t\to0}(e^t+1)=2$.
\end{itemize}
\[ \lim_{t\to0}\mathbf F(t)=\mathbf i+8\,\mathbf j+2\,\mathbf k . \]
\textit{Not:} Tahtada limit ``$t\to t_0$'' diye yazılmıştır; sonuç yalnızca $t\to0$ için geçerlidir.
\end{cevap}

\begin{soru}[Soru 6 \foto{IMG-20261009-WA0019}, \foto{IMG-20261023-WA0046}, \foto{IMG-20261023-WA0037}]
$\mathbf F(t)=\sin t\,\mathbf i+\cos t\,\mathbf j+(1+t)\,\mathbf k$, $\mathbf G(t)=\cos2t\,\mathbf i+\sin3t\,\mathbf j+(e^t+1)\,\mathbf k$ için
$\displaystyle\lim_{t\to0}\mathbf F(t)\cdot\mathbf G(t)$ ve $\displaystyle\lim_{t\to0}\mathbf F(t)\times\mathbf G(t)$'yi bulunuz.
\end{soru}
\begin{cevap}[Çözüm (adım adım)]
\textbf{Adım 1.} $\lim_{t\to0}\mathbf F(t)=(\sin0,\cos0,1+0)=(0,1,1)$; \ $\lim_{t\to0}\mathbf G(t)=(\cos0,\sin0,e^0+1)=(1,0,2)$.

\textbf{Adım 2.} Limit kuralları: çarpımın limiti, limitlerin çarpımıdır.
\[ \lim_{t\to0}\mathbf F\cdot\mathbf G=(0,1,1)\cdot(1,0,2)=0+0+2=2 . \]
\[ \lim_{t\to0}\mathbf F\times\mathbf G=\begin{vmatrix}\mathbf i&\mathbf j&\mathbf k\\0&1&1\\1&0&2\end{vmatrix}=\mathbf i(2-0)-\mathbf j(0-1)+\mathbf k(0-1)=2\,\mathbf i+\mathbf j-\mathbf k . \]
\end{cevap}

\begin{soru}[Soru 7 \foto{IMG-20261009-WA0019}, \foto{IMG-20261009-WA0021}, \foto{IMG-20261023-WA0037}]
$\mathbf F(t)=(|t|+1)\,\mathbf i+t\,[\![t]\!]\,\mathbf j+(2+t^2)\,\mathbf k$ fonksiyonu $t=0$ ve $t=1$ noktalarında sürekli midir?
($[\![t]\!]$: tam değer fonksiyonu.)
\end{soru}
\begin{cevap}[Çözüm (adım adım)]
$\mathbf F$'nin sürekliliği bileşenlerin sürekliliğine denktir. $|t|+1$ ve $2+t^2$ her yerde süreklidir; yalnızca
$g(t)=t\,[\![t]\!]$ incelenir.

\textbf{$t=0$:} $g(0)=0$. Soldan ($-1<t<0$): $[\![t]\!]=-1$, $g(t)=-t\to0$. Sağdan ($0<t<1$): $[\![t]\!]=0$, $g(t)=0$.
Sağ limit $=$ sol limit $=g(0)=0$ $\Rightarrow$ $\mathbf F$, $t=0$'da \textbf{süreklidir}.

\textbf{$t=1$:} Soldan ($0<t<1$): $g(t)=t\cdot0\to0$. Sağdan ($1<t<2$): $g(t)=t\cdot1\to1$. Sol ve sağ limitler farklı, limit yok
$\Rightarrow$ $\mathbf F$, $t=1$'de \textbf{sürekli değildir}.
\end{cevap}

%======================================================================
\subsection{Uzay eğrileri}

\begin{soru}[Soru 8 \foto{IMG-20261009-WA0021}, \foto{IMG-20261023-WA0051}, \foto{IMG-20261023-WA0054}]
$\mathbf r(t)=(2+3t)\,\mathbf i+(4+5t)\,\mathbf j+(6+4t)\,\mathbf k$ eğrisinin cinsi nedir?
\end{soru}
\begin{cevap}
$x=2+3t$, $y=4+5t$, $z=6+4t$; her birinden $t$ çekilip eşitlenirse
\[ \frac{x-2}{3}=\frac{y-4}{5}=\frac{z-6}{4}\ (=t) . \]
Bu, $(2,4,6)$ noktasından geçen ve $\mathbf v=(3,5,4)$ vektörüne paralel \textbf{doğrunun} denklemidir.
\end{cevap}

\begin{soru}[Soru 9 \foto{IMG-20261009-WA0022}, \foto{IMG-20261023-WA0051}, \foto{IMG-20261023-WA0054}]
$\mathbf r(t)=4\cos t\,\mathbf i+4\sin t\,\mathbf j$, $0\le t\le2\pi$ eğrisini tanıyınız.
\end{soru}
\begin{cevap}
$x=4\cos t$, $y=4\sin t$ $\Rightarrow$ $x^2+y^2=16(\cos^2t+\sin^2t)=16$. $xy$ düzleminde, merkezi orijin, yarıçapı $4$ olan \textbf{çember}.
\end{cevap}

\begin{soru}[Soru 10 \foto{IMG-20261009-WA0022}, \foto{IMG-20261023-WA0054}, \foto{IMG-20261023-WA0059}]
$\mathbf r(t)=\cos t\,\mathbf i+\sin t\,\mathbf j+3\,\mathbf k$, $0\le t\le2\pi$ eğrisini tanıyınız.
\end{soru}
\begin{cevap}
$x=\cos t$, $y=\sin t$, $z=3$. $z$ sabit: eğri $z=3$ düzleminde. $x^2+y^2=\cos^2t+\sin^2t=1$ (tahtada ``$y^2+y^2$'' yazılmış,
doğrusu $x^2+y^2$). Sonuç: $z=3$ düzleminde, merkezi $(0,0,3)$, yarıçapı $1$ olan \textbf{çember}.
\end{cevap}

\begin{soru}[Soru 11 \foto{IMG-20261009-WA0025}, \foto{IMG-20261023-WA0054}]
$\mathbf r(t)=2\cos t\,\mathbf i+2\sin t\,\mathbf j+4t\,\mathbf k$ eğrisini inceleyiniz.
\end{soru}
\begin{cevap}[Çözüm (adım adım)]
\textbf{Adım 1.} $x=2\cos t$, $y=2\sin t$, $z=4t$.

\textbf{Adım 2.} $x^2+y^2=4(\cos^2t+\sin^2t)=4$: noktalar $x^2+y^2=4$ silindiri üzerindedir. (Defter ve tahtada
``$4(\cos t+\sin t)$'' yazılmış; kareler unutulmuştur.)

\textbf{Adım 3.} $z=4t$ düzgün artar: $t=1\Rightarrow z=4$, $t=2\Rightarrow z=8$. Bir turda ($t:0\to2\pi$) yükselme $4\cdot2\pi=8\pi$.

\textbf{Sonuç:} Silindire sarılan, adımı $8\pi$ olan \textbf{dairesel helis}.
\end{cevap}

\begin{soru}[Soru 12 \foto{IMG-20261009-WA0025}, \foto{IMG-20261009-WA0020}, \foto{IMG-20261023-WA0039}, \foto{IMG-20261023-WA0059}]
$\mathbf r(t)=(t^2-4)\,\mathbf i+(t^3-4t)\,\mathbf j$, $-2\le t\le2$ eğrisi kapalı mıdır? Basit midir? Orijinde köşe var mıdır?
\end{soru}
\begin{cevap}[Çözüm (adım adım)]
\textbf{Adım 1 (kapalılık).} $\mathbf r(-2)=(0,0)$ ve $\mathbf r(2)=(0,0)$; uçlar çakışır, eğri \textbf{kapalıdır}.

\textbf{Adım 2 (katlı nokta araması).} $\mathbf r(t_1)=\mathbf r(t_2)$ olsun. $x$'lerden $t_1^2=t_2^2$, yani $t_1=t_2$ (aynı nokta) ya da
$t_2=-t_1$. İkinci durumda $y$'ler:
\[ t_1^3-4t_1=-t_1^3+4t_1\ \Rightarrow\ 2t_1^3-8t_1=0\ \Rightarrow\ t_1(t_1^2-4)=0\ \Rightarrow\ t_1\in\{0,\,2,\,-2\}. \]
$t_1=0$ ise $t_2=0$ (aynı parametre); $t_1=\pm2$ ise $t_2=\mp2$ (uç noktalar). Uçlar dışında kesişme yok:
eğri \textbf{basit kapalı eğridir}.

\textbf{Adım 3 (köşe).} $\mathbf r'(t)=2t\,\mathbf i+(3t^2-4)\,\mathbf j$; $\mathbf r'(2)=4\mathbf i+8\mathbf j$, $\mathbf r'(-2)=-4\mathbf i+8\mathbf j$.
Bu vektörler paralel değil (biri diğerinin skaler katı değil): eğri orijinde \textbf{köşe yapar}, düzgün kapanmaz.
\end{cevap}

\begin{soru}[Soru 13 \foto{IMG-20261009-WA0020}, \foto{IMG-20261023-WA0041}, \foto{IMG-20261023-WA0059}]
$\mathbf r(t)=3\cos t\,\mathbf i+3\sin t\,\mathbf j$, $0\le t\le2\pi$ eğrisinin grafiği ve yönü nedir?
\end{soru}
\begin{cevap}
$x^2+y^2=9$: yarıçapı $3$, merkezi orijin olan çember. $t=0:(3,0)$, $t=\pi/2:(0,3)$, $t=\pi:(-3,0)$, $t=3\pi/2:(0,-3)$,
$t=2\pi:(3,0)$. Nokta çemberi \textbf{saat yönünün tersine} dolaşır; bu, pozitif yön kabul edilir.
\end{cevap}

%======================================================================
\subsection{Türev ve türev kuralları}

\begin{soru}[Soru 14 \foto{IMG-20261009-WA0023}, \foto{IMG-20261023-WA0062}]
$\mathbf F(t)=t^2\,\mathbf i+e^{-t}\,\mathbf j+\sin t\,\mathbf k$ için birinci, ikinci ve üçüncü türevleri bulunuz.
\end{soru}
\begin{cevap}
\begin{align*}
\frac{d\mathbf F}{dt}&=2t\,\mathbf i-e^{-t}\,\mathbf j+\cos t\,\mathbf k, &
\frac{d^2\mathbf F}{dt^2}&=2\,\mathbf i+e^{-t}\,\mathbf j-\sin t\,\mathbf k, &
\frac{d^3\mathbf F}{dt^3}&=-e^{-t}\,\mathbf j-\cos t\,\mathbf k .
\end{align*}
(Paydalar $dt^2$, $dt^3$ olmalıdır; tahtada ``$dt$'' yazılmıştır.)
\end{cevap}

\begin{soru}[Soru 15 \foto{IMG-20261009-WA0024}, \foto{IMG-20261023-WA0062}, \foto{IMG-20261023-WA0064}, \foto{IMG-20261023-WA0068}]
$\mathbf r(t)=(1-\cos t)\,\mathbf i+\sin t\,\mathbf j+t\,\mathbf k$ için $\mathbf r'$, $\mathbf r''$ ve $t=\pi/2$'deki değerlerini bulunuz.
\end{soru}
\begin{cevap}
$\mathbf r'(t)=\sin t\,\mathbf i+\cos t\,\mathbf j+\mathbf k$, \ $\mathbf r'(\pi/2)=\mathbf i+0\,\mathbf j+\mathbf k=\mathbf i+\mathbf k$.

$\mathbf r''(t)=\cos t\,\mathbf i-\sin t\,\mathbf j$, \ $\mathbf r''(\pi/2)=-\mathbf j$.

(Tahtada $\mathbf r'(\pi/2)$'de $\mathbf k$ katsayısı ``$0$'' gibi okunuyor; $\mathbf k$ bileşeni sabit $1$'dir.)
\end{cevap}

\begin{soru}[Soru 16 \foto{IMG-20261009-WA0024}, \foto{IMG-20261023-WA0064}, \foto{IMG-20261023-WA0068}]
$\mathbf r(t)=a\cos t\,\mathbf i+a\sin t\,\mathbf j+bt\,\mathbf k$ için $\mathbf r'$, $\mathbf r''$ ve $\mathbf r'\times\mathbf r''$'yü bulunuz.
\end{soru}
\begin{cevap}[Çözüm (adım adım)]
$\mathbf r'(t)=-a\sin t\,\mathbf i+a\cos t\,\mathbf j+b\,\mathbf k$, \quad $\mathbf r''(t)=-a\cos t\,\mathbf i-a\sin t\,\mathbf j$.
\begin{align*}
\mathbf r'\times\mathbf r''&=\begin{vmatrix}\mathbf i&\mathbf j&\mathbf k\\-a\sin t&a\cos t&b\\-a\cos t&-a\sin t&0\end{vmatrix}\\
&=\mathbf i\big(a\cos t\cdot0-b(-a\sin t)\big)-\mathbf j\big(-a\sin t\cdot0-b(-a\cos t)\big)+\mathbf k\big(a^2\sin^2t+a^2\cos^2t\big)\\
&=ab\sin t\,\mathbf i-ab\cos t\,\mathbf j+a^2\,\mathbf k .
\end{align*}
\end{cevap}

\begin{soru}[Soru 17 \foto{IMG-20261009-WA0024}, \foto{IMG-20261023-WA0064}, \foto{IMG-20261023-WA0068}]
$\mathbf u(t)=t\,\mathbf i-t\,\mathbf j+\frac{t^2}{2}\,\mathbf k$, $\mathbf v(t)=\frac{t^4}{4}\,\mathbf i+\frac{t^3}{3}\,\mathbf j+\frac{t^2}{2}\,\mathbf k$ için
$\dfrac{d}{dt}(\mathbf u\times\mathbf v)$'yi bulunuz.
\end{soru}
\begin{cevap}[Çözüm (iki yol)]
\textbf{1. yol.}
\[ \mathbf u\times\mathbf v=\begin{vmatrix}\mathbf i&\mathbf j&\mathbf k\\t&-t&\frac{t^2}{2}\\\frac{t^4}{4}&\frac{t^3}{3}&\frac{t^2}{2}\end{vmatrix}
=\Big(-\frac{t^3}{2}-\frac{t^5}{6}\Big)\mathbf i+\Big(\frac{t^6}{8}-\frac{t^3}{2}\Big)\mathbf j+\Big(\frac{t^4}{3}+\frac{t^5}{4}\Big)\mathbf k . \]
\[ \frac{d}{dt}(\mathbf u\times\mathbf v)=\Big(-\frac{3t^2}{2}-\frac{5t^4}{6}\Big)\mathbf i+\Big(\frac{3t^5}{4}-\frac{3t^2}{2}\Big)\mathbf j+\Big(\frac{4t^3}{3}+\frac{5t^4}{4}\Big)\mathbf k . \]
\textbf{2. yol.} $(\mathbf u\times\mathbf v)'=\mathbf u'\times\mathbf v+\mathbf u\times\mathbf v'$, \ $\mathbf u'=(1,-1,t)$, $\mathbf v'=(t^3,t^2,t)$:
\begin{align*}
\mathbf u'\times\mathbf v&=\Big(-\frac{t^2}{2}-\frac{t^4}{3},\ \frac{t^5}{4}-\frac{t^2}{2},\ \frac{t^3}{3}+\frac{t^4}{4}\Big),\\
\mathbf u\times\mathbf v'&=\Big(-t^2-\frac{t^4}{2},\ \frac{t^5}{2}-t^2,\ t^3+t^4\Big).
\end{align*}
Toplam 1. yoldaki sonucun aynısıdır. \checkmark
\end{cevap}

\begin{dikkat}[Defterdeki hata]
09.10 defterinde (\foto{IMG-20261009-WA0024}) $\mathbf j$ bileşeninin türevi $\frac{3t^5}{2}$ ile yazılmıştır;
$\frac{d}{dt}\frac{t^6}{8}=\frac{3t^5}{4}$'tür. 23.10 tahtasında doğru ($\frac68t^5$) yazılmıştır.
\end{dikkat}

%======================================================================
\subsection{Frenet üçlüsü, hız--ivme ve eğrilik}

\begin{soru}[Soru 18 \foto{IMG-20261015-WA0034}, \foto{IMG-20261023-WA0047}, \foto{IMG-20261023-WA0067}]
$\mathbf r(t)=a\cos t\,\mathbf i+a\sin t\,\mathbf j$ ($a>0$) çemberinin $\mathbf T$, $\mathbf N$, $\mathbf B$ vektörlerini bulunuz.
\end{soru}
\begin{cevap}[Çözüm (adım adım)]
\textbf{Adım 1.} $\mathbf r'(t)=-a\sin t\,\mathbf i+a\cos t\,\mathbf j$, \ $\|\mathbf r'(t)\|=\sqrt{a^2(\sin^2t+\cos^2t)}=a$.

\textbf{Adım 2.} $\mathbf T=\dfrac{\mathbf r'}{\|\mathbf r'\|}=-\sin t\,\mathbf i+\cos t\,\mathbf j$ (defterde işaret unutulmuş: ``$\sin t\,\mathbf i+\cos t\,\mathbf j$'').

\textbf{Adım 3.} $\mathbf T'(t)=-\cos t\,\mathbf i-\sin t\,\mathbf j$, \ $\|\mathbf T'\|=1$ $\Rightarrow$ $\mathbf N=-\cos t\,\mathbf i-\sin t\,\mathbf j=-\frac1a\mathbf r(t)$ (merkeze doğru).

\textbf{Adım 4.}
\[ \mathbf B=\mathbf T\times\mathbf N=\begin{vmatrix}\mathbf i&\mathbf j&\mathbf k\\-\sin t&\cos t&0\\-\cos t&-\sin t&0\end{vmatrix}=(\sin^2t+\cos^2t)\,\mathbf k=\mathbf k . \]
\end{cevap}

\begin{soru}[Soru 19 \foto{IMG-20261015-WA0036}, \foto{IMG-20261023-WA0049}, \foto{IMG-20261023-WA0067}]
$a>0$, $\mathbf r(t)=a\cos t\,\mathbf i+a\sin t\,\mathbf j$ olan hareketlinin hız vektörü konum vektörüne dik midir?
\end{soru}
\begin{cevap}
$\mathbf v=\mathbf r'=-a\sin t\,\mathbf i+a\cos t\,\mathbf j$; \ $\mathbf v\cdot\mathbf r=-a^2\sin t\cos t+a^2\sin t\cos t=0$.
Her $t$ için $\mathbf v\perp\mathbf r$ (çemberde hız yarıçapa diktir).
\end{cevap}

\begin{soru}[Soru 20 \foto{IMG-20261015-WA0036}, \foto{IMG-20261023-WA0049}, \foto{IMG-20261023-WA0067}]
$\mathbf r(t)=5\cos t\,\mathbf i+5\sin t\,\mathbf j+6t^2\,\mathbf k$ olan hareketlinin hızı ve ivmesi nedir? Hız ve ivme ne zaman diktir?
\end{soru}
\begin{cevap}
$\mathbf v(t)=-5\sin t\,\mathbf i+5\cos t\,\mathbf j+12t\,\mathbf k$, \ $\mathbf a(t)=-5\cos t\,\mathbf i-5\sin t\,\mathbf j+12\,\mathbf k$.
\[ \mathbf a\cdot\mathbf v=25\sin t\cos t-25\sin t\cos t+144t=144t=0\ \Rightarrow\ t=0 . \]
\end{cevap}

\begin{soru}[Soru 21 \foto{IMG-20261015-WA0036}, \foto{IMG-20261023-WA0049}, \foto{IMG-20261023-WA0038}]
Parametrik denklemleri $x=e^{-t}$, $y=2\cos3t$, $z=2\sin3t$ olan eğri boyunca bir cisim hareket ediyor ($t$: zaman).
(a) Hız ve ivme? (b) $t=0$ iken büyüklükleri? (c) $t=0$ iken birim teğet vektör?
\end{soru}
\begin{cevap}[Çözüm (adım adım)]
$\mathbf r(t)=e^{-t}\,\mathbf i+2\cos3t\,\mathbf j+2\sin3t\,\mathbf k$.

\textbf{(a)} Zincir kuralıyla ($\frac{d}{dt}\cos3t=-3\sin3t$):
$\mathbf v=-e^{-t}\,\mathbf i-6\sin3t\,\mathbf j+6\cos3t\,\mathbf k$, \ $\mathbf a=e^{-t}\,\mathbf i-18\cos3t\,\mathbf j-18\sin3t\,\mathbf k$.

\textbf{(b)} $\mathbf v(0)=-\mathbf i+6\mathbf k$, $\|\mathbf v(0)\|=\sqrt{37}$; \ $\mathbf a(0)=\mathbf i-18\mathbf j$, $\|\mathbf a(0)\|=\sqrt{325}=5\sqrt{13}$.

\textbf{(c)} $\mathbf T(0)=\dfrac{\mathbf v(0)}{\|\mathbf v(0)\|}=-\dfrac{1}{\sqrt{37}}\,\mathbf i+\dfrac{6}{\sqrt{37}}\,\mathbf k$.
\end{cevap}

\begin{soru}[Soru 22 \foto{IMG-20261015-WA0035}, \foto{IMG-20261023-WA0052}]
Bir nokta $\mathbf r=\mathbf r(t)$ eğrisi boyunca hareket ediyor; $s$ alınan yol, $\upsilon=\|d\mathbf r/dt\|=ds/dt$.
$\mathbf a=\dfrac{d\upsilon}{dt}\mathbf T+\dfrac{\upsilon^2}{\rho}\mathbf N$ olduğunu gösteriniz.
\end{soru}
\begin{cevap}[Çözüm (adım adım)]
\textbf{Adım 1.} $\mathbf v=\dfrac{d\mathbf r}{dt}=\dfrac{d\mathbf r}{ds}\dfrac{ds}{dt}=\upsilon\,\mathbf T$.

\textbf{Adım 2.} Çarpım kuralı: $\mathbf a=\dfrac{d\mathbf v}{dt}=\dfrac{d\upsilon}{dt}\mathbf T+\upsilon\dfrac{d\mathbf T}{dt}$.

\textbf{Adım 3.} Zincir kuralı: $\dfrac{d\mathbf T}{dt}=\dfrac{d\mathbf T}{ds}\dfrac{ds}{dt}=\upsilon\dfrac{d\mathbf T}{ds}$; eğrilik tanımından $\dfrac{d\mathbf T}{ds}=\kappa\mathbf N=\dfrac1\rho\mathbf N$.

\textbf{Adım 4.} Yerine yazılırsa $\mathbf a=\dfrac{d\upsilon}{dt}\mathbf T+\upsilon^2\dfrac{d\mathbf T}{ds}=\dfrac{d\upsilon}{dt}\mathbf T+\dfrac{\upsilon^2}{\rho}\mathbf N$. \hfill$\blacksquare$
\end{cevap}

\begin{soru}[Soru 23 \foto{IMG-20261015-WA0038}, \foto{IMG-20261023-WA0052}]
$\mathbf r(t)=a\cos t\,\mathbf i+a\sin t\,\mathbf j$ çemberinin eğriliğini ve eğrilik yarıçapını bulunuz; $t$ zaman ise ivmeyi
$\mathbf T$, $\mathbf N$ bileşenlerine ayırınız.
\end{soru}
\begin{cevap}[Çözüm (adım adım)]
Soru 18'den $\mathbf T'(t)=-\cos t\,\mathbf i-\sin t\,\mathbf j$, $\frac{ds}{dt}=a$.
\[ \frac{d\mathbf T}{ds}=\frac{d\mathbf T/dt}{ds/dt}=\frac1a(-\cos t\,\mathbf i-\sin t\,\mathbf j)\ \Rightarrow\ \kappa=\frac1a,\quad\rho=a . \]
Sürat $\upsilon=a$ sabit $\Rightarrow\frac{d\upsilon}{dt}=0$. \
$\mathbf a(t)=\mathbf r''(t)=-a\cos t\,\mathbf i-a\sin t\,\mathbf j=a\,\mathbf N=\dfrac{\upsilon^2}{\rho}\mathbf N$ ($a^2/a=a$ \checkmark).
\end{cevap}

%======================================================================
\subsection{Düzgün eğri, integral ve yay uzunluğu}

\begin{soru}[Soru 24 \foto{IMG-20261015-WA0038}, \foto{IMG-20261023-WA0053}]
$\mathbf r(t)=a(t-\sin t)\,\mathbf i+a(1-\cos t)\,\mathbf j$, $-2\pi\le t\le2\pi$ eğrisi düzgün müdür, parçalı düzgün müdür?
\end{soru}
\begin{cevap}[Çözüm (adım adım)]
\textbf{Adım 1.} $\mathbf r'(t)=a(1-\cos t)\,\mathbf i+a\sin t\,\mathbf j$ süreklidir.

\textbf{Adım 2.} $\mathbf r'=\mathbf 0\iff\cos t=1$ ve $\sin t=0\iff t=2k\pi$. Aralıkta: $t=-2\pi,\,0,\,2\pi$.

\textbf{Adım 3.} $t=0$ iç nokta olduğundan eğri $[-2\pi,2\pi]$'de düzgün değildir. Ama $[-2\pi,0]$ ve $[0,2\pi]$ alt aralıklarının
iç noktalarında $\mathbf r'\neq\mathbf 0$: her biri düzgündür. Eğri \textbf{parçalı düzgündür}; $x=-2\pi a,0,2\pi a$ sivri uçlardır.
\end{cevap}

\begin{soru}[Soru 25 \foto{IMG-20261015-WA0037}, \foto{IMG-20261023-WA0055}]
$\mathbf r(t)=2t\,\mathbf i+e^t\,\mathbf j+4\sin2t\,\mathbf k$ fonksiyonunun belirsiz integralini bulunuz.
\end{soru}
\begin{cevap}
\[ \int\mathbf r\,dt=t^2\,\mathbf i+e^t\,\mathbf j-2\cos2t\,\mathbf k+\mathbf C=(t^2+c_1)\,\mathbf i+(e^t+c_2)\,\mathbf j+(-2\cos2t+c_3)\,\mathbf k . \]
($4\int\sin2t\,dt=4\cdot\big(-\frac12\cos2t\big)=-2\cos2t$.)
\end{cevap}

\begin{soru}[Soru 26 \foto{IMG-20261015-WA0037}, \foto{IMG-20261023-WA0055}]
$\mathbf r'(t)=2t\,\mathbf i+\sin t\,\mathbf j+e^t\,\mathbf k$ ve $\mathbf r(0)=2\,\mathbf i+2\,\mathbf j$ koşullarını sağlayan $\mathbf r(t)$'yi bulunuz.
\end{soru}
\begin{cevap}[Çözüm (adım adım)]
\textbf{Adım 1.} $\mathbf r(t)=t^2\,\mathbf i-\cos t\,\mathbf j+e^t\,\mathbf k+\mathbf C$.

\textbf{Adım 2.} $t=0$: $-\mathbf j+\mathbf k+\mathbf C=2\mathbf i+2\mathbf j$ $\Rightarrow$ $\mathbf C=2\mathbf i+3\mathbf j-\mathbf k$.

\textbf{Adım 3.} $\boxed{\mathbf r(t)=(t^2+2)\,\mathbf i+(3-\cos t)\,\mathbf j+(e^t-1)\,\mathbf k}$. \
Kontrol: $\mathbf r(0)=2\mathbf i+2\mathbf j+0\mathbf k$ \checkmark
\end{cevap}

\begin{soru}[Soru 27 \foto{IMG-20261015-WA0039}, \foto{IMG-20261023-WA0040}]
$\mathbf F(t)=2t\,\mathbf i+(1+t^2)\,\mathbf j+\sin\pi t\,\mathbf k$ için $\displaystyle\int_0^1\mathbf F(t)\,dt$ nedir?
\end{soru}
\begin{cevap}
\[ \int_0^1\mathbf F\,dt=t^2\Big|_0^1\mathbf i+\Big(t+\frac{t^3}{3}\Big)\Big|_0^1\mathbf j-\frac{\cos\pi t}{\pi}\Big|_0^1\mathbf k
=\mathbf i+\frac43\,\mathbf j-\frac{-1-1}{\pi}\,\mathbf k=\mathbf i+\frac43\,\mathbf j+\frac2\pi\,\mathbf k . \]
\end{cevap}

\begin{soru}[Soru 28 \foto{IMG-20261015-WA0039}, \foto{IMG-20261023-WA0040}, \foto{IMG-20261023-WA0058}]
$\mathbf r(t)=5\cos t\,\mathbf i+5\sin t\,\mathbf j+12t\,\mathbf k$, $0\le t\le2\pi$ helis parçasının uzunluğu nedir?
\end{soru}
\begin{cevap}[Çözüm (adım adım)]
\textbf{Adım 1.} $\mathbf r'(t)=-5\sin t\,\mathbf i+5\cos t\,\mathbf j+12\,\mathbf k$.

\textbf{Adım 2.} $\|\mathbf r'(t)\|=\sqrt{25\sin^2t+25\cos^2t+144}=\sqrt{169}=13$.

\textbf{Adım 3.} $\ell=\int_0^{2\pi}13\,dt=26\pi$.

\textbf{Sağlama (tahtadaki gibi):} Silindir açılınca helis, kenarları $2\pi\cdot5=10\pi$ ve $12\cdot2\pi=24\pi$ olan dikdörtgenin
köşegenidir: $\sqrt{(10\pi)^2+(24\pi)^2}=26\pi$ \checkmark
\end{cevap}

%======================================================================
\subsection{Çalışma soruları}

Aşağıdaki problemler 15.10 defterinde ve 23.10 tahtasında (\foto{IMG-20261015-WA0039}, \foto{IMG-20261015-WA0040},
\foto{IMG-20261023-WA0058}) çözümsüz olarak verilmiştir.

\begin{soru}[Problem 1 (derste çalışma sorusu olarak verildi) \foto{IMG-20261015-WA0039}, \foto{IMG-20261023-WA0058}]
Türevlerini bulunuz:
(a) $\mathbf F(t)=t\,\mathbf i+t^2\,\mathbf j+t^3\,\mathbf k$; \
(b) $\mathbf F(t)=(t\ln t-t)\,\mathbf i+e^t\,\mathbf j+e^t\,\mathbf k$; \
(c) $\mathbf F(t)=\arcsin(4t)\,\mathbf i-3\arctan(2t-1)\,\mathbf j+7\ln(3t^2)\,\mathbf k$.
\end{soru}
\begin{cevap}[Çözüm (adım adım)]
\textbf{(a)} $\mathbf F'(t)=\mathbf i+2t\,\mathbf j+3t^2\,\mathbf k$.

\textbf{(b)} $\frac{d}{dt}(t\ln t-t)=\ln t+t\cdot\frac1t-1=\ln t$ (çarpım kuralı). Böylece $\mathbf F'(t)=\ln t\,\mathbf i+e^t\,\mathbf j+e^t\,\mathbf k$ \ ($t>0$).

\textbf{(c)} $\frac{d}{dt}\arcsin u=\frac{u'}{\sqrt{1-u^2}}$, $\frac{d}{dt}\arctan u=\frac{u'}{1+u^2}$, $\frac{d}{dt}\ln u=\frac{u'}{u}$:
\[ \mathbf F'(t)=\frac{4}{\sqrt{1-16t^2}}\,\mathbf i-\frac{6}{1+(2t-1)^2}\,\mathbf j+\frac{14}{t}\,\mathbf k\qquad\Big(\tfrac{7\cdot6t}{3t^2}=\tfrac{14}{t}\Big). \]
\textit{Not:} 15.10 defterinde (b)'nin $\mathbf k$ bileşenindeki üs belirsizdir (``$e^{-t}$'' de olabilir); 23.10 tahtasında açıkça $e^t$
okunmaktadır. $e^{-t}$ olsaydı türevi $-e^{-t}\,\mathbf k$ olurdu.
\end{cevap}

\begin{soru}[Problem 2 (derste çalışma sorusu olarak verildi) \foto{IMG-20261015-WA0039}, \foto{IMG-20261023-WA0058}]
$\mathbf F$ iki kez türevlenebilen bir fonksiyon olsun. $[\mathbf F(t)\times\mathbf F'(t)]'=\mathbf F(t)\times\mathbf F''(t)$ olduğunu gösteriniz.
\end{soru}
\begin{cevap}
Vektörel çarpımın türev kuralıyla
\[ [\mathbf F\times\mathbf F']'=\mathbf F'\times\mathbf F'+\mathbf F\times\mathbf F'' . \]
Bir vektörün kendisiyle vektörel çarpımı sıfırdır ($\sin0=0$): $\mathbf F'\times\mathbf F'=\mathbf 0$. Dolayısıyla
$[\mathbf F\times\mathbf F']'=\mathbf F\times\mathbf F''$. \hfill$\blacksquare$
\end{cevap}

\begin{soru}[Problem 3 (derste çalışma sorusu olarak verildi) \foto{IMG-20261015-WA0040}, \foto{IMG-20261023-WA0058}]
$\mathbf r(t)=e^t\,\mathbf i+e^{-t}\,\mathbf j+\sqrt2\,t\,\mathbf k$ eğrisinin $t=0$ noktasındaki Frenet üçlüsünü hesaplayınız.
\end{soru}
\begin{cevap}[Çözüm (adım adım)]
\textbf{Adım 1 ($\mathbf T$).} $\mathbf r'(t)=(e^t,-e^{-t},\sqrt2)$, $\mathbf r'(0)=(1,-1,\sqrt2)$, $\|\mathbf r'(0)\|=\sqrt{1+1+2}=2$.
\[ \mathbf T(0)=\Big(\frac12,\,-\frac12,\,\frac{\sqrt2}{2}\Big). \]
\textbf{Adım 2 ($\mathbf B$).} $\mathbf r''(t)=(e^t,e^{-t},0)$, $\mathbf r''(0)=(1,1,0)$.
\[ \mathbf r'\times\mathbf r''=\begin{vmatrix}\mathbf i&\mathbf j&\mathbf k\\1&-1&\sqrt2\\1&1&0\end{vmatrix}=(0-\sqrt2)\,\mathbf i-(0-\sqrt2)\,\mathbf j+(1+1)\,\mathbf k=(-\sqrt2,\sqrt2,2), \]
normu $\sqrt{2+2+4}=2\sqrt2$. \quad $\mathbf B(0)=\Big(-\dfrac12,\,\dfrac12,\,\dfrac{\sqrt2}{2}\Big)$.

\textbf{Adım 3 ($\mathbf N$).} $\mathbf N=\mathbf B\times\mathbf T$:
\[ \mathbf N(0)=\begin{vmatrix}\mathbf i&\mathbf j&\mathbf k\\-\frac12&\frac12&\frac{\sqrt2}2\\\frac12&-\frac12&\frac{\sqrt2}2\end{vmatrix}
=\Big(\frac{\sqrt2}{4}+\frac{\sqrt2}{4}\Big)\mathbf i-\Big(-\frac{\sqrt2}{4}-\frac{\sqrt2}{4}\Big)\mathbf j+\Big(\frac14-\frac14\Big)\mathbf k=\Big(\frac{\sqrt2}{2},\frac{\sqrt2}{2},0\Big). \]
\textbf{Sağlama (tanımla).} $\|\mathbf r'(t)\|=\sqrt{e^{2t}+e^{-2t}+2}=e^t+e^{-t}$; bu fonksiyonun $t=0$'daki türevi $0$ olduğundan
$\mathbf T'(0)=\frac{\mathbf r''(0)}{\|\mathbf r'(0)\|}=\frac12(1,1,0)$ ve $\mathbf N(0)=\frac{(1,1,0)}{\sqrt2}$ \checkmark. Ek olarak
$\kappa(0)=\frac{2\sqrt2}{2^3}=\frac{\sqrt2}{4}$.
\end{cevap}

\begin{soru}[Problem 4 (derste çalışma sorusu olarak verildi) \foto{IMG-20261015-WA0040}, \foto{IMG-20261023-WA0058}]
İntegralleri hesaplayınız:
(a) $\displaystyle\int_0^1\big[4t^3\,\mathbf i+3\,\mathbf j+2(t+1)\,\mathbf k\big]dt$; \
(b) $\displaystyle\int_1^2\Big[(6-6t)\,\mathbf i+3\sqrt t\,\mathbf j+\frac4{t^2}\,\mathbf k\Big]dt$; \
(c) $\displaystyle\int_{-\pi/4}^{\pi/4}\big[\sin t\,\mathbf i+(1+\cos t)\,\mathbf j+\sec^2t\,\mathbf k\big]dt$.
\end{soru}
\begin{cevap}[Çözüm (adım adım)]
\textbf{(a)} $\int_0^14t^3dt=1$; $\int_0^13\,dt=3$; $\int_0^12(t+1)dt=[t^2+2t]_0^1=3$. \ Sonuç: $\mathbf i+3\,\mathbf j+3\,\mathbf k$.

\textbf{(b)} $\int_1^2(6-6t)dt=[6t-3t^2]_1^2=0-3=-3$; \ $\int_1^23\sqrt t\,dt=[2t^{3/2}]_1^2=4\sqrt2-2$; \
$\int_1^2\frac4{t^2}dt=\big[-\frac4t\big]_1^2=-2+4=2$. \ Sonuç: $-3\,\mathbf i+(4\sqrt2-2)\,\mathbf j+2\,\mathbf k$.

\textbf{(c)} $\sin t$ tek fonksiyon, simetrik aralıkta integrali $0$; \
$\int_{-\pi/4}^{\pi/4}(1+\cos t)dt=\frac\pi2+[\sin t]_{-\pi/4}^{\pi/4}=\frac\pi2+\sqrt2$; \
$\int_{-\pi/4}^{\pi/4}\sec^2t\,dt=[\tan t]_{-\pi/4}^{\pi/4}=2$. \ Sonuç: $\big(\frac\pi2+\sqrt2\big)\mathbf j+2\,\mathbf k$.
\end{cevap}

%======================================================================
\subsection{Ek alıştırmalar}

\begin{soru}[\ekal Ek alıştırma 1]
$\mathbf F(t)=\sqrt{t-1}\,\mathbf i+\dfrac{1}{t-3}\,\mathbf j+\ln(5-t)\,\mathbf k$ fonksiyonunun tanım kümesini bulunuz.
\end{soru}
\begin{cevap}
$T_f=[1,\infty)$, $T_g=\mathbb R\setminus\{3\}$, $T_h=(-\infty,5)$. Kesişim: $T_{\mathbf F}=[1,3)\cup(3,5)$.
\end{cevap}

\begin{soru}[\ekal Ek alıştırma 2]
$\displaystyle\lim_{t\to0}\Big[\frac{1-\cos t}{t^2}\,\mathbf i+\frac{e^{2t}-1}{t}\,\mathbf j+\frac{\tan t}{t}\,\mathbf k\Big]$ limitini bulunuz.
\end{soru}
\begin{cevap}
Her bileşen $\frac00$ belirsizliğidir; L'Hospital (ya da bilinen limitler):
$\frac{1-\cos t}{t^2}\to\frac{\sin t}{2t}\to\frac12$; \ $\frac{e^{2t}-1}{t}\to\frac{2e^{2t}}{1}\to2$; \ $\frac{\tan t}{t}\to\frac{\sec^2t}{1}\to1$.
Limit: $\frac12\,\mathbf i+2\,\mathbf j+\mathbf k$.
\end{cevap}

\begin{soru}[\ekal Ek alıştırma 3]
$\mathbf r''(t)=6t\,\mathbf i+2\,\mathbf j+e^t\,\mathbf k$, $\mathbf r'(0)=\mathbf j+\mathbf k$, $\mathbf r(0)=\mathbf i$ ise $\mathbf r(t)$'yi bulunuz.
\end{soru}
\begin{cevap}
$\mathbf r'=3t^2\,\mathbf i+2t\,\mathbf j+e^t\,\mathbf k+\mathbf C_1$; $\mathbf r'(0)=\mathbf k+\mathbf C_1=\mathbf j+\mathbf k\Rightarrow\mathbf C_1=\mathbf j$:
$\mathbf r'=3t^2\,\mathbf i+(2t+1)\,\mathbf j+e^t\,\mathbf k$.

$\mathbf r=t^3\,\mathbf i+(t^2+t)\,\mathbf j+e^t\,\mathbf k+\mathbf C_2$; $\mathbf r(0)=\mathbf k+\mathbf C_2=\mathbf i\Rightarrow\mathbf C_2=\mathbf i-\mathbf k$:
\[ \mathbf r(t)=(t^3+1)\,\mathbf i+(t^2+t)\,\mathbf j+(e^t-1)\,\mathbf k . \]
\end{cevap}

\begin{soru}[\ekal Ek alıştırma 4]
$\mathbf r(t)=t^3\,\mathbf i+3t^2\,\mathbf j+6t\,\mathbf k$, $0\le t\le1$ eğrisinin uzunluğunu bulunuz.
\end{soru}
\begin{cevap}
$\mathbf r'=(3t^2,6t,6)$, $\|\mathbf r'\|^2=9t^4+36t^2+36=9(t^2+2)^2$, $\|\mathbf r'\|=3(t^2+2)$.
\[ \ell=\int_0^13(t^2+2)\,dt=3\Big(\frac13+2\Big)=7 . \]
\end{cevap}

\begin{soru}[\ekal Ek alıştırma 5]
$\mathbf r(t)=3\cos t\,\mathbf i+3\sin t\,\mathbf j+4t\,\mathbf k$ helisinin $\mathbf T$, $\mathbf N$, $\mathbf B$ vektörlerini, eğriliğini ve burulmasını bulunuz.
\end{soru}
\begin{cevap}[Çözüm (adım adım)]
$\mathbf r'=(-3\sin t,3\cos t,4)$, $\|\mathbf r'\|=5$: \ $\mathbf T=\frac15(-3\sin t,3\cos t,4)$.

$\mathbf T'=\frac15(-3\cos t,-3\sin t,0)$, $\|\mathbf T'\|=\frac35$: \ $\mathbf N=(-\cos t,-\sin t,0)$.

$\mathbf B=\mathbf T\times\mathbf N=\frac15\begin{vmatrix}\mathbf i&\mathbf j&\mathbf k\\-3\sin t&3\cos t&4\\-\cos t&-\sin t&0\end{vmatrix}=\frac15(4\sin t,-4\cos t,3)$.

$\kappa=\dfrac{\|\mathbf T'\|}{\|\mathbf r'\|}=\dfrac{3/5}{5}=\dfrac{3}{25}$ \ (helis formülü $\frac{a}{a^2+b^2}=\frac{3}{25}$ \checkmark). \
$\tau=\dfrac{b}{a^2+b^2}=\dfrac{4}{25}$ \ (kontrol: $(\mathbf r'\times\mathbf r'')\cdot\mathbf r'''=a^2b=36$, $\|\mathbf r'\times\mathbf r''\|^2=a^2(a^2+b^2)=225$, $36/225=4/25$ \checkmark).
\end{cevap}

\begin{soru}[\ekal Ek alıştırma 6]
$\mathbf r(t)=t\,\mathbf i+2t\,\mathbf j+t^2\,\mathbf k$ boyunca hareket eden cismin ivmesinin teğet ve normal bileşenlerini bulunuz.
\end{soru}
\begin{cevap}
$\mathbf v=(1,2,2t)$, $\mathbf a=(0,0,2)$, $\|\mathbf v\|=\sqrt{5+4t^2}$.
\[ a_T=\frac{\mathbf v\cdot\mathbf a}{\|\mathbf v\|}=\frac{4t}{\sqrt{5+4t^2}},\qquad
\mathbf v\times\mathbf a=(4,-2,0),\quad a_N=\frac{\|\mathbf v\times\mathbf a\|}{\|\mathbf v\|}=\frac{2\sqrt5}{\sqrt{5+4t^2}} . \]
Kontrol: $a_T^2+a_N^2=\frac{16t^2+20}{5+4t^2}=4=\|\mathbf a\|^2$ \checkmark
\end{cevap}

\begin{soru}[\ekal Ek alıştırma 7]
$y=\ln x$ eğrisinin eğriliğini bulunuz. Eğrilik hangi noktada en büyüktür?
\end{soru}
\begin{cevap}
$f'=\frac1x$, $f''=-\frac1{x^2}$:
\[ \kappa(x)=\frac{1/x^2}{(1+1/x^2)^{3/2}}=\frac{x}{(x^2+1)^{3/2}}\quad(x>0). \]
$\kappa'(x)=\dfrac{(x^2+1)-3x^2}{(x^2+1)^{5/2}}=\dfrac{1-2x^2}{(x^2+1)^{5/2}}=0\Rightarrow x=\dfrac1{\sqrt2}$ (işaret $+$'dan $-$'ye döner: maksimum).
En büyük eğrilik $\big(\frac1{\sqrt2},-\frac12\ln2\big)$ noktasındadır: $\kappa_{\max}=\dfrac{1/\sqrt2}{(3/2)^{3/2}}=\dfrac{2\sqrt3}{9}$.
\end{cevap}

\begin{soru}[\ekal Ek alıştırma 8]
Bir top yerden $v_0=20$ m/s süratle, yatayla $60^\circ$ açıyla atılıyor ($g\approx10$ m/s$^2$). Konum vektörünü, uçuş süresini, menzili ve
en büyük yüksekliği bulunuz.
\end{soru}
\begin{cevap}
$\mathbf a=-10\,\mathbf j$, $\mathbf v(0)=20\cos60^\circ\,\mathbf i+20\sin60^\circ\,\mathbf j=10\,\mathbf i+10\sqrt3\,\mathbf j$, $\mathbf r(0)=\mathbf 0$.
İki kez integral: $\mathbf r(t)=10t\,\mathbf i+(10\sqrt3\,t-5t^2)\,\mathbf j$.
\begin{itemize}
\item Uçuş süresi: $10\sqrt3\,t-5t^2=0\Rightarrow t_u=2\sqrt3\approx3{,}46$ s.
\item Menzil: $x(2\sqrt3)=20\sqrt3\approx34{,}6$ m \ ($\frac{v_0^2\sin120^\circ}{g}=\frac{400\cdot\sqrt3/2}{10}$ \checkmark).
\item En büyük yükseklik: $y'=10\sqrt3-10t=0\Rightarrow t=\sqrt3$, $H=30-15=15$ m.
\end{itemize}
\end{cevap}

\clearpage\section{Konu 4: Çok Değişkenli Fonksiyonlar, Limit ve Süreklilik}

Bu bölümde derste (15.10 ve 22.10 defter notları, 23.10 ve 08.11 tahta fotoğrafları) çözülen bütün örnekler, işlenme
sırasıyla ve adım adım çözümleriyle verilmiştir. Aynı örnek birden çok fotoğrafta geçiyorsa bir kez yazılmış, bütün
kaynak fotoğraflar belirtilmiştir. Sonda kaynak kitaplardaki tipik alıştırmalara benzer ek alıştırmalar vardır.

\begin{bilgi}[Çalışma soruları hakkında]
Bu konu için tahtada ayrıca ``çalışma sorusu / problemler'' listesi verilmemiştir. \foto{IMG-20261015-WA0040} ve
\foto{IMG-20261023-WA0042}, \foto{IMG-20261023-WA0043} kenarlarında görünen problem listesi (Frenet üçlüsü, vektör
integralleri) Konu 3'e aittir ve orada çözülmüştür.
\end{bilgi}

% =====================================================================
\subsection{Fonksiyon kavramı, tanım ve görüntü kümeleri}

\begin{soru}[Soru 1 \foto{IMG-20261015-WA0040}, \foto{IMG-20261023-WA0043}]
Aşağıdaki büyüklükleri fonksiyon olarak yazınız ve kaç değişkenli olduklarını belirtiniz:
(a) tabanı $a$, yüksekliği $h$ olan üçgenin alanı; (b) $A$ anaparanın $t$ faiz oranıyla $n$ yıllık basit faizi;
(c) kenarları $x,y$ olan dikdörtgenin alanı; (d) boyutları $x,y,h$ olan kutunun hacmi.
\end{soru}
\begin{cevap}
\begin{description}
  \item[(a)] $A(a,h)=\dfrac{a\cdot h}{2}$: iki bağımsız değişken ($a,h$), \textbf{iki değişkenli}. Fiziksel anlamı
  gereği tanım kümesi $a>0,\ h>0$ (birinci bölgenin içi).
  \item[(b)] $f(A,n,t)=\dfrac{A\cdot n\cdot t}{100}$ ($t$ yüzde olarak): \textbf{üç değişkenli}.
  \item[(c)] $f(x,y)=x\cdot y$: \textbf{iki değişkenli}.
  \item[(d)] $f(x,y,h)=x\cdot y\cdot h$: \textbf{üç değişkenli}.
\end{description}
Genel olarak $D\subset\mathbb R^2$ ise $z=f(x,y)$, $D\subset\mathbb R^3$ ise $w=f(x,y,z)$ yazılır; her girdi noktasına
\emph{tek} bir reel sayı karşılık gelir.
\end{cevap}

\begin{soru}[Soru 2 \foto{IMG-20261015-WA0041}, \foto{IMG-20261023-WA0042}]
$f(x,y)=\sqrt{4-x^2-y^2}$ eşitliği ile tanımlanan $f$ reel değerli fonksiyonunun tanım kümesini analitik düzlemde
gösteriniz.
\end{soru}
\begin{cevap}[Çözüm (adım adım)]
\textbf{1. Koşulu yaz.} Karekökün içi negatif olamaz: $4-x^2-y^2\ge0$.

\textbf{2. Düzenle.} $x^2+y^2\le4=2^2$.

\textbf{3. Geometrik yorum.} $x^2+y^2=4$ merkezi orijin, yarıçapı $2$ olan çemberdir; ``$\le$'' çemberin içi ve kendisi
demektir:
\[
T_f=\{(x,y): x^2+y^2\le4\}\quad\text{(kapalı disk, sınır çember dahil)}.
\]
\textbf{4. Çizim.} Çember sürekli çizgiyle çizilir, içi taranır.

\textbf{Ek (görüntü kümesi).} Disk içinde $0\le4-x^2-y^2\le4$ olduğundan $0\le f\le2$; orijinde $2$, çember üzerinde
$0$ değeri alınır ve arada her değer alınır: $G_f=[0,2]$. Grafik, yarıçapı $2$ olan kürenin üst yarısıdır.
\end{cevap}

\begin{soru}[Soru 3 \foto{IMG-20261015-WA0041}, \foto{IMG-20261023-WA0042}]
$f(x,y)=\arccos(x-y)$ fonksiyonunun tanım ve görüntü kümesi nedir?
\end{soru}
\begin{cevap}[Çözüm (adım adım)]
\textbf{1. Hatırlatma.} $\cos:\mathbb R\to[-1,1]$; $\arccos:[-1,1]\to[0,\pi]$. Yani $\arccos$'un girdisi $[-1,1]$
aralığında olmalıdır.

\textbf{2. Koşul.} $-1\le x-y\le1$.

\textbf{3. İki eşitsizliğe ayır.}
\[
x-y\ge-1\ \Rightarrow\ y\le x+1, \qquad x-y\le1\ \Rightarrow\ y\ge x-1 .
\]
\textbf{4. Tanım kümesi.}
\[
T_f=\{(x,y): x-1\le y\le x+1\}:
\]
eğimi $1$ olan paralel $y=x-1$ ve $y=x+1$ doğruları arasındaki şerit; sınır doğruları dahil.

\textbf{5. Görüntü kümesi.} Şerit üzerinde $u=x-y$ ifadesi $[-1,1]$'in her değerini alır (ör. $(x,y)=(u,0)$).
$\arccos u$ bu aralıkta $\arccos(-1)=\pi$'den $\arccos1=0$'a sürekli ve azalan biçimde iner. Dolayısıyla
\[
G_f=[0,\pi].
\]
\emph{Not:} Tahtadaki $\arccos u$ çizimi kabataslaktır ve etiketi \okunamadi{``arccos4'' gibi görünüyor; büyük olasılıkla
``$\arccos u$''}; gerçek eğri $(0,\pi/2)$'den geçer.
\end{cevap}

\begin{soru}[Soru 4 \foto{IMG-20261015-WA0042}, \foto{IMG-20261023-WA0044}]
$f(x,y)=\sqrt{xy}+\log(1-x^2-y^2)$ şeklinde tanımlı $f$ fonksiyonunun tanım kümesi nedir?
\end{soru}
\begin{cevap}[Çözüm (adım adım)]
\textbf{1. Kök koşulu.} $xy\ge0$: çarpım negatif değilse $x$ ve $y$ aynı işaretli ya da biri sıfırdır:
$(x\ge0$ ve $y\ge0)$ veya $(x\le0$ ve $y\le0)$. Bu, I. ve III. bölgeler (eksenler dahil).

\textbf{2. Logaritma koşulu.} $1-x^2-y^2>0\Rightarrow x^2+y^2<1$: birim çemberin içi, çember \textbf{hariç} (kesin
eşitsizlik).

\textbf{3. Kesişim.}
\[
T_f=\{(x,y): xy\ge0,\ x^2+y^2<1\}.
\]
Birim dairenin I. ve III. bölgedeki çeyrekleri; eksen parçaları dahil, çember hariç (çizimde çember kesikli).
Orijin de dahildir: $\sqrt0+\log1=0$. Logaritmanın tabanı tanım kümesini etkilemez.
\end{cevap}

\begin{soru}[Soru 5 \foto{IMG-20261022-WA0053}, \foto{IMG-20261108-WA0088}]
Aşağıdaki fonksiyonların tanım kümelerini bulunuz.
(a) $f(x,y)=\ln(x^2+y^2+4)$ \quad (b) $f(x,y)=\sqrt{xy-x^4}+xy$ \quad (c) $f(x,y)=\sqrt{x^2-y}+x^2+y^2+1$
\end{soru}
\begin{cevap}[Çözüm (adım adım)]
\begin{description}
  \item[(a)] Logaritma koşulu: $x^2+y^2+4>0$. $x^2+y^2\ge0$ olduğundan sol taraf her zaman en az $4$'tür; koşul her
  $(x,y)$ için doğrudur. $T_f=\mathbb R^2$ (bütün düzlem).
  \item[(b)] Kök koşulu: $xy-x^4\ge0$. Ortak çarpan $x$: $x(y-x^3)\ge0$. Çarpımın işaret tablosu:
  \begin{itemize}
    \item (1) $x\ge0$ ve $y-x^3\ge0$, yani $x\ge0$ ve $y\ge x^3$: sağ yarı düzlemde $y=x^3$ eğrisinin üstü;
    \item (2) $x\le0$ ve $y-x^3\le0$, yani $x\le0$ ve $y\le x^3$: sol yarı düzlemde $y=x^3$ eğrisinin altı.
  \end{itemize}
  $T_f$ = (1) $\cup$ (2); sınırlar ($y=x^3$ eğrisi ve $y$ ekseninin tamamı) dahildir. Kontrol: $(1,2)$ için
  $2-1=1\ge0$ \checkmark; $(1,0)$ için $0-1<0$, dışarıda \checkmark; $(-1,-2)$ için $2-1=1\ge0$ \checkmark.
  \item[(c)] Kök koşulu: $x^2-y\ge0\Rightarrow y\le x^2$. Polinom kısmı her yerde tanımlı. $T_f$: $y=x^2$ parabolünün
  altı ve kendisi. Kontrol: $(0,-1)$: $0+1=1\ge0$ \checkmark; $(0,1)$: $-1<0$, dışarıda \checkmark.
\end{description}
\end{cevap}

% =====================================================================
\subsection{Seviye eğrileri ve grafik}

\begin{soru}[Soru 6 \foto{IMG-20261015-WA0042}, \foto{IMG-20261015-WA0043}, \foto{IMG-20261023-WA0044}]
$f(x,y)=x^2+y^2$ fonksiyonunun bazı seviye eğrilerini bulunuz. Bundan yararlanarak fonksiyon grafiğini çiziniz.
\end{soru}
\begin{cevap}[Çözüm (adım adım)]
\textbf{1. Seviye eğrisi denklemi.} $f(x,y)=c\ \Rightarrow\ x^2+y^2=c$.

\textbf{2. $c$'ye göre durum.}
\begin{itemize}
  \item $c<0$: $x^2+y^2\ge0$ olduğundan çözüm yok (boş küme).
  \item $c=0$: yalnızca $(0,0)$.
  \item $c>0$: merkezi orijin, yarıçapı $\sqrt c$ olan çember. Örneğin $c=1$: birim çember; $c=4$: yarıçapı $2$;
        $c=9$: yarıçapı $3$.
\end{itemize}
\textbf{3. Grafiği kurma.} Yüzeyin $z=c$ yüksekliğindeki kesiti yarıçapı $\sqrt c$ olan yatay bir çemberdir. Yükseklik
arttıkça çemberler büyür; $x=0$ düzlemiyle kesit $z=y^2$, $y=0$ düzlemiyle kesit $z=x^2$ parabolüdür. Bu bilgiler
birleşince orijinden yukarı açılan kâse biçimli \textbf{paraboloid} $z=x^2+y^2$ elde edilir.

\textbf{Gösterim notu.} Tahtada ``$x^2+y^2=c^2$ alındığında yarıçapı $c$ olan çember'' yazılmıştır; bu, seviye değeri
$c^2$ seçildiğinde doğrudur. $f=c$ yazımıyla doğrusu yarıçapın $\sqrt c$ olmasıdır; tahtadaki şekildeki $c=1,4,9$
etiketleri (yarıçap $1,2,3$) de bunu gösterir. Ayrıca tahtadaki paraboloid düz kenarlı (koni gibi) çizilmiştir;
gerçek yüzeyin kenarları eğridir.
\end{cevap}

% =====================================================================
\subsection{Limit hesabı ve limitin yokluğu}

\begin{soru}[Soru 7 \foto{IMG-20261022-WA0058}, \foto{IMG-20261108-WA0094}]
$\displaystyle\lim_{(x,y)\to(0,0)}\frac{2xy}{x^2+y^2}$ limitini inceleyiniz.
\end{soru}
\begin{cevap}[Çözüm (adım adım)]
\textbf{1. Doğrudan yerine koyma} $\frac00$ belirsizliği verir; fonksiyon orijinde tanımsız.

\textbf{2. Kutupsal koordinatlar.} $x=r\cos\theta$, $y=r\sin\theta$; $(x,y)\to(0,0)\Leftrightarrow r\to0$:
\[
\frac{2xy}{x^2+y^2}=\frac{2r^2\cos\theta\sin\theta}{r^2(\cos^2\theta+\sin^2\theta)}=2\cos\theta\sin\theta=\sin2\theta .
\]
\textbf{3. Yorum.} Sonuç $r$'den bağımsız ama $\theta$'ya bağlı: $\theta=0$ doğrultusunda $0$, $\theta=\pi/4$
doğrultusunda $1$, $\theta=3\pi/4$ doğrultusunda $-1$. Farklı doğrultularda farklı değer $\Rightarrow$ \textbf{limit
yoktur}.

\textbf{4. Doğrularla doğrulama.} $y=mx$:
\[
\frac{2x\cdot mx}{x^2+m^2x^2}=\frac{2m}{1+m^2};\qquad m=0\Rightarrow0,\quad m=1\Rightarrow1 .
\]
İki farklı yol, iki farklı değer: limit yoktur.
\end{cevap}

\begin{soru}[Soru 8 \foto{IMG-20261022-WA0053}]
$P(x,y)=5x^2y^2+4xy^3-8xy^2+3$ için $\displaystyle\lim_{(x,y)\to(1,2)}P(x,y)$ değerini bulunuz.
\end{soru}
\begin{cevap}
Polinomlar için $\lim_{(x,y)\to(a,b)}x^my^n=a^mb^n$ ve limit kurallarından $\lim P(x,y)=P(a,b)$ (doğrudan yerine koyma):
\[
P(1,2)=5\cdot1\cdot4+4\cdot1\cdot8-8\cdot1\cdot4+3=20+32-32+3=\boxed{23}.
\]
\end{cevap}

\begin{soru}[Soru 9 \foto{IMG-20261022-WA0053}, \foto{IMG-20261022-WA0056}, \foto{IMG-20261108-WA0088}, \foto{IMG-20261108-WA0089}]
Aşağıdaki limitleri hesaplayınız.
\begin{description}
  \item[(a)] $\displaystyle\lim_{(x,y)\to(0,0)}\frac{2x^2+y^2+1}{x^2+y^2+3}$ \qquad
  (b) $\displaystyle\lim_{(x,y)\to(0,0)}\frac{2-\sqrt{xy+4}}{xy}$
  \item[(c)] $\displaystyle\lim_{(x,y)\to(0,0)}\frac{\sin(xy)}{x}$ \qquad
  (d) $\displaystyle\lim_{(x,y)\to(0,\ln2)}e^{x-y}$
  \item[(e)] $\displaystyle\lim_{(x,y)\to(1,0)}\frac{x\sin y}{x^2+1}$ \qquad
  (f) $\displaystyle\lim_{(x,y)\to(0,0)}\frac{e^y\sin x}{x}$
\end{description}
\end{soru}
\begin{cevap}[Çözüm (adım adım)]
\begin{description}
  \item[(a)] Payda $(0,0)$'da $3\neq0$. Rasyonel fonksiyon, doğrudan yerine koyma:
  $\dfrac{0+0+1}{0+0+3}=\dfrac13$.
  \item[(b)] \textbf{1.} Yerine koyma: $\frac{2-\sqrt4}{0}=\frac00$, belirsiz.
  \textbf{2.} Payın eşleniği $2+\sqrt{xy+4}$ ile pay ve paydayı çarp ($xy\neq0$ noktalarında çalışıyoruz):
  \[
  \frac{(2-\sqrt{xy+4})(2+\sqrt{xy+4})}{xy\,(2+\sqrt{xy+4})}=\frac{4-(xy+4)}{xy\,(2+\sqrt{xy+4})}
  =\frac{-xy}{xy\,(2+\sqrt{xy+4})}=\frac{-1}{2+\sqrt{xy+4}} .
  \]
  \textbf{3.} Artık payda sıfırdan farklı; yerine koy: $\dfrac{-1}{2+\sqrt4}=-\dfrac14$.
  \item[(c)] \textbf{1.} $\frac00$ belirsizliği. \textbf{2.} Pay ve paydayı $y$ ile çarpıp $u=xy\to0$ kullan:
  \[
  \frac{\sin(xy)}{x}=y\cdot\frac{\sin(xy)}{xy}\ \Rightarrow\
  \lim\frac{\sin(xy)}{xy}\cdot\lim y=1\cdot0=0 .
  \]
  \textbf{Titiz alternatif:} $|\sin u|\le|u|$ ile $\left|\frac{\sin(xy)}{x}\right|\le\frac{|x||y|}{|x|}=|y|\to0$;
  sıkıştırma ile limit $0$ (bu yol $y=0$ noktalarında tanımsız bölme yaratmaz).
  \item[(d)] $e^{x-y}=e^x\cdot e^{-y}$:
  \[
  \lim_{x\to0}e^x\cdot\lim_{y\to\ln2}e^{-y}=1\cdot e^{-\ln2}=1\cdot\frac12=\frac12 .
  \]
  \emph{Düzeltme:} 22.10 defter notunda (\foto{IMG-20261022-WA0056}) ikinci çarpan $e^{y}$ yazılmış; doğrusu
  $e^{-y}$'dir (08.11 tahtasında, \foto{IMG-20261108-WA0089}, doğru yazılmış). $e^y$ ile hesaplansaydı $2$ çıkardı.
  Sonuç $\tfrac12$ doğrudur.
  \item[(e)] Payda $x^2+1\ge1$ asla sıfır olmaz; değişkenlere ayır:
  \[
  \lim_{x\to1}\frac{x}{x^2+1}\cdot\lim_{y\to0}\sin y=\frac{1}{2}\cdot0=0 .
  \]
  \item[(f)] $\dfrac{e^y\sin x}{x}=e^y\cdot\dfrac{\sin x}{x}$; iki tek değişkenli limit de var:
  \[
  \lim_{x\to0}\frac{\sin x}{x}\cdot\lim_{y\to0}e^y=1\cdot1=1 .
  \]
  (22.10 defterinde yarım kalan bu şıkkın çözümü 08.11 tahtasından alınmıştır.)
\end{description}
\end{cevap}

\begin{soru}[Soru 10 \foto{IMG-20261022-WA0056}, \foto{IMG-20261022-WA0054}, \foto{IMG-20261108-WA0089}]
Aşağıdaki fonksiyonların $(0,0)$ noktasında limite sahip olmadığını gösteriniz.
\begin{description}
  \item[(a)] $f(x,y)=\dfrac{x^4-y^2}{x^4+y^2}$ \qquad (b) $f(x,y)=\dfrac{x-y}{x+y}$
  \item[(c)] $f(x,y)=\dfrac{x}{\sqrt{x^2+y^2}}$ \qquad (d) $f(x,y)=\dfrac{x^2+y^2}{y^2}$
\end{description}
\end{soru}
\begin{cevap}[Çözüm (adım adım)]
Strateji: iki farklı yol bulup farklı limit değerleri elde etmek.
\begin{description}
  \item[(a)] Paydada $x^4$ ve $y^2$ olduğu için $y=ax^2$ parabolleri $x^4$ ile $y^2$'yi aynı mertebeye getirir:
  \[
  \lim_{x\to0}\frac{x^4-a^2x^4}{x^4+a^2x^4}=\frac{1-a^2}{1+a^2}.
  \]
  $a=0$ ($x$ ekseni) $\Rightarrow1$; $a=1$ ($y=x^2$) $\Rightarrow0$. Farklı $\Rightarrow$ limit yok.
  \item[(b)] $y=mx$ ($m\neq-1$, çünkü $y=-x$ üzerinde payda sıfır):
  \[
  \lim_{x\to0}\frac{x-mx}{x+mx}=\lim_{x\to0}\frac{x(1-m)}{x(1+m)}=\frac{1-m}{1+m}.
  \]
  $m=0\Rightarrow1$, $m=1\Rightarrow0$. Limit yok.
  \item[(c)] Kutupsal ($r>0$ olduğundan $\sqrt{r^2}=r$):
  \[
  \frac{r\cos\theta}{\sqrt{r^2(\cos^2\theta+\sin^2\theta)}}=\frac{r\cos\theta}{r}=\cos\theta .
  \]
  $\theta=0$'da $1$, $\theta=\pi/2$'de $0$, $\theta=\pi$'de $-1$; $\theta$'ya bağlı $\Rightarrow$ limit yok.
  \item[(d)] $y=mx$ ($m\neq0$; $y=0$ üzerinde $f$ tanımsız):
  \[
  \lim_{x\to0}\frac{x^2+m^2x^2}{m^2x^2}=\frac{x^2(1+m^2)}{x^2m^2}=\frac{1+m^2}{m^2}.
  \]
  $m=1\Rightarrow2$, $m=2\Rightarrow\frac54$. Limit yok.
\end{description}
\end{cevap}

% =====================================================================
\subsection{Süreklilik}

\begin{soru}[Soru 11 \foto{IMG-20261022-WA0054}, \foto{IMG-20261108-WA0090}]
\[
f(x,y)=\begin{cases}\dfrac{\sin(xy^2)}{x^2+y^2}, & (x,y)\neq(0,0)\\[6pt] 0, & (x,y)=(0,0)\end{cases}
\]
fonksiyonunun her yerde sürekli olduğunu gösteriniz.
\end{soru}
\begin{cevap}[Çözüm (adım adım)]
Süreklilik tanımı: $\lim_{(x,y)\to(a,b)}f(x,y)=f(a,b)$ (\foto{IMG-20261022-WA0053}, \foto{IMG-20261108-WA0088}).

\textbf{1. $(x,y)\neq(0,0)$.} $\sin(xy^2)$, sürekli $\sin$ ile $xy^2$ polinomunun bileşkesi olduğundan süreklidir;
$x^2+y^2$ polinomdur ve bu noktalarda payda sıfır değildir. Bölüm sürekli $\Rightarrow f$ süreklidir.

\textbf{2. $(0,0)$ noktası.} $f(0,0)=0$ tanımlı. Limit:
\[
\lim_{(x,y)\to(0,0)}\frac{\sin(xy^2)}{x^2+y^2}=\lim\frac{\sin(xy^2)}{xy^2}\cdot\lim\frac{xy^2}{x^2+y^2}.
\]
Birinci çarpan: $u=xy^2\to0$ ve $\frac{\sin u}{u}\to1$. İkinci çarpan:
\begin{itemize}
  \item[(1)] $y=mx$: $\dfrac{x\cdot m^2x^2}{x^2+m^2x^2}=\dfrac{m^2x}{1+m^2}\to0$.
  \item[(2)] Kutupsal: $\dfrac{r\cos\theta\cdot r^2\sin^2\theta}{r^2}=r\cos\theta\sin^2\theta$ ve
  $\boxed{|r\cos\theta\sin^2\theta|\le r}\to0$.
\end{itemize}
Limit $=1\cdot0=0=f(0,0)$; $f$ orijinde de süreklidir. \textbf{Sonuç:} $f$ her yerde süreklidir.

\textbf{Titizlik notu.} (1)'deki doğru yolu tek başına kanıt değildir (doğrularda $0$ çıkıp eğride farklı değer çıkan
fonksiyonlar vardır); kanıtı $\theta$'dan bağımsız (2) sınırı tamamlar. Ayrıca $\frac{\sin(xy^2)}{xy^2}$ eksenler
üzerinde tanımsızdır. Bölme yapmadan tek satırlık kanıt: $|\sin u|\le|u|$ ile
\[
|f(x,y)|\le\frac{|x|\,y^2}{x^2+y^2}\le|x|\to0 .
\]
\end{cevap}

% =====================================================================
\subsection{Ek alıştırmalar}

\begin{soru}[\ekal Ek alıştırma 1]
$f(x,y)=\sqrt{1-x^2}+\sqrt{1-y^2}$ fonksiyonunun tanım ve görüntü kümesini bulunuz.
\end{soru}
\begin{cevap}
Kök koşulları: $1-x^2\ge0\Rightarrow-1\le x\le1$ ve $1-y^2\ge0\Rightarrow-1\le y\le1$.
$T_f=[-1,1]\times[-1,1]$: köşeleri $(\pm1,\pm1)$ olan kapalı kare.

Her terim $[0,1]$ aralığında değer alır; toplam $[0,2]$ içindedir. $f(\pm1,\pm1)=0$, $f(0,0)=2$; $y=0$ alıp
$x$'i $1$'den $0$'a götürdüğümüzde $f=\sqrt{1-x^2}+1$ değeri $1$'den $2$'ye, $x=1$ alıp $y$'yi değiştirdiğimizde
$0$'dan $1$'e sürekli değişir. Böylece $G_f=[0,2]$.
\end{cevap}

\begin{soru}[\ekal Ek alıştırma 2]
$f(x,y)=x^2-y^2$ fonksiyonunun seviye eğrilerini $c=-1,0,1$ için bulunuz ve grafiğin biçimini yorumlayınız.
\end{soru}
\begin{cevap}
\begin{itemize}
  \item $c=1$: $x^2-y^2=1$; tepe noktaları $(\pm1,0)$ olan, $x$ ekseni boyunca açılan hiperbol.
  \item $c=0$: $x^2=y^2\Rightarrow y=\pm x$; orijinde kesişen iki doğru.
  \item $c=-1$: $y^2-x^2=1$; tepe noktaları $(0,\pm1)$ olan, $y$ ekseni boyunca açılan hiperbol.
\end{itemize}
Grafik $z=x^2-y^2$, $x$ yönünde yukarı ($z=x^2$), $y$ yönünde aşağı ($z=-y^2$) açılan bir \textbf{eyer yüzeyi}
(hiperbolik paraboloid) belirtir.
\end{cevap}

\begin{soru}[\ekal Ek alıştırma 3]
$\varepsilon$--$\delta$ tanımıyla $\displaystyle\lim_{(x,y)\to(1,2)}(2x+3y)=8$ olduğunu gösteriniz.
\end{soru}
\begin{cevap}
$\varepsilon>0$ verilsin ve $d=\sqrt{(x-1)^2+(y-2)^2}$ olsun. $|x-1|\le d$ ve $|y-2|\le d$ olduğundan
\[
|2x+3y-8|=|2(x-1)+3(y-2)|\le2|x-1|+3|y-2|\le5d .
\]
$\delta=\varepsilon/5$ seçilirse $0<d<\delta$ iken $|2x+3y-8|<5\delta=\varepsilon$ olur. Limit $8$'dir.
\end{cevap}

\begin{soru}[\ekal Ek alıştırma 4]
Limitleri hesaplayınız:
(a) $\displaystyle\lim_{(x,y)\to(1,1)}\frac{x^3-y^3}{x-y}$ \qquad
(b) $\displaystyle\lim_{(x,y)\to(0,0)}\frac{x^2+y^2}{\sqrt{x^2+y^2+1}-1}$
\end{soru}
\begin{cevap}
\begin{description}
  \item[(a)] $\frac00$ belirsizliği ($x\neq y$). $x^3-y^3=(x-y)(x^2+xy+y^2)$:
  \[
  \lim_{(x,y)\to(1,1)}(x^2+xy+y^2)=1+1+1=3 .
  \]
  \item[(b)] $\frac00$ belirsizliği. Paydanın eşleniği $\sqrt{x^2+y^2+1}+1$ ile çarp:
  \[
  \frac{(x^2+y^2)\big(\sqrt{x^2+y^2+1}+1\big)}{(x^2+y^2+1)-1}=\sqrt{x^2+y^2+1}+1\ \longrightarrow\ 1+1=2 .
  \]
\end{description}
\end{cevap}

\begin{soru}[\ekal Ek alıştırma 5]
Aşağıdaki limitlerin olmadığını gösteriniz:
(a) $\displaystyle\lim_{(x,y)\to(0,0)}\frac{x^2-y^2}{x^2+y^2}$ \qquad
(b) $\displaystyle\lim_{(x,y)\to(0,0)}\frac{xy^2}{x^2+y^4}$
\end{soru}
\begin{cevap}
\begin{description}
  \item[(a)] $x$ ekseni ($y=0$): $\frac{x^2}{x^2}=1$. $y$ ekseni ($x=0$): $\frac{-y^2}{y^2}=-1$. Farklı; limit yok.
  \item[(b)] Doğrular $y=mx$: $\dfrac{m^2x^3}{x^2+m^4x^4}=\dfrac{m^2x}{1+m^4x^2}\to0$; eksenlerde de $0$. Ama $x=y^2$
  parabolü boyunca: $\dfrac{y^2\cdot y^2}{y^4+y^4}=\dfrac12$. Doğrularda $0$, parabolde $\frac12$; limit yok.
  (Bütün doğrularda aynı değerin çıkması limitin varlığı için yetmez!)
\end{description}
\end{cevap}

\begin{soru}[\ekal Ek alıştırma 6]
$\displaystyle\lim_{(x,y)\to(0,0)}\frac{x^4+y^4}{x^2+y^2}$ limitinin var olduğunu gösterip değerini bulunuz.
\end{soru}
\begin{cevap}
$x^4+y^4\le x^4+2x^2y^2+y^4=(x^2+y^2)^2$ olduğundan $(x,y)\neq(0,0)$ için
\[
0\le\frac{x^4+y^4}{x^2+y^2}\le\frac{(x^2+y^2)^2}{x^2+y^2}=x^2+y^2\to0 .
\]
Sıkıştırma teoremiyle limit $0$. (Kutupsal: $r^2(\cos^4\theta+\sin^4\theta)\le r^2$.)
\end{cevap}

\begin{soru}[\ekal Ek alıştırma 7]
\[
f(x,y)=\begin{cases}\dfrac{e^{x^2+y^2}-1}{x^2+y^2}, & (x,y)\neq(0,0)\\[6pt] k, & (x,y)=(0,0)\end{cases}
\]
fonksiyonu her yerde sürekli olacak şekilde $k$ sabitini bulunuz.
\end{soru}
\begin{cevap}
$(x,y)\neq(0,0)$'da $f$, sürekli fonksiyonların bölümü (payda $\neq0$) olduğundan süreklidir. Orijinde $u=x^2+y^2\to0^+$:
\[
\lim_{(x,y)\to(0,0)}f(x,y)=\lim_{u\to0^+}\frac{e^u-1}{u}=1
\]
(L'Hospital ya da $e^u$'nun $u=0$'daki türevinin $1$ olması). Süreklilik için $k=f(0,0)=1$.
\end{cevap}

\begin{soru}[\ekal Ek alıştırma 8]
$\displaystyle\lim_{(x,y,z)\to(0,0,0)}\frac{x^2+y^2-z^2}{x^2+y^2+z^2}$ limitini inceleyiniz.
\end{soru}
\begin{cevap}
$x$ ekseni ($y=z=0$): $\frac{x^2}{x^2}=1$. $z$ ekseni ($x=y=0$): $\frac{-z^2}{z^2}=-1$. İki farklı yol, iki farklı değer;
limit \textbf{yoktur}.
\end{cevap}

\begin{soru}[\ekal Ek alıştırma 9]
$g(x,y)=\dfrac{x^2y}{x^2+y^2}$ ($(x,y)\neq(0,0)$), $g(0,0)=0$ fonksiyonunun sürekli olduğu kümeyi bulunuz.
\end{soru}
\begin{cevap}
$(x,y)\neq(0,0)$'da rasyonel ve paydası sıfırdan farklı olduğundan süreklidir. Orijinde:
$\dfrac{x^2}{x^2+y^2}\le1$ olduğundan $|g(x,y)|\le|y|\to0=g(0,0)$. O hâlde $g$ \textbf{bütün düzlemde} süreklidir.
\end{cevap}

\clearpage\section{Konu 5: Kısmi Türevler, Teğet Düzlemler ve Lineer Yaklaşımlar}

Bu bölümde 22.10 (defter notu), 05.11 ve 08.11 derslerinde tahtada çözülen tüm sorular derste işlenme sırasıyla ve adım
adım çözülmüştür. 08.11 dersinde 22.10 konuları (kısmi türev, Schwarz teoremi, lineerleştirme) ve 05.11 problemleri
tahtada yeniden yazıldığı için aynı soru birden fazla fotoğrafta görünüyorsa tüm fotoğraflar birlikte verilmiştir.
Bu konuya ait tahtada sorusu yazılıp çözülmeden bırakılmış bir çalışma sorusu yoktur; ``Problemler'' başlığıyla verilen
soruların hepsi derste çözülmüştür. Sonda kaynak kitaplardaki tipik alıştırmalara benzer ek alıştırmalar vardır.

%-------------------------------------------------------------------------------------------
\subsection{Kısmi türevler ve yüksek mertebeden kısmi türevler (22.10 dersi)}

\begin{soru}[Soru 1 \foto{IMG-20261022-WA0055}, \foto{IMG-20261108-WA0090}]
$f(x,y)=x^2y^3+3x$ fonksiyonunun $f_x$ ve $f_y$ kısmi türevlerini bulunuz.
\end{soru}
\begin{cevap}[Çözüm (adım adım)]
\textbf{Kural:} Bir değişkene göre kısmi türev alırken diğer değişken sabit (bir sayı gibi) kabul edilir.

\textbf{$f_x$ ($y$ sabit):}
\begin{itemize}
  \item $x^2y^3$: $y^3$ sabit katsayıdır, $\dfrac{d}{dx}(x^2)=2x$ $\Rightarrow$ $2xy^3$.
  \item $3x$: türevi $3$.
\end{itemize}
\[ f_x=2xy^3+3 . \]
\textbf{$f_y$ ($x$ sabit):}
\begin{itemize}
  \item $x^2y^3$: $x^2$ sabit katsayıdır, $\dfrac{d}{dy}(y^3)=3y^2$ $\Rightarrow$ $3x^2y^2$.
  \item $3x$: $y$ içermediği için $y$'ye göre sabittir; türevi $0$.
\end{itemize}
\[ f_y=3x^2y^2 . \]
\textbf{Ek kontrol (Schwarz teoremi):} $f_{xy}=(2xy^3+3)_y=6xy^2$, $f_{yx}=(3x^2y^2)_x=6xy^2$; karma türevler eşittir
(fonksiyon simetrik olmadığı hâlde).
\end{cevap}

\begin{soru}[Soru 2 \foto{IMG-20261022-WA0055}, \foto{IMG-20261022-WA0059}, \foto{IMG-20261108-WA0091}, \foto{IMG-20261108-WA0092}]
$f(x,y)=\dfrac{xy}{x^2+y^2}$ fonksiyonunun 2. mertebeden kısmi türevlerini hesaplayınız.
\end{soru}
\begin{cevap}[Çözüm (adım adım)]
Fonksiyon $(x,y)\ne(0,0)$ için tanımlıdır; tüm hesaplar bu noktalarda geçerlidir. Bölüm kuralı:
$\left(\dfrac uv\right)'=\dfrac{u'v-uv'}{v^2}$.

\textbf{Adım 1 -- $f_x$.} $y$ sabit: $u=xy\Rightarrow u_x=y$; $v=x^2+y^2\Rightarrow v_x=2x$.
\[
f_x=\frac{y(x^2+y^2)-2x\cdot xy}{(x^2+y^2)^2}=\frac{yx^2+y^3-2x^2y}{(x^2+y^2)^2}=\frac{y^3-x^2y}{(x^2+y^2)^2}.
\]
\textbf{Adım 2 -- $f_{xx}$.} $f_x$'i $x$'e göre türevle: $u=y^3-x^2y\Rightarrow u_x=-2xy$;
$v=(x^2+y^2)^2\Rightarrow v_x=2(x^2+y^2)\cdot2x$.
\[
f_{xx}=\frac{-2xy(x^2+y^2)^2-2(x^2+y^2)\cdot2x\,(y^3-x^2y)}{(x^2+y^2)^4}.
\]
Pay ve payda $(x^2+y^2)$ ile sadeleşir:
\[
f_{xx}=\frac{-2xy(x^2+y^2)-4x(y^3-x^2y)}{(x^2+y^2)^3}=\frac{-2x^3y-2xy^3-4xy^3+4x^3y}{(x^2+y^2)^3}
=\frac{2x^3y-6xy^3}{(x^2+y^2)^3}.
\]
\textbf{Adım 3 -- $f_{xy}$.} $f_x$'i $y$'ye göre türevle: $u=y^3-x^2y\Rightarrow u_y=3y^2-x^2$;
$v_y=2(x^2+y^2)\cdot2y$.
\begin{align*}
f_{xy}&=\frac{(3y^2-x^2)(x^2+y^2)^2-2(x^2+y^2)\cdot2y\,(y^3-x^2y)}{(x^2+y^2)^4}
=\frac{(3y^2-x^2)(x^2+y^2)-4y(y^3-x^2y)}{(x^2+y^2)^3}\\
&=\frac{3x^2y^2+3y^4-x^4-x^2y^2-4y^4+4x^2y^2}{(x^2+y^2)^3}
=\boxed{\frac{-x^4-y^4+6x^2y^2}{(x^2+y^2)^3}} .
\end{align*}
\textbf{Adım 4 -- $f_y$.} $x$ sabit: $u_y=x$, $v_y=2y$.
\[
f_y=\frac{x(x^2+y^2)-2y\cdot xy}{(x^2+y^2)^2}=\frac{x^3-xy^2}{(x^2+y^2)^2}.
\]
\textbf{Adım 5 -- $f_{yx}$.} $f_y$'yi $x$'e göre türevle: $u=x^3-xy^2\Rightarrow u_x=3x^2-y^2$; $v_x=2(x^2+y^2)\cdot2x$.
\begin{align*}
f_{yx}&=\frac{(3x^2-y^2)(x^2+y^2)^2-2(x^2+y^2)\cdot2x\,(x^3-xy^2)}{(x^2+y^2)^4}
=\frac{(3x^2-y^2)(x^2+y^2)-4x(x^3-xy^2)}{(x^2+y^2)^3}\\
&=\frac{3x^4+3x^2y^2-x^2y^2-y^4-4x^4+4x^2y^2}{(x^2+y^2)^3}
=\boxed{\frac{-x^4-y^4+6x^2y^2}{(x^2+y^2)^3}} .
\end{align*}
\textbf{Adım 6 -- $f_{yy}$.} $f_y$'yi $y$'ye göre türevle: $u_y=-2xy$, $v_y=2(x^2+y^2)\cdot2y$.
\[
f_{yy}=\frac{-2xy(x^2+y^2)-4y(x^3-xy^2)}{(x^2+y^2)^3}=\frac{-2x^3y-2xy^3-4x^3y+4xy^3}{(x^2+y^2)^3}
=\frac{2xy^3-6x^3y}{(x^2+y^2)^3}.
\]
\textbf{Sonuç:} $f_{xy}(x,y)=f_{yx}(x,y)$. Bu eşitlik, karma türevler $(0,0)$ dışında sürekli olduğu için
\textbf{Schwarz (Clairaut) teoreminin} beklediği sonuçtur: \emph{$f_x$, $f_y$ türevlenebilir ve karma türevler sürekli
ise $f_{xy}=f_{yx}$}.
\end{cevap}

\begin{dikkat}[Tahtadaki notlarla ilgili düzeltmeler]
\begin{itemize}
  \item $f_{xy}$'nin ara satırında ``$-4y^2$'' yazılmıştı (\foto{IMG-20261022-WA0055}, \foto{IMG-20261108-WA0091});
        doğrusu $-4y\cdot y^3=-4y^4$'tür. Defter notundaki $f_{xx}$ ara satırının ilk terimi de ``$2x^3y$'' olarak
        görünüyor; doğrusu $-2x^3y$ (\foto{IMG-20261108-WA0091}'de doğru). Kutulanan sonuçlar doğrudur.
  \item Defter notunda (\foto{IMG-20261022-WA0059}) ``$f_{xy}=f_{yx}$, eğer $f(x,y)=f(y,x)$ ise geçerlidir'' yazılmıştı.
        Bu yanlıştır: karma türevlerin eşitliği simetriye değil, karma türevlerin \textbf{sürekliliğine} bağlıdır
        (Soru 1'deki simetrik olmayan fonksiyonda da $f_{xy}=f_{yx}$).
\end{itemize}
\end{dikkat}

%-------------------------------------------------------------------------------------------
\subsection{Lineerleştirme ve hata sınırı (22.10 dersi)}

\begin{soru}[Soru 3 \foto{IMG-20261022-WA0060}, \foto{IMG-20261108-WA0093}]
$f(x,y)=x^2+3xy+2y^2+4$ fonksiyonunun $(2,1)$ noktasındaki lineerleştirmesini bulunuz.
$B=\{(x,y):|x-2|\le0{,}01,\ |y-1|\le0{,}01\}$ dikdörtgensel bölgesi üzerindeki hatanın bir üst sınırını hesaplayınız.
\end{soru}
\begin{cevap}[Çözüm (adım adım)]
\textbf{Kullanılan formüller:}
$L(x,y)=f(x_0,y_0)+f_x(x_0,y_0)(x-x_0)+f_y(x_0,y_0)(y-y_0)$ ve
$|E(x,y)|\le\dfrac12M\big(|x-x_0|+|y-y_0|\big)^2$, burada $M$, $B$ üzerinde $|f_{xx}|,|f_{xy}|,|f_{yy}|$'nin üst sınırı.

\textbf{Adım 1 -- Fonksiyon değeri:} $f(2,1)=2^2+3\cdot2\cdot1+2\cdot1^2+4=4+6+2+4=16$.

\textbf{Adım 2 -- Birinci kısmi türevler:}
$f_x=2x+3y\Rightarrow f_x(2,1)=4+3=7$;\quad $f_y=3x+4y\Rightarrow f_y(2,1)=6+4=10$.

\textbf{Adım 3 -- Lineerleştirme:}
\[ L(x,y)=16+7(x-2)+10(y-1)=16+7x-14+10y-10=7x+10y-8 . \]
\textbf{Adım 4 -- İkinci türevler ve $M$:} $f_{xx}=2$, $f_{xy}=3$, $f_{yy}=4$. Sabit oldukları için $B$'nin her yerinde
aynıdır; $M=\max\{2,3,4\}=4$.

\textbf{Adım 5 -- Hata sınırı:} $B$'de $|x-2|\le0{,}01$ ve $|y-1|\le0{,}01$:
\[
|E(x,y)|\le\frac12\cdot4\,\big(|x-2|+|y-1|\big)^2\le2\,(0{,}01+0{,}01)^2=2\cdot0{,}0004=0{,}0008 .
\]
\textbf{Yorum:} $B$ içinde $f$ yerine $7x+10y-8$ kullanılırsa hata en fazla $0{,}0008$ olur. (Bu $f$ ikinci derece polinom
olduğundan hatanın tam ifadesi $E=(x-2)^2+3(x-2)(y-1)+2(y-1)^2$'dir; $B$'deki en büyük değeri $0{,}0006$'dır, yani bulunan
sınır geçerlidir.)
\end{cevap}

%-------------------------------------------------------------------------------------------
\subsection{Tam diferansiyel ve diferansiyelle yaklaşık hesap -- Problemler (05.11 dersi)}

\begin{soru}[Soru 4 \foto{IMG-20261105-WA0034}, \foto{IMG-20261108-WA0082} (Problem 1a)]
$f(x,y)=x^3+5xy^2+3xy+4x-3y+5$ fonksiyonunun birinci ve ikinci mertebeden tam diferansiyellerini bulunuz.
\end{soru}
\begin{cevap}[Çözüm (adım adım)]
\textbf{Formüller:} $df=f_x\,dx+f_y\,dy$,\quad $d^2f=f_{xx}\,dx^2+2f_{xy}\,dx\,dy+f_{yy}\,dy^2$.

\textbf{Birinci türevler:} $f_x=3x^2+5y^2+3y+4$ ($-3y$ ve $5$'in $x$'e göre türevi $0$),\quad
$f_y=10xy+3x-3$ ($x^3$, $4x$, $5$'in $y$'ye göre türevi $0$).
\[ df=(3x^2+5y^2+3y+4)\,dx+(10xy+3x-3)\,dy . \]
\textbf{İkinci türevler:} $f_{xx}=6x$,\quad $f_{xy}=(f_x)_y=10y+3$ \ (kontrol: $f_{yx}=(10xy+3x-3)_x=10y+3$ $\checkmark$),\quad
$f_{yy}=10x$.
\[ d^2f=6x\,dx^2+2(10y+3)\,dx\,dy+10x\,dy^2 . \]
\end{cevap}

\begin{soru}[Soru 5 \foto{IMG-20261105-WA0034}, \foto{IMG-20261105-WA0035}, \foto{IMG-20261108-WA0082}, \foto{IMG-20261108-WA0083} (Problem 1b)]
$w=f(x,y,z)=3x^2yz^2-3xy+2yz^2+xyz$ fonksiyonunun birinci ve ikinci mertebeden tam diferansiyellerini bulunuz.
\end{soru}
\begin{cevap}[Çözüm (adım adım)]
\textbf{Formüller (üç değişken):} $dw=w_x\,dx+w_y\,dy+w_z\,dz$,
\[ d^2w=w_{xx}dx^2+w_{yy}dy^2+w_{zz}dz^2+2w_{xy}dx\,dy+2w_{xz}dx\,dz+2w_{yz}dy\,dz . \]
\textbf{Birinci türevler} (her seferinde diğer iki değişken sabit):
\begin{align*}
w_x&=6xyz^2-3y+yz, & w_y&=3x^2z^2-3x+2z^2+xz, & w_z&=6x^2yz+4yz+xy .
\end{align*}
\[ dw=(6xyz^2-3y+yz)\,dx+(3x^2z^2-3x+2z^2+xz)\,dy+(6x^2yz+4yz+xy)\,dz . \]
\textbf{İkinci türevler:}
\begin{align*}
w_{xx}&=(w_x)_x=6yz^2, & w_{yy}&=(w_y)_y=0, & w_{zz}&=(w_z)_z=6x^2y+4y,\\
w_{xy}&=(w_x)_y=6xz^2-3+z, & w_{xz}&=(w_x)_z=12xyz+y, & w_{yz}&=(w_y)_z=6x^2z+4z+x .
\end{align*}
(Kontrol: $(w_y)_x=6xz^2-3+z=w_{xy}$, $(w_z)_x=12xyz+y=w_{xz}$, $(w_z)_y=6x^2z+4z+x=w_{yz}$ $\checkmark$.)

$w_{yy}=0$ olduğundan $dy^2$ terimi yoktur:
\begin{align*}
d^2w={}&6yz^2\,dx^2+(6x^2y+4y)\,dz^2+2(6xz^2+z-3)\,dx\,dy\\
&+2(12xyz+y)\,dx\,dz+2(6x^2z+x+4z)\,dy\,dz .
\end{align*}
\end{cevap}

\begin{soru}[Soru 6 \foto{IMG-20261105-WA0035}, \foto{IMG-20261108-WA0083}, \foto{IMG-20261108-WA0084} (Problem 2)]
Bir dik dairesel silindirin yarıçapı 10 cm'den 9,9 cm'ye azalıyor, yüksekliği 30 cm'den 30,3 cm'ye çıkıyor. Hacimdeki
değişimi diferansiyel ile yaklaşık olarak hesaplayınız.
\end{soru}
\begin{cevap}[Çözüm (adım adım)]
\textbf{Adım 1 -- Fonksiyon:} $V=\pi r^2h$.

\textbf{Adım 2 -- Tam diferansiyel:} $V_r=2\pi rh$, $V_h=\pi r^2$ $\Rightarrow$ $dV=2\pi rh\,dr+\pi r^2\,dh$.

\textbf{Adım 3 -- Değerler:} $r=10$, $h=30$; $dr=9{,}9-10=-0{,}1$ (azalma, negatif), $dh=30{,}3-30=0{,}3$.

\textbf{Adım 4 -- Hesap:}
\[
dV=2\pi\cdot10\cdot30\cdot(-0{,}1)+\pi\cdot100\cdot0{,}3=-60\pi+30\pi=-30\pi\ \text{cm}^3\approx-94{,}25\ \text{cm}^3 .
\]
\textbf{Adım 5 -- Yeni hacim:} $V=\pi\cdot100\cdot30=3000\pi$ $\Rightarrow$ yeni hacim $\cong3000\pi-30\pi=2970\pi$ cm$^3$.

Hacim yaklaşık $30\pi$ cm$^3$ azalır. (Kontrol: gerçek yeni hacim $\pi\cdot9{,}9^2\cdot30{,}3\approx2969{,}70\pi$ cm$^3$,
gerçek değişim $\approx-30{,}30\pi$ cm$^3$.)
\end{cevap}

\begin{soru}[Soru 7 \foto{IMG-20261105-WA0036}, \foto{IMG-20261108-WA0085} (Problem 3)]
Yarıçapı $r=1$, yüksekliği $h=5$ olan bir silindirin yarıçapı $0{,}03$ arttırılıp, yüksekliği $0{,}1$ azaltılırsa
hacimdeki yaklaşık değişim ve yüzde değişim nedir?
\end{soru}
\begin{cevap}[Çözüm (adım adım)]
\textbf{Adım 1:} $dV=2\pi rh\,dr+\pi r^2\,dh$; $r=1$, $h=5$, $dr=0{,}03$, $dh=-0{,}1$.

\textbf{Adım 2:} $dV=2\pi\cdot1\cdot5\cdot0{,}03+\pi\cdot1\cdot(-0{,}1)=0{,}3\pi-0{,}1\pi=0{,}2\pi\approx0{,}63$.

\textbf{Adım 3 -- Bağıl ve yüzde değişim:} Başlangıç hacmi $V=\pi\cdot1^2\cdot5=5\pi$.
\[ \frac{dV}{V}=\frac{0{,}2\pi}{5\pi}=0{,}04\quad\Rightarrow\quad\text{hacim yaklaşık \%4 artar.} \]
\textbf{Alternatif (logaritmik diferansiyel):} $\ln V=\ln\pi+2\ln r+\ln h$ $\Rightarrow$
$\dfrac{dV}{V}=2\dfrac{dr}{r}+\dfrac{dh}{h}=2(0{,}03)+\dfrac{-0{,}1}{5}=0{,}06-0{,}02=0{,}04$.

(Kontrol: gerçek değişim $\pi(1{,}03^2\cdot4{,}9-5)\approx0{,}198\pi$, yani yaklaşık \%3,97.)
\end{cevap}
Not: \foto{IMG-20261108-WA0085}'in sol kenarında bu sorunun yalnızca kesik bir parçası görünür; soru metni ve çözüm
\foto{IMG-20261105-WA0036}'dan alınmıştır.

%-------------------------------------------------------------------------------------------
\subsection{Lineerleştirme ve hata sınırı -- Problemler (05.11 dersi)}

\begin{soru}[Soru 8 \foto{IMG-20261105-WA0037}, \foto{IMG-20261108-WA0085} (Problem 4)]
$f(x,y)=e^{2y-x}$ fonksiyonunun $P(0,0)$ ve $Q(1,2)$ noktalarındaki lineerleştirmelerini bulunuz.
\end{soru}
\begin{cevap}[Çözüm (adım adım)]
\textbf{Kısmi türevler} (zincir kuralı; iç fonksiyon $u=2y-x$, $u_x=-1$, $u_y=2$):
$f_x=-e^{2y-x}$,\quad $f_y=2e^{2y-x}$.

\textbf{$P(0,0)$'da:} $2y-x=0$; $f=e^0=1$, $f_x=-1$, $f_y=2$.
\[ L(x,y)=1+(-1)(x-0)+2(y-0)=1-x+2y . \]
\textbf{$Q(1,2)$'de:} $2y-x=3$; $f=e^3$, $f_x=-e^3$, $f_y=2e^3$.
\[ L(x,y)=e^3-e^3(x-1)+2e^3(y-2)=e^3\big(1-x+1+2y-4\big)=e^3(2y-x-2). \]
\end{cevap}

\begin{soru}[Soru 9 \foto{IMG-20261105-WA0038}, \foto{IMG-20261105-WA0039}, \foto{IMG-20261108-WA0085}, \foto{IMG-20261108-WA0086} (Problem 5a)]
$f(x,y)=x^2-3xy+y^2$ fonksiyonunun $P(1,2)$ noktasındaki $L(x,y)$ lineerleştirmesini ve
$B=\{(x,y):|x-1|\le0{,}1,\ |y-2|\le0{,}1\}$ bölgesindeki hatanın bir üst sınırını bulunuz.
\end{soru}
\begin{cevap}[Çözüm (adım adım)]
\textbf{Adım 1:} $f(1,2)=1-6+4=-1$.

\textbf{Adım 2:} $f_x=2x-3y\Rightarrow f_x(1,2)=2-6=-4$;\quad $f_y=-3x+2y\Rightarrow f_y(1,2)=-3+4=1$.

\textbf{Adım 3:} $L(x,y)=-1-4(x-1)+1\cdot(y-2)=-1-4x+4+y-2=-4x+y+1$.

\textbf{Adım 4:} $f_{xx}=2$, $f_{xy}=-3$, $f_{yy}=2$. $M$ mutlak değerlerden seçilir:
$M=\max\{|2|,|-3|,|2|\}=3$.

\textbf{Adım 5:}
\[ |E(x,y)|\le\frac12\cdot3\,(0{,}1+0{,}1)^2=1{,}5\cdot0{,}04=0{,}06 . \]
(Tahtada ``$f_x=2x-3y=f_x(1,2)=-4$'' biçiminde ``$\Rightarrow$'' yerine ``$=$'' yazılmıştı; anlamı
$f_x(1,2)=-4$'tür.)
\end{cevap}

\begin{soru}[Soru 10 \foto{IMG-20261105-WA0040}, \foto{IMG-20261108-WA0086} (Problem 5b)]
$f(x,y)=1+y+x\cos y$ fonksiyonunun $P(0,0)$ noktasındaki lineerleştirmesini ve
$B=\{(x,y):|x|\le0{,}2,\ |y|\le0{,}2\}$ bölgesindeki hatanın bir üst sınırını bulunuz.
\end{soru}
\begin{cevap}[Çözüm (adım adım)]
\textbf{Adım 1:} $f(0,0)=1+0+0=1$.

\textbf{Adım 2:} $f_x=\cos y\Rightarrow f_x(0,0)=1$;\quad $f_y=1-x\sin y\Rightarrow f_y(0,0)=1$.

\textbf{Adım 3:} $L(x,y)=1+1\cdot x+1\cdot y=1+x+y$.

\textbf{Adım 4 -- İkinci türevler:} $f_{xx}=0$, $f_{xy}=(\cos y)_y=-\sin y$, $f_{yy}=(1-x\sin y)_y=-x\cos y$.

\textbf{Adım 5 -- $B$ üzerinde üst sınır:} $|\sin y|\le|y|$ ve $|\cos y|\le1$ eşitsizlikleriyle
\[
|f_{xy}|=|\sin y|\le|y|\le0{,}2,\qquad |f_{yy}|=|x|\,|\cos y|\le0{,}2\cdot1=0{,}2,\qquad |f_{xx}|=0
\ \Rightarrow\ M=0{,}2 .
\]
\textbf{Adım 6:}
\[ |E(x,y)|\le\frac12\cdot0{,}2\,(0{,}2+0{,}2)^2=0{,}1\cdot0{,}16=0{,}016 . \]
\textbf{Dikkat:} İkinci türevler yalnızca $P(0,0)$'da hesaplansaydı hepsi $0$ olurdu; $M$, $B$'nin \emph{tamamındaki}
üst sınır olmalıdır.
\end{cevap}

%-------------------------------------------------------------------------------------------
\subsection{Tam diferansiyel ve yüksek mertebeden diferansiyeller (08.11 dersi)}

\begin{soru}[Soru 11 \foto{IMG-20261108-WA0065}, \foto{IMG-20261108-WA0066}, \foto{IMG-20261108-WA0067}]
Yarıçapı $r=3$ cm, yüksekliği $h=10$ cm olan bir dik dairesel silindirin yarıçapı $0{,}01$ cm arttırılıp yüksekliği
$0{,}1$ cm azaltılırsa hacmindeki değişim nedir?
\end{soru}
\begin{cevap}[Çözüm (adım adım)]
\textbf{Adım 1:} $V=\pi r^2h\Rightarrow dV=V_r\,dr+V_h\,dh=2\pi rh\,dr+\pi r^2\,dh$.

\textbf{Adım 2:} $r=3$, $h=10$, $dr=0{,}01$, $dh=-0{,}1$.

\textbf{Adım 3:}
\[
dV=2\pi\cdot3\cdot10\,(0{,}01)+\pi\cdot9\,(-0{,}1)=0{,}6\pi-0{,}9\pi=-0{,}3\pi\ \text{cm}^3\approx-0{,}94\ \text{cm}^3 .
\]
\textbf{Adım 4:} $V=\pi\cdot9\cdot10=90\pi$ $\Rightarrow$ yeni hacim $\approx90\pi-0{,}3\pi=89{,}7\pi$ cm$^3\approx281{,}8$ cm$^3$.

(Kontrol: gerçek değer $\pi\cdot3{,}01^2\cdot9{,}9\approx89{,}695\pi$ cm$^3$.)
\end{cevap}
\begin{dikkat}[Tahtada düzeltilen hata]
Tahtada son satır ``$90\pi-0{,}3\pi=89{,}7$ cm$^3$'' olarak yazılmıştı; $\pi$ düşmüştür. Doğrusu $89{,}7\pi$ cm$^3$
($\approx281{,}8$ cm$^3$).
\end{dikkat}

\begin{soru}[Soru 12 \foto{IMG-20261108-WA0068}, \foto{IMG-20261108-WA0069}]
$z=x^2y$ fonksiyonunun ikinci mertebeden diferansiyelini ($d^2z$) bulunuz.
\end{soru}
\begin{cevap}[Çözüm (adım adım)]
\textbf{Formül:} $dx$, $dy$ sabit artımlar kabul edilerek $d^2z=d(dz)=z_{xx}\,dx^2+2z_{xy}\,dx\,dy+z_{yy}\,dy^2$
(ikinci türevler sürekli olduğundan $z_{xy}=z_{yx}$).

\textbf{Türevler:} $z_x=2xy$, $z_y=x^2$; $z_{xx}=2y$, $z_{xy}=2x$, $z_{yy}=0$.
\[ d^2z=2y\,dx^2+2\cdot2x\,dx\,dy+0=2y\,dx^2+4x\,dx\,dy . \]
\end{cevap}

%-------------------------------------------------------------------------------------------
\subsection{Ek alıştırmalar}

\begin{soru}[\ekal Ek alıştırma 1 (tanımla kısmi türev)]
$f(x,y)=xy^2-x^3$ için $f_x(1,-1)$'i limit tanımıyla hesaplayınız ve kuralla doğrulayınız.
\end{soru}
\begin{cevap}
$f(1,-1)=1-1=0$; $f(1+h,-1)=(1+h)-(1+h)^3=1+h-1-3h-3h^2-h^3=-2h-3h^2-h^3$.
\[ f_x(1,-1)=\lim_{h\to0}\frac{-2h-3h^2-h^3}{h}=\lim_{h\to0}(-2-3h-h^2)=-2 . \]
Kuralla: $f_x=y^2-3x^2\Rightarrow f_x(1,-1)=1-3=-2$ $\checkmark$.
\end{cevap}

\begin{soru}[\ekal Ek alıştırma 2 (ikinci mertebe ve Clairaut)]
$f(x,y)=x\,e^{y}+y\ln x$ $(x>0)$ fonksiyonunun tüm ikinci mertebe kısmi türevlerini bulunuz; $f_{xy}=f_{yx}$
olduğunu gösteriniz.
\end{soru}
\begin{cevap}
$f_x=e^y+\dfrac yx$, \quad $f_y=xe^y+\ln x$.
\[
f_{xx}=-\frac{y}{x^2},\qquad f_{xy}=e^y+\frac1x,\qquad f_{yx}=e^y+\frac1x,\qquad f_{yy}=xe^y .
\]
$f_{xy}=f_{yx}$; karma türevler $x>0$ bölgesinde sürekli olduğundan bu Schwarz teoremiyle uyumludur.
\end{cevap}

\begin{soru}[\ekal Ek alıştırma 3 (teğet düzlem ve normal doğru)]
$z=3x^2-xy+y^2$ yüzeyinin $(1,2,5)$ noktasındaki teğet düzleminin ve normal doğrusunun denklemlerini bulunuz.
\end{soru}
\begin{cevap}
$f(1,2)=3-2+4=5$ (nokta yüzey üzerinde). $f_x=6x-y\Rightarrow f_x(1,2)=4$; $f_y=-x+2y\Rightarrow f_y(1,2)=3$.
\[ z-5=4(x-1)+3(y-2)\quad\Longrightarrow\quad z=4x+3y-5 . \]
Normal doğru (doğrultu $\langle4,3,-1\rangle$): $x=1+4t$, $y=2+3t$, $z=5-t$.
\end{cevap}

\begin{soru}[\ekal Ek alıştırma 4 (lineerleştirme ve hata sınırı)]
$f(x,y)=\ln x+\ln y$ fonksiyonunun $(1,1)$'deki lineerleştirmesini ve $B:\ |x-1|\le0{,}2,\ |y-1|\le0{,}2$ üzerindeki
hata sınırını bulunuz.
\end{soru}
\begin{cevap}
$f(1,1)=0$, $f_x=\dfrac1x\Rightarrow1$, $f_y=\dfrac1y\Rightarrow1$: \ $L(x,y)=(x-1)+(y-1)=x+y-2$.\\
$f_{xx}=-\dfrac1{x^2}$, $f_{yy}=-\dfrac1{y^2}$, $f_{xy}=0$. $B$'de $x,y\ge0{,}8$ olduğundan
$|f_{xx}|,|f_{yy}|\le\dfrac1{0{,}64}=1{,}5625=M$ (en büyük değer kutunun $0{,}8$ kenarında alınır).
\[ |E|\le\frac{1{,}5625}{2}\,(0{,}2+0{,}2)^2=0{,}78125\cdot0{,}16=0{,}125 . \]
\end{cevap}

\begin{soru}[\ekal Ek alıştırma 5 (birinci ve ikinci diferansiyel)]
$z=x\sin y$ için $dz$ ve $d^2z$'yi bulunuz.
\end{soru}
\begin{cevap}
$z_x=\sin y$, $z_y=x\cos y$ $\Rightarrow$ $dz=\sin y\,dx+x\cos y\,dy$.\\
$z_{xx}=0$, $z_{xy}=\cos y$, $z_{yy}=-x\sin y$ $\Rightarrow$ $d^2z=2\cos y\,dx\,dy-x\sin y\,dy^2$.
\end{cevap}

\begin{soru}[\ekal Ek alıştırma 6 (diferansiyelle yaklaşık değer)]
$(1{,}02)^{3{,}01}$ sayısını diferansiyel yardımıyla yaklaşık hesaplayınız.
\end{soru}
\begin{cevap}
$f(x,y)=x^y$, $(x_0,y_0)=(1,3)$, $dx=0{,}02$, $dy=0{,}01$. $f(1,3)=1$;
$f_x=y\,x^{y-1}\Rightarrow3$; $f_y=x^y\ln x\Rightarrow1\cdot\ln1=0$.
\[ df=3(0{,}02)+0(0{,}01)=0{,}06\quad\Rightarrow\quad(1{,}02)^{3{,}01}\approx1{,}06 . \]
(Hesap makinesi: $\approx1{,}0614$.)
\end{cevap}

\begin{soru}[\ekal Ek alıştırma 7 (yüzde değişim)]
Bir dirençte harcanan güç $P=\dfrac{V^2}{R}$'dir. Gerilim \%2, direnç \%1 artarsa güç yaklaşık yüzde kaç değişir?
\end{soru}
\begin{cevap}
$\ln P=2\ln V-\ln R\Rightarrow\dfrac{dP}{P}=2\dfrac{dV}{V}-\dfrac{dR}{R}=2(0{,}02)-0{,}01=0{,}03$.
Güç yaklaşık \textbf{\%3 artar}.
\end{cevap}

\begin{soru}[\ekal Ek alıştırma 8 (üç değişkenli ikinci diferansiyel)]
$w=xyz$ için $dw$ ve $d^2w$'yi bulunuz.
\end{soru}
\begin{cevap}
$w_x=yz$, $w_y=xz$, $w_z=xy$ $\Rightarrow$ $dw=yz\,dx+xz\,dy+xy\,dz$.\\
$w_{xx}=w_{yy}=w_{zz}=0$; $w_{xy}=z$, $w_{xz}=y$, $w_{yz}=x$ $\Rightarrow$
$d^2w=2z\,dx\,dy+2y\,dx\,dz+2x\,dy\,dz$.
\end{cevap}

\clearpage\section{Konu 6: Zincir Kuralı, Yönlü Türevler ve Gradyan Vektörü}

Bu bölümde derste (5 ve 8 Kasım 2026) tahtada çözülen tüm sorular, işlenme sırasıyla ve adım adım çözümleriyle yer alıyor.
Yönlü türev, gradyan ve teğet düzlem konuları derste tahtada işlenmediği için bunlara ait sorular sondaki ek
alıştırmalarda verilmiştir.

\begin{bilgi}[Kullanılan temel formüller]
\begin{itemize}
  \item $z=f(x,y)$, $x(t),y(t)$: $\dfrac{dz}{dt}=z_x\dfrac{dx}{dt}+z_y\dfrac{dy}{dt}$.
  \item $z=f(x,y)$, $x(u,v),y(u,v)$: $z_u=z_xx_u+z_yy_u$, $z_v=z_xx_v+z_yy_v$ (ağaçta $u$'ya inen tüm yollar).
  \item Kapalı türev: $F(x,y)=0\Rightarrow y'=-F_x/F_y$; $F(x,y,z)=0\Rightarrow z_x=-F_x/F_z$, $z_y=-F_y/F_z$.
  \item $D_{\mathbf u}f=\nabla f\cdot\mathbf u$ ($|\mathbf u|=1$); en büyük yönlü türev $|\nabla f|$, $\nabla f$ yönünde.
\end{itemize}
\end{bilgi}

\subsection{5 Kasım dersi: zincir kuralı problemleri}

\begin{soru}[Soru 1 (tahtada Problem 6) \foto{IMG-20261105-WA0041}, \foto{IMG-20261108-WA0087}]
$\dfrac{dz}{dt}$'yi (i) zincir kuralıyla, (ii) $z$'yi $t$ cinsinden yazıp doğrudan türev alarak bulunuz.\\
\textbf{a)} $z=x^2+y^2,\quad x=\cos t+\sin t,\quad y=\cos t-\sin t$\\
\textbf{b)} $z=\arctan\dfrac{y}{x},\quad x=\cos t,\quad y=\sin t$
\end{soru}
\begin{cevap}[Çözüm (adım adım)]
\textbf{a) (i) Zincir kuralı.} $z$, $x$ ve $y$ üzerinden yalnızca $t$'ye bağlı olduğundan Durum 1 uygulanır:
$\frac{dz}{dt}=z_x\frac{dx}{dt}+z_y\frac{dy}{dt}$.
\begin{itemize}
  \item Dış türevler: $z_x=2x$, $z_y=2y$.
  \item İç türevler: $x'=-\sin t+\cos t$, $y'=-\sin t-\cos t$.
\end{itemize}
\begin{align*}
\frac{dz}{dt}&=2(\cos t+\sin t)(\cos t-\sin t)+2(\cos t-\sin t)(-\sin t-\cos t)\\
&=2(\cos^2t-\sin^2t)-2(\cos t-\sin t)(\cos t+\sin t)\\
&=2(\cos^2t-\sin^2t)-2(\cos^2t-\sin^2t)=\boxed{0}.
\end{align*}
\textbf{(ii) Doğrudan.} $(\cos t\pm\sin t)^2=1\pm2\sin t\cos t$ olduğundan $z=(1+2\sin t\cos t)+(1-2\sin t\cos t)=2$.
Sabit fonksiyonun türevi $0$'dır: $dz/dt=0$. İki yol uyuşuyor.

\textbf{b) (i) Zincir kuralı.} $\dfrac{d}{dw}\arctan w=\dfrac{1}{1+w^2}$ ve $w=y/x$ ile:
\[
z_x=\frac{1}{1+\frac{y^2}{x^2}}\cdot\left(-\frac{y}{x^2}\right)=-\frac{y}{x^2+y^2},\qquad
z_y=\frac{1}{1+\frac{y^2}{x^2}}\cdot\frac{1}{x}=\frac{x}{x^2+y^2}.
\]
$x=\cos t$, $y=\sin t$ için $x^2+y^2=1$, $x'=-\sin t$, $y'=\cos t$:
\[
\frac{dz}{dt}=\underbrace{(-\sin t)}_{z_x}\underbrace{(-\sin t)}_{x'}+\underbrace{\cos t}_{z_y}\underbrace{\cos t}_{y'}
=\sin^2t+\cos^2t=\boxed{1}.
\]
\textbf{(ii) Doğrudan.} $z=\arctan\dfrac{\sin t}{\cos t}=\arctan(\tan t)=t$ ($-\frac\pi2<t<\frac\pi2$), dolayısıyla
$dz/dt=1$. (Tahtada b şıkkının (ii) yolu yazılmamıştı; kaynaklara göre tamamlandı.)
\end{cevap}

\begin{soru}[Soru 2 (tahtada Problem 7) \foto{IMG-20261105-WA0042}]
$x=e^t$ dönüşümü ile $x^2\dfrac{d^2y}{dx^2}+x\dfrac{dy}{dx}+y=0$ denklemini $t$ değişkenine göre yazınız.
\end{soru}
\begin{cevap}[Çözüm (adım adım)]
\textbf{1. Dönüşümün türevi.} $x=e^t\iff t=\ln x$ ($x>0$); $\dfrac{dt}{dx}=\dfrac1x=e^{-t}$.

\textbf{2. Birinci türev (zincir kuralı).} $y$, $t$ üzerinden $x$'e bağlıdır:
$\dfrac{dy}{dx}=\dfrac{dy}{dt}\cdot\dfrac{dt}{dx}=e^{-t}\dfrac{dy}{dt}$.

\textbf{3. İkinci türev.} $\frac{dy}{dx}=e^{-t}\frac{dy}{dt}$ bir $t$ fonksiyonudur; $x$'e göre türevi için önce $t$'ye göre
türevleyip (çarpım kuralı) $\frac{dt}{dx}$ ile çarparız:
\begin{align*}
\frac{d^2y}{dx^2}&=\frac{d}{dt}\left(e^{-t}\frac{dy}{dt}\right)\cdot\frac{dt}{dx}
=\left(-e^{-t}\frac{dy}{dt}+e^{-t}\frac{d^2y}{dt^2}\right)e^{-t}
=e^{-2t}\left(\frac{d^2y}{dt^2}-\frac{dy}{dt}\right).
\end{align*}
\textbf{4. Katsayılarla çarpım.} $x^2=e^{2t}$, $x=e^t$:
\[
x^2\frac{d^2y}{dx^2}=\frac{d^2y}{dt^2}-\frac{dy}{dt},\qquad x\frac{dy}{dx}=\frac{dy}{dt}.
\]
\textbf{5. Denklem.}
$\dfrac{d^2y}{dt^2}-\dfrac{dy}{dt}+\dfrac{dy}{dt}+y=0\ \Longrightarrow\ \boxed{\dfrac{d^2y}{dt^2}+y=0}$.

Bu sabit katsayılı denklemin genel çözümü $y=c_1\cos t+c_2\sin t$, yani $y=c_1\cos(\ln x)+c_2\sin(\ln x)$'tir.

\textbf{Not:} Tahtada son satırda ``$x\cdot\frac{dx}{dy}$'' yazılmıştı; doğrusu $x\cdot\frac{dy}{dx}$'tir. Sonuç denklemi
fotoğrafta kadraj dışında kaldı \okunamadi{sağ kenardaki kutulu sonuç}; kaynaklara göre doğrusu $\frac{d^2y}{dt^2}+y=0$'dır.
\end{cevap}

\begin{soru}[Soru 3 (tahtada Problem 8, kısmen görünüyor) \foto{IMG-20261105-WA0042}]
\okunamadi{sorunun metni ve çözümü fotoğrafın sağ kenarında kesik; yalnızca ``8) a, \dots'', ``$x-a\dots$'',
``oldu\dots'', ``$F(x\dots$'', ``$F_x$'', ``$\partial z/\partial x$'', ``$\frac{a}{a-b\dots}$'' parçaları seçilebiliyor}.
Görünen parçalar kaynaklarda yaygın olan şu soruyla uyumludur (\textbf{tahmini metin}):
$a,b$ sabit, $\varphi$ türevlenebilir bir fonksiyon ve $x-az=\varphi(y-bz)$ ise $a\dfrac{\partial z}{\partial x}+b\dfrac{\partial z}{\partial y}=1$
olduğunu gösteriniz.
\end{soru}
\begin{cevap}[Çözüm (adım adım)]
\textbf{1. Kapalı biçim.} $F(x,y,z)=x-az-\varphi(y-bz)=0$. Burada $\varphi'$, $w=y-bz$ noktasında hesaplanır.

\textbf{2. Kısmi türevler} (Durum 4: $\varphi$ tek değişkenli, $w$ iki değişkenli):
\[
F_x=1,\qquad F_y=-\varphi'(w)\cdot w_y=-\varphi',\qquad F_z=-a-\varphi'(w)\cdot w_z=-a+b\varphi'=-(a-b\varphi').
\]
\textbf{3. Kapalı türev formülleri} ($a-b\varphi'\ne0$):
\[
\frac{\partial z}{\partial x}=-\frac{F_x}{F_z}=\frac{1}{a-b\varphi'},\qquad
\frac{\partial z}{\partial y}=-\frac{F_y}{F_z}=-\frac{\varphi'}{a-b\varphi'}.
\]
\textbf{4. İstenen ifade.}
\[
a\frac{\partial z}{\partial x}+b\frac{\partial z}{\partial y}=\frac{a}{a-b\varphi'}-\frac{b\varphi'}{a-b\varphi'}=\frac{a-b\varphi'}{a-b\varphi'}=1 .
\]
Tahtada görünen ``$\frac{a}{a-b\dots}$'' parçası, $a\,\partial z/\partial x$ teriminin bu biçimine karşılık geliyor.
\end{cevap}

\subsection{8 Kasım dersi: zincir kuralı durumları ve kapalı türev}

\begin{soru}[Soru 4 \foto{IMG-20261108-WA0072}, \foto{IMG-20261108-WA0073}, \foto{IMG-20261108-WA0074}]
$z=\ln(x^2+y^2)$, $x=e^{v}\cos u$, $y=e^{v}\sin u$ ise $\dfrac{\partial z}{\partial u}$ ve $\dfrac{\partial z}{\partial v}$'yi bulunuz.
\end{soru}
\begin{cevap}[Çözüm (adım adım)]
\textbf{1. Ağaç.} $z\to x,y$; $x\to u,v$; $y\to u,v$. Formüller:
$z_u=z_xx_u+z_yy_u$, $\ z_v=z_xx_v+z_yy_v$.

\textbf{2. Dış türevler.} $z_x=\dfrac{2x}{x^2+y^2}$, $z_y=\dfrac{2y}{x^2+y^2}$; ayrıca $x^2+y^2=e^{2v}(\cos^2u+\sin^2u)=e^{2v}$.

\textbf{3. İç türevler.} $x_u=-e^{v}\sin u$, $x_v=e^{v}\cos u$, $y_u=e^{v}\cos u$, $y_v=e^{v}\sin u$.

\textbf{4. Hesap.}
\begin{align*}
\frac{\partial z}{\partial u}&=\frac{2e^{v}\cos u\,(-e^{v}\sin u)+2e^{v}\sin u\;e^{v}\cos u}{e^{2v}}=\frac{-2e^{2v}\sin u\cos u+2e^{2v}\sin u\cos u}{e^{2v}}=\boxed{0},\\[4pt]
\frac{\partial z}{\partial v}&=\frac{2e^{v}\cos u\;e^{v}\cos u+2e^{v}\sin u\;e^{v}\sin u}{e^{2v}}=\frac{2e^{2v}(\cos^2u+\sin^2u)}{e^{2v}}=\boxed{2}.
\end{align*}
\textbf{5. Kontrol.} $z=\ln(e^{2v})=2v$ $\Rightarrow$ $z_u=0$, $z_v=2$. \checkmark
(Tahtadaki kutuda üs ``$e^{2u}$'' gibi okunuyor; doğrusu $e^{2v}$.)
\end{cevap}

\begin{soru}[Soru 5 \foto{IMG-20261108-WA0075}]
$z=\dfrac{1}{x^2+y^2}$, $x=\cos2t$, $y=\sin2t$ ise $\dfrac{dz}{dt}$'yi bulunuz.
\end{soru}
\begin{cevap}[Çözüm (adım adım)]
\textbf{1.} $x^2+y^2=\cos^22t+\sin^22t=1$.

\textbf{2. Kısmi türevler.} $z=(x^2+y^2)^{-1}$ olduğundan
$z_x=-(x^2+y^2)^{-2}\cdot2x=\dfrac{-2x}{(x^2+y^2)^2}$, $z_y=\dfrac{-2y}{(x^2+y^2)^2}$.

\textbf{3. İç türevler.} $x'=-2\sin2t$, $y'=2\cos2t$ (içteki $2t$'nin türevi $2$ unutulmamalı).

\textbf{4. Durum 1.}
\[
\frac{dz}{dt}=(-2\cos2t)(-2\sin2t)+(-2\sin2t)(2\cos2t)=4\sin2t\cos2t-4\sin2t\cos2t=\boxed{0}.
\]
\textbf{5. Kontrol.} $z=\dfrac{1}{\cos^22t+\sin^22t}=1\Rightarrow dz/dt=0$. (Tahtada ``$z=$'' yazılıp bırakılmıştı.)
\end{cevap}

\begin{soru}[Soru 6 \foto{IMG-20261108-WA0077}, \foto{IMG-20261108-WA0078}]
$s=\ln(x^2+y^2+z^2)$, $x=u-v$, $y=u+v$, $z=2\sqrt{uv}$ ise $\dfrac{\partial s}{\partial u}$ ve $\dfrac{\partial s}{\partial v}$'yi bulunuz.
\end{soru}
\begin{cevap}[Çözüm (adım adım)]
\textbf{1. Ağaç (Durum 3).} $s\to x,y,z$; her biri $\to u,v$. Her türevde üç terim:
$s_u=s_xx_u+s_yy_u+s_zz_u$, $\ s_v=s_xx_v+s_yy_v+s_zz_v$.

\textbf{2. Dış türevler.} $s_x=\dfrac{2x}{N}$, $s_y=\dfrac{2y}{N}$, $s_z=\dfrac{2z}{N}$, burada
$N=x^2+y^2+z^2=(u-v)^2+(u+v)^2+4uv=2u^2+2v^2+4uv=2(u+v)^2$.

\textbf{3. İç türevler.} $x_u=1,\ x_v=-1,\ y_u=1,\ y_v=1,\ z_u=\dfrac{v}{\sqrt{uv}},\ z_v=\dfrac{u}{\sqrt{uv}}$
($z=2\sqrt{uv}$ için $z_u=2\cdot\frac{v}{2\sqrt{uv}}$).

\textbf{4. Hesap.} $2z\,z_u=2\cdot2\sqrt{uv}\cdot\frac{v}{\sqrt{uv}}=4v$ ve $2z\,z_v=4u$ olduğundan
\begin{align*}
\frac{\partial s}{\partial u}&=\frac{2(u-v)+2(u+v)+4v}{2(u+v)^2}=\frac{4u+4v}{2(u+v)^2}=\boxed{\frac{2}{u+v}},\\
\frac{\partial s}{\partial v}&=\frac{-2(u-v)+2(u+v)+4u}{2(u+v)^2}=\frac{4u+4v}{2(u+v)^2}=\boxed{\frac{2}{u+v}}.
\end{align*}
\textbf{5. Kontrol.} $s=\ln\big(2(u+v)^2\big)=\ln2+2\ln(u+v)\Rightarrow s_u=s_v=\dfrac{2}{u+v}$. \checkmark
\end{cevap}

\begin{soru}[Soru 7 \foto{IMG-20261108-WA0079}]
$y=\sin x$, $x=u^2v$ ise $y_u$ ve $y_v$'yi bulunuz.
\end{soru}
\begin{cevap}[Çözüm (adım adım)]
\textbf{1. Ağaç (Durum 4).} $y\to x$ (adi türev $dy/dx$), $x\to u,v$ (kısmi türevler).

\textbf{2.} $\dfrac{dy}{dx}=\cos x=\cos(u^2v)$; $x_u=2uv$, $x_v=u^2$.

\textbf{3.} $y_u=\dfrac{dy}{dx}x_u=\boxed{2uv\cos(u^2v)}$,\qquad $y_v=\dfrac{dy}{dx}x_v=\boxed{u^2\cos(u^2v)}$.
\end{cevap}

\begin{soru}[Soru 8 \foto{IMG-20261108-WA0080}, \foto{IMG-20261108-WA0081}]
$2x^4z^3+3xy^2+4xz+3y+5x-23=0$ denklemi $z=z(x,y)$ fonksiyonunu tanımlıyor. $(1,2)$ noktasında $z_x$ ve $z_y$'yi bulunuz.
\end{soru}
\begin{cevap}[Çözüm (adım adım)]
\textbf{1. Formül (WA0080).} $F(x,y,z)=0$'ın $x$'e göre türevi: $F_x\cdot1+F_y\cdot0+F_zz_x=0\Rightarrow z_x=-F_x/F_z$;
benzer şekilde $z_y=-F_y/F_z$ ($F_z\ne0$).

\textbf{2. Noktayı tamamla.} $x=1,\ y=2$: $2z^3+12+4z+6+5-23=0\Rightarrow2z^3+4z=0\Rightarrow2z(z^2+2)=0$.
$z^2+2>0$ olduğundan tek reel çözüm $z=0$; nokta $(1,2,0)$.

\textbf{3. Kısmi türevler} (değişkenler bağımsızmış gibi):
$F_x=8x^3z^3+3y^2+4z+5$, $F_y=6xy+3$, $F_z=6x^4z^2+4x$.

\textbf{4. Noktada:} $F_x=17$, $F_y=15$, $F_z=4\ne0$.

\textbf{5. Sonuç:} $z_x=\boxed{-\dfrac{17}{4}}$, $\quad z_y=\boxed{-\dfrac{15}{4}}$.
\end{cevap}

\subsection{Ek alıştırmalar}

\begin{soru}[\ekal Ek alıştırma 1 (Durum 2)]
$w=xy+z$, $x=\cos t$, $y=\sin t$, $z=t$ ise $dw/dt$'yi bulunuz ve $t=0$'daki değerini hesaplayınız.
\end{soru}
\begin{cevap}
$w_x=y,\ w_y=x,\ w_z=1$; $x'=-\sin t,\ y'=\cos t,\ z'=1$.
\[
\frac{dw}{dt}=\sin t(-\sin t)+\cos t\cos t+1=\cos^2t-\sin^2t+1=\cos2t+1 .
\]
$t=0$'da $dw/dt=2$. Kontrol: $w=\sin t\cos t+t=\frac12\sin2t+t$, $w'=\cos2t+1$. \checkmark
\end{cevap}

\begin{soru}[\ekal Ek alıştırma 2 (kutupsal koordinatlar)]
$z=x^2+xy$, $x=r\cos\theta$, $y=r\sin\theta$ ise $z_r$ ve $z_\theta$'yı zincir kuralıyla bulunuz.
\end{soru}
\begin{cevap}
$z_x=2x+y$, $z_y=x$; $x_r=\cos\theta$, $y_r=\sin\theta$, $x_\theta=-r\sin\theta$, $y_\theta=r\cos\theta$.
\begin{align*}
z_r&=(2x+y)\cos\theta+x\sin\theta=(2r\cos\theta+r\sin\theta)\cos\theta+r\cos\theta\sin\theta
=2r\cos^2\theta+r\sin2\theta,\\
z_\theta&=-(2x+y)\,r\sin\theta+x\,r\cos\theta=-2r^2\sin\theta\cos\theta-r^2\sin^2\theta+r^2\cos^2\theta
=r^2(\cos2\theta-\sin2\theta).
\end{align*}
Kontrol: $z=r^2\cos^2\theta+r^2\sin\theta\cos\theta=r^2\left(\cos^2\theta+\tfrac12\sin2\theta\right)$; doğrudan türevler
aynı sonucu verir. \checkmark
\end{cevap}

\begin{soru}[\ekal Ek alıştırma 3 (Durum 4, keyfi fonksiyon)]
$f$ türevlenebilir ve $z=f(x^2-y^2)$ ise $y\,z_x+x\,z_y=0$ olduğunu gösteriniz.
\end{soru}
\begin{cevap}
$w=x^2-y^2$ dersek $z=f(w)$ (Durum 4): $z_x=f'(w)\,w_x=2x\,f'(w)$, $z_y=f'(w)\,w_y=-2y\,f'(w)$. Buna göre
$y\,z_x+x\,z_y=2xy\,f'-2xy\,f'=0$.
\end{cevap}

\begin{soru}[\ekal Ek alıştırma 4 (kapalı türev)]
$x^3+y^3+z^3+6xyz=1$ denklemi $z=z(x,y)$'yi tanımlıyorsa $z_x$ ve $z_y$'yi bulunuz.
\end{soru}
\begin{cevap}
$F=x^3+y^3+z^3+6xyz-1$; $F_x=3x^2+6yz$, $F_y=3y^2+6xz$, $F_z=3z^2+6xy$. $F_z\ne0$ iken
\[
z_x=-\frac{3x^2+6yz}{3z^2+6xy}=-\frac{x^2+2yz}{z^2+2xy},\qquad z_y=-\frac{y^2+2xz}{z^2+2xy}.
\]
\end{cevap}

\begin{soru}[\ekal Ek alıştırma 5 (Euler--Cauchy)]
$x=e^t$ dönüşümüyle $x^2y''+3xy'+y=0$ denklemini $t$'ye göre yazınız ve çözünüz.
\end{soru}
\begin{cevap}
$x^2y''=\ddot y-\dot y$ ve $xy'=\dot y$ (nokta $t$'ye göre türev) kullanılırsa
\[
\ddot y-\dot y+3\dot y+y=0\ \Longrightarrow\ \ddot y+2\dot y+y=0 .
\]
Karakteristik denklem $m^2+2m+1=(m+1)^2=0$, çift kök $m=-1$: $y=(c_1+c_2t)e^{-t}$. $t=\ln x$ ile
$y=\dfrac{c_1+c_2\ln x}{x}$. (Kontrol: $y=1/x$ için $x^2\cdot\frac{2}{x^3}+3x\cdot\left(-\frac{1}{x^2}\right)+\frac1x=0$.)
\end{cevap}

\begin{soru}[\ekal Ek alıştırma 6 (yönlü türev)]
$f(x,y)=x^2y^3-4y$ fonksiyonunun $(2,-1)$ noktasında $\mathbf v=2\mathbf i+5\mathbf j$ yönündeki yönlü türevini bulunuz.
\end{soru}
\begin{cevap}
$f_x=2xy^3$, $f_y=3x^2y^2-4$; $(2,-1)$'de $f_x=2\cdot2\cdot(-1)=-4$, $f_y=3\cdot4\cdot1-4=8$; $\nabla f=\langle-4,8\rangle$.
$\mathbf v$ birim değil: $|\mathbf v|=\sqrt{29}$, $\mathbf u=\frac{1}{\sqrt{29}}\langle2,5\rangle$.
\[
D_{\mathbf u}f(2,-1)=\frac{-4\cdot2+8\cdot5}{\sqrt{29}}=\boxed{\frac{32}{\sqrt{29}}}\approx5{,}94 .
\]
\end{cevap}

\begin{soru}[\ekal Ek alıştırma 7 (en hızlı artış, verilen noktaya doğru yön)]
$f(x,y)=x\,e^{y}$ ve $P(2,0)$ olsun. (a) $P$'den $Q(\frac12,2)$'ye doğru yöndeki yönlü türevi bulunuz. (b) $f$ en hızlı
hangi yönde artar, en büyük artış hızı nedir? (c) Hangi yönlerde değişim hızı sıfırdır?
\end{soru}
\begin{cevap}
$\nabla f=\langle e^{y},\ x e^{y}\rangle$, $\nabla f(2,0)=\langle1,2\rangle$.

\textbf{(a)} $\overrightarrow{PQ}=\langle-\frac32,2\rangle$, $|\overrightarrow{PQ}|=\frac52$, $\mathbf u=\langle-\frac35,\frac45\rangle$;
$D_{\mathbf u}f=-\frac35+\frac85=\boxed{1}$.

\textbf{(b)} $\nabla f(2,0)=\langle1,2\rangle$ yönünde (birim: $\frac{1}{\sqrt5}\langle1,2\rangle$); en büyük hız $|\nabla f|=\sqrt5$.

\textbf{(c)} $\mathbf u\perp\langle1,2\rangle$: $\mathbf u=\pm\frac{1}{\sqrt5}\langle2,-1\rangle$.
\end{cevap}

\begin{soru}[\ekal Ek alıştırma 8 (üç değişken)]
$f(x,y,z)=xy+yz+zx$ için $P(1,-1,2)$ noktasında (a) $Q(2,1,4)$'e doğru yönlü türevi, (b) maksimum yönlü türevi bulunuz.
\end{soru}
\begin{cevap}
$\nabla f=\langle y+z,\ x+z,\ x+y\rangle$, $\nabla f(P)=\langle1,3,0\rangle$.

\textbf{(a)} $\overrightarrow{PQ}=\langle1,2,2\rangle$, uzunluğu $3$; $D_{\mathbf u}f=\frac{1\cdot1+3\cdot2+0\cdot2}{3}=\boxed{\frac73}$.

\textbf{(b)} $\max D_{\mathbf u}f=|\nabla f(P)|=\sqrt{1+9+0}=\boxed{\sqrt{10}}$, $\langle1,3,0\rangle$ yönünde.
\end{cevap}

\begin{soru}[\ekal Ek alıştırma 9 (teğet düzlem ve normal doğru)]
$x^2+y^2-z^2=4$ tek kanatlı hiperboloidinin $(2,1,1)$ noktasındaki teğet düzlemini ve normal doğrusunu bulunuz.
\end{soru}
\begin{cevap}
Kontrol: $4+1-1=4$ \checkmark. $F=x^2+y^2-z^2$, $\nabla F=\langle2x,2y,-2z\rangle$, $\nabla F(2,1,1)=\langle4,2,-2\rangle\parallel\langle2,1,-1\rangle$.

Teğet düzlem: $2(x-2)+(y-1)-(z-1)=0\iff\boxed{2x+y-z=4}$.

Normal doğru: $x=2+2t,\quad y=1+t,\quad z=1-t$ \ (simetrik: $\frac{x-2}{2}=\frac{y-1}{1}=\frac{z-1}{-1}$).
\end{cevap}

\end{document}
